% !TeX program = xelatex
% Сборка: xelatex coveo_session_recommender_math.tex (до стабилизации ссылок).
% Требуются XeLaTeX, polyglossia и стандартные шрифты Windows: Times New Roman, Arial и Consolas.
% Самостоятельный файл; внешние .bib, изображения и включаемые файлы не нужны.
\documentclass[12pt,a4paper]{article}
\usepackage{fontspec}
\setmainfont{Times New Roman}
\setsansfont{Arial}
\setmonofont{Consolas}[Scale=MatchLowercase]
\newfontfamily\cyrillicfont{Times New Roman}[Script=Default]
\newfontfamily\cyrillicfontsf{Arial}[Script=Default]
\newfontfamily\cyrillicfonttt{Consolas}[Script=Default,Scale=MatchLowercase]
\usepackage{polyglossia}
\setdefaultlanguage{russian}
\usepackage{amsmath,amssymb,amsthm,mathtools}
\usepackage{array,booktabs,longtable,geometry,enumitem}
\usepackage{microtype,xurl,hyperref,ragged2e}
\AtBeginDocument{\RaggedRight}
\geometry{margin=2.2cm}
\hypersetup{hidelinks,pdftitle={Рекомендательная система Coveo: математическая спецификация}}
\setlength{\parindent}{0pt}
\setlength{\parskip}{0.32\baselineskip}
\setlength{\emergencystretch}{3em}
\setlist{nosep,leftmargin=2em}
\setcounter{tocdepth}{2}
\setcounter{secnumdepth}{3}
\allowdisplaybreaks[2]
\newcolumntype{L}[1]{>{\raggedright\arraybackslash}p{#1}}
\newcommand{\code}[1]{{\ttfamily\small\nolinkurl{#1}}}
\newcommand{\R}{\mathbb R}
\newcommand{\N}{\mathbb N}
\newcommand{\E}{\mathbb E}
\newcommand{\Prob}{\mathbb P}
\newcommand{\ind}{\mathbf 1}
\newcommand{\eps}{\varepsilon}
\newcommand{\norm}[1]{\left\lVert#1\right\rVert}
\newcommand{\abs}[1]{\left|#1\right|}
\newcommand{\pos}[1]{\left(#1\right)_+}
\newcommand{\given}{\,\vert\,}
\DeclareMathOperator*{\argmin}{arg\,min}
\DeclareMathOperator*{\argmax}{arg\,max}
\DeclareMathOperator{\Var}{Var}
\DeclareMathOperator{\Cov}{Cov}
\DeclareMathOperator{\sigmoid}{sigmoid}
\DeclareMathOperator{\Top}{Top}
\DeclareMathOperator{\clip}{clip}
\newtheoremstyle{plainru}{0.45\baselineskip}{0.45\baselineskip}{\normalfont}{}{\bfseries}{.}{0.5em}{}
\theoremstyle{plainru}
\newtheorem{definition}{Определение}[section]
\newtheorem{assumption}{Предположение}[section]
\newtheorem{proposition}{Утверждение}[section]
\newtheorem{invariant}{Инвариант}[section]
\newtheorem{testcase}{Проверка}[section]
\newtheorem{remark}{Замечание}[section]
\makeatletter
\renewcommand*\l@section[2]{%
  \ifnum \c@tocdepth >\z@
    \addpenalty\@secpenalty\addvspace{0.6em \@plus\p@}%
    \setlength\@tempdima{2.8em}%
    \begingroup\parindent\z@\rightskip\@pnumwidth
      \parfillskip-\@pnumwidth\leavevmode\bfseries
      \advance\leftskip\@tempdima\hskip-\leftskip
      #1\nobreak\hfil\nobreak\hb@xt@\@pnumwidth{\hss#2}\par
    \endgroup\fi}
\renewcommand*\l@subsection{\@dottedtocline{2}{2.8em}{3.4em}}
\makeatother
\begin{document}
\begin{center}
{\Large\bfseries Рекомендательная система для интернет-магазина:\\[0.3em]
математическая спецификация на данных Coveo SIGIR eCommerce\par}
\vspace{0.7em}
{\large Граф взаимодействий, признаки товаров\par
и поисковый контекст сессии\par}
\vspace{0.7em}
{Версия 2.10.\par}
{Условия демонстрационного сервиса уточнены 8 октября 2026 года.\par}
\end{center}
\begin{abstract}
Первый вариант использует граф переходов, готовые векторы описаний
и изображений, поиск и ранжирование LambdaMART. Нейросетевое расширение
обучает товарные представления и намерения сессии.
Для каждого этапа определены входы, порядок обучения, проверка
и условия перехода к следующему этапу.
Бизнес-цель --- дополнительные сессии с покупкой на 1000 допущенных сессий.
Первый результат на Coveo --- охват зарегистрированных будущих покупок
первой выдачей среди сессий с такой покупкой. Полнота телеметрии
остаётся неизвестной; оценка по всем сессиям требует её подтверждения.
Данные позволяют измерить совпадение с историческими покупками;
причинный эффект требует рандомизированного эксперимента с показами.
Спецификация не содержит результатов обучения или нагрузочных измерений.
\end{abstract}

\section*{Порядок реализации}
\phantomsection\label{sec:reading-route}
\label{sec:first-run}\label{sec:stage-acceptance}
\addcontentsline{toc}{section}{Порядок реализации}

\begin{longtable}{L{0.22\textwidth}L{0.69\textwidth}}
\toprule Этап & Разделы и результат \\\midrule\endhead
Базовая система & Часть~\ref{part:core}, рецепт~\ref{app:first-config}
и проверки~\ref{app:tests-base}. Первый отчёт --- на
$C_{\mathrm{smoke}}$; без подбора настроек и открытия $D$. \\
Выбор и итоговая оценка & Порядок периодов и выбора ---
раздел~\ref{sec:eval-primary-comparison}; состояния результата ---
раздел~\ref{sec:acceptance}. \\
Нейросеть & Часть~\ref{part:neural}, параметры~\ref{app:neural-config},
проверки~\ref{app:tests-neural}. Требуются воспроизводимая база,
диагностика её потерь и измеренная стоимость. \\
Отдельные исследования & Последовательность~\ref{sec:sequence-control},
обновление графа~\ref{sec:freshness}, SID~\ref{sec:sid-overview},
другие модели~\ref{sec:model-extensions}. Для каждого заранее задают
вопрос, контроль и бюджет. \\
Эксперимент с показами & Протокол~\ref{sec:eval-online},
после проверки обслуживания. Измеряет причинный эффект. \\
\bottomrule
\end{longtable}

Условия перехода между этапами заданы один раз в
паспорте~\ref{app:run-stages}. Значения параметров ---
в приложении~\ref{app:parameters}, состав отчёта ---
в~\ref{app:diagnostic-output}. Если $D$ уже использован для оценки базы,
следующее расширение требует нового итогового периода.

Текущая цель реализации --- локальная историческая оценка до $D$
по протоколу \code{observed_purchases_v1}
(раздел~\ref{sec:eval-observed}) в режиме \code{ideal_deterministic}.
Он задаёт единые правила записанной истории, целей, метрик и приёмки.
Оценка частоты покупок среди всех сессий потребовала бы подтверждённой
полноты наблюдений и отдельной постановки; здесь она не выполняется.
Готовность сервиса проверяется
отдельно и не является условием этого маршрута. Нейросетевые исследования,
сервис и A/B остаются последующими этапами со всеми требованиями
соответствующих разделов.

\clearpage
\begingroup
\small
\setlength{\parskip}{0.1\baselineskip}
\tableofcontents
\endgroup
\clearpage

\part{Основная система: данные, обучение и проверка}\label{part:core}
\section{Назначение и выбранная постановка}\label{sec:scope}

Рекомендации строятся по текущей сессии без постоянного идентификатора
покупателя. На вход поступают доступная история, последний поисковый
запрос и версия товарных признаков. Ответ содержит упорядоченный список,
версию модели, источники кандидатов и состояние расчёта.

Первая реализация --- политика \(\pi_{\mathrm{base}}\): графовые
переходы, сходство предоставленных векторов товаров, поисковый контекст
и ранжирование LambdaMART. Далее она обозначается GCS+Rk:
G --- граф, C --- признаки описаний и изображений,
S --- поиск, Rk --- ранжировщик.
Первый запуск определён в части~I, нейросетевые представления ---
в части~II, дополнительные исследования --- в приложениях.

Начальные настройки и нагрузочные цели задают требования;
производительность и бизнес-эффект предстоит измерить.

\begin{longtable}{L{0.28\textwidth}L{0.64\textwidth}}
\toprule Объект & Основной вариант \\\midrule\endhead
Цель внедрения & Увеличение доли подходящих сессий с покупкой при ограничениях
на задержку, ошибки и стоимость. \\
Единица результата & Сессия, допущенная по правилам, проверяемым до первого
показа экспериментальных рекомендаций. \\
Отбор товаров первого запуска & Одношаговые переходы по историческому
графу, содержательные и поисковые кандидаты, недавние товары и популярность. \\
Ранжирование & Деревья с парным обучением; основная метка --- последующая
покупка товара. Клик и корзина рассматриваются отдельно. \\
Историческая проверка & Единые протокол, сессии и запросы для всех методов;
первая оценка Coveo --- по разделу~\ref{sec:eval-observed}. \\
Доказательство эффекта & Рандомизированный эксперимент с учётом всех
назначенных сессий, включая ошибки обслуживания. \\
Подвыборка для прототипа & Полные сессии из заранее заданного набора;
результат не объявляется официальным результатом соревнования. \\
\bottomrule
\end{longtable}

\begin{invariant}[Разделение трёх величин]
Баллы модели, совпадение выдачи с историческими покупками и изменение
конверсии после внедрения --- разные величины. Между ними не вводится
равенство без дополнительных данных и предположений.
\end{invariant}

\section{Обозначения и единицы}\label{sec:symbols}

\(\ind\{A\}\) --- индикатор условия \(A\), \(\E\) --- математическое
ожидание, \(\Prob\) --- вероятность. \(\Top_m\) возвращает не более
\(m\) элементов с наибольшим баллом; равенства разрешаются
лексикографическим порядком идентификаторов. Время измеряется в секундах.

\small
\begin{longtable}{L{0.23\textwidth}L{0.69\textwidth}}
\toprule Обозначение & Смысл \\\midrule\endhead
\(s,i,j\) & Сессия и идентификаторы товаров. \\
\(e,t(e),a(e)\) & Событие, его время и время доступности системе. \\
\(q,s(q),t_q\) & Запрос рекомендаций, его сессия и время. \\
\(t_{0,s},H,g\) & Первый допустимый запрос, горизонт результата и интервал
неактивности, равный 30 минутам для сессий Coveo. \\
\(\mathcal H_q,\mathcal F_q\) & Доступная история и вся информация,
которой располагает алгоритм при запросе. \\
\(r_G,r_N,r(q),\mathcal I_q\) & Время графового снимка, публикации
обученного отбора, использованного снимка и допустимый каталог. \\
\(\mathcal Q_s,Q_{\max}\) & Запросы сессии и их заранее установленный предел. \\
\(\pi,\pi_0\) & Полная политика рекомендаций и политика сравнения. \\
\(\pi_{\mathrm{base}}\) & Первый реализуемый вариант GCS+Rk. \\
\(\mathcal C_q(\pi),M\) & Кандидаты и максимальное число кандидатов. \\
\(L_q(\pi),K\) & Итоговый список и предел его длины. \\
\(Y_q,y_{qi}\) & Множество будущих покупок после запроса и бинарная метка
кандидата \(i\). \\
\(Y_s(\pi)\) & Индикатор покупки в сессии при фактическом использовании
политики \(\pi\). Отличается от множества \(Y_q\). \\
\(B(\pi),\Delta B\) & Покупающие сессии на 1000 допущенных сессий и их причинный прирост. \\
\(\operatorname{Hit}_s(\pi),\operatorname{Hit@20}_J\) & Попадание первой выдачи
в зарегистрированные будущие покупки и его доля среди сессий с такой целью. \\
\(f_q(i)\) & Балл ранжирования, не вероятность покупки. \\
\(Z_s\) & Назначенный вариант эксперимента с показами, 0 или 1. \\
\(\xi,\widehat\xi\) & Конфигурация процедуры и выбранная конфигурация. \\
\(A_G,A_N,B,C,E,D\) & Блоки топологии и статистик, обучения нейросетевого
отбора, обучения ранжировщика, выбора, резерва и итоговой проверки;
\(A=A_G\cup A_N\). \\
\bottomrule
\end{longtable}
\normalsize

Индекс \(r\) у агрегата ограничивает события доступными к снимку.
Сложение балла и денежной стоимости требует обоснованного коэффициента
перевода. Политику сравнения \(\pi_0\) фиксируют заранее;
она не обязана совпадать с \(\pi_{\mathrm{base}}\).
Блок \(A_N\) зарезервирован для нейросети и не дополняет граф
или товарные статистики первого запуска.

\section{Данные Coveo и границы наблюдаемости}\label{sec:data}

\subsection{Наблюдаемые объекты}

Публикация содержит таблицы поведения, поисковых событий и товарных
признаков~\cite{coveo-data,coveo-paper}. Для дальнейших формул зададим
собственное представление события:
\begin{equation}\label{eq:data-event}
e=(s,i,t,a,\kappa,\eta),\qquad
\kappa\in\{\mathrm{detail},\mathrm{add},\mathrm{remove},
\mathrm{purchase},\mathrm{search},\mathrm{page}\}.
\end{equation}
Здесь \(i=\varnothing\) разрешён для страницы без товара и поиска;
\(\eta\) содержит дополнительные поля. Время \(t\) получается делением
\code{server_timestamp_epoch_ms} на 1000. Времена загрузки \(a\)
в исходном наборе не даны: в историческом воспроизведении основная модель
предполагает \(a=t\), а задержки рассматриваются отдельным испытанием.

Поиск представлен как
\begin{equation}\label{eq:data-search}
v=(s,t,u,R,C),
\end{equation}
где \(u\) --- \code{query_vector}, \(R\) --- список
\code{product_skus_hash}, \(C\) --- \code{clicked_skus_hash}.
Содержательный вектор товара \(d_i\) берётся из
\code{description_vector}, вектор изображения \(v_i\) --- из
\code{image_vector}. Категория \(c_i\) задаётся хешированной иерархией,
\(b_i^{\mathrm{price}}\) --- порядковым ценовым интервалом.

\subsection{Что не восстанавливается из набора}

\begin{longtable}{L{0.32\textwidth}L{0.60\textwidth}}
\toprule Отсутствующий объект & Следствие для постановки \\\midrule\endhead
Фактическая цена заказа, количество и маржа & Нельзя вычислить денежную
выручку и прибыль из номера ценового интервала. \\
Журнал показов нашего рекомендательного блока & Совпадение с покупкой
не устанавливает влияние рекомендации. \\
Вероятности назначения рекомендаций & Основной исторический эксперимент
не использует причинное перевзвешивание показов. \\
Исторический склад и даты активации товаров & Допустимость в каталоге
задаётся правилом наблюдаемости, а не наличием на складе. \\
Время каждого клика внутри строки поиска & Поле \(C\) не входит
в признаки текущего или следующего запроса без отдельного датированного события. \\
Начальный состав и количества корзины & Состояние корзины означает
последнее наблюдаемое действие, а не точный остаток. \\
Исторические версии содержательного кодировщика & Работа с предоставленными
векторами условна относительно опубликованного снимка признаков. \\
Подтверждение полноты телеметрии & Пустое множество записанных покупок
не доказывает отсутствия покупки. Оценка по положительно наблюдаемым
целям определена в разделе~\ref{sec:eval-observed}. \\
\bottomrule
\end{longtable}

\begin{assumption}[Доступность статических признаков]\label{ass:static-content}
В исторической проверке предоставленные векторы и категории
считаются доступными для товаров исторического каталога. Это допущение
о боковой информации, а не доказательство её фактической доступности
в прошлом. Обязательное сравнение без предоставленных метаданных проверяет
чувствительность к этому допущению (пункт~\ref{sec:metadata-sensitivity}).
Оно не восстанавливает неизвестную историческую версию векторов.
\end{assumption}

Ценовые интервалы кодируются как категории или сравниваются по порядку
и равенству; их разности и отношения не имеют денежного смысла.
Хешированные категории сравниваются по равенству и известной иерархии.

\subsection{Снимок и воспроизводимая подвыборка}

Снимок данных определяется кортежем
\begin{equation}\label{eq:data-snapshot}
\mathsf D=(\mathrm{version},\mathrm{hashes},\mathrm{schema},
T_{\min},T_{\max},\mathrm{sampling},\mathrm{cleaning}).
\end{equation}
Внутри заранее фиксированных временных блоков сессии сортируются
по устойчивому псевдослучайному ключу \(h_{\mathrm{seed}}(s)\).
Начальный алгоритм --- SHA-256 от конкатенации seed как unsigned
64-bit big-endian, длины ID в байтах как unsigned 64-bit big-endian
и точных UTF-8 байтов ID без изменения регистра или Unicode-нормализации.
Сортируют по 32 байтам digest лексикографически, при совпадении ---
по тем же байтам ID; seed равен 17. Размер $n$ задаётся отдельно
для каждого периода до просмотра его исходов.
Берутся первые \(n\) идентификаторов со всеми строками поведения и поиска,
независимо от будущих покупок. Так получаются вложенные обучающие выборки,
например 200 и 500 тысяч сессий.

Публикуются числа сессий, событий, товаров, покупающих сессий и запросов
поиска, доли пропущенных векторов и неизвестных тестовых товаров.
Покупающие сессии считаются по идентификаторам, не по строкам
\code{purchase}; размер каталога указывается отдельно от размера выборки.

\section{Время, очистка и доступная история}\label{sec:time}

\subsection{Порядок событий}

Поведение и поиск объединяются по \(s\) и времени; одновременные события
образуют блок. Номер строки задаёт порядок хранения, но не причинный порядок.
При историческом \(a=t\) блок содержит все строки этого времени;
в рабочем сервисе его завершённость задаёт закрытый конверт
раздела~\ref{sec:ops-transport-block}.

Один товарный просмотр может сопровождаться \code{detail} и
\code{pageview}. Основная товарная последовательность использует
\code{product_action}; страница без действия не создаёт второй
товарный переход. Все исходные строки сохраняются для аудита.
При отсутствии идентификатора события точное совпадение строк само по себе
не доказывает технический дубль. Основной вариант использует объявленную
канонизацию \code{event_semantics=set_v1}: одному сочетанию
\((s,t,i,\kappa)\) соответствует одно каноническое событие с сохранённой
кратностью исходных строк. Для поиска его полезная часть дополнительно
сохраняет множество различных значений \code{query_vector} с кратностью
каждого, а не первое или последнее значение.
\label{sec:data-search-canonicalization}
Корректные векторы сравниваются до нормировки: по размерности и точным
координатам FP64, причём $-0$ и $+0$ совпадают. Отсутствующее и
некорректное поле образуют отдельные маркеры с кодом причины;
исходные строки сохраняются.
Поля $R,C$ из~\eqref{eq:data-search} не участвуют в этом представлении.
Поисковый блок однозначен, только если содержит одно корректное значение
без маркеров. Два или более различных корректных вектора дают
признак неоднозначности; наличие маркера также отключает этот контекст
с соответствующей причиной. Повторы одного вектора неоднозначности
не создают. Порядок строк не выбирает один из запросов.
Далее $\mathcal H_q$, счётчики просмотров,
история $\mathcal D_q$ и последовательный кодировщик используют именно
канонические события. Это допущение о представлении, не восстановление
истинного числа действий; исходные строки и их кратность не удаляются.

На $C_{\mathrm{smoke}}$ обязательна проверка чувствительности
\code{event_semantics=multiset_v1}: кратность учитывается в счётчиках,
истории и усреднении временных групп. Бинарные покупочные цели остаются
множествами, а поддержка ребра по-прежнему считает различные сессии.
Отчёт содержит долю затронутых сессий, изменение признаков и охвата.
Выбор семантики фиксируется до $C_{\mathrm{select}}$ и одинаков для
сравниваемых моделей. Транспортные повторы рабочего сервиса удаляются по
устойчивому event ID до этой операции; совпадение полей его не заменяет.

\begin{invariant}[Одновременность]
Между событиями одного временного блока не строятся направленные
переходы и не назначаются цели <<после запроса>>. Запрос выполняется
после обработки всего доступного блока; его будущие цели имеют
строго большее время.
\end{invariant}

\subsection{Информация на момент решения}

Пусть \(a_q\) --- время начала расчёта. Тогда
\begin{equation}\label{eq:time-history}
\mathcal H_q=\{e:s(e)=s(q),\ t(e)\le t_q,\ a(e)\le a_q\}.
\end{equation}
В \(\mathcal F_q\) входят \(\mathcal H_q\), доступные снимки
справочников, опубликованные параметры модели и состояние сервиса.
Для обслуживания и независимого применения ранжировщика на $C,E,D$ требуется
\begin{equation}\label{eq:time-adapted}
\mathcal C_q,\ f_q,\ L_q\quad\text{являются функциями только }\mathcal F_q.
\end{equation}
Будущая длина сессии, итоговая конверсия, следующий поиск и будущая
популярность запрещены. Снимок \(r(q)<t_q\) выбирается по времени
публикации. В $B$ ранжировщик использует исторически допустимые признаки
и кандидатов, но обучающие предсказания не являются независимой оценкой.
Нейросеть обучается по продолжениям примеров \(A_N\) и проверяется
на последующих независимых периодах.

Первый и нейросетевой варианты используют только \(u\), без выдачи \(R\),
зависящей от неизвестной политики поискового движка.
Переупорядочивание полученного \(R\) --- отдельная задача;
\(C\) остаётся последующим откликом. Её результаты не объединяются
с результатами рекомендаций.

\subsection{Наблюдаемое состояние корзины}

Определим \(b_{qi}\in\{\mathrm{add},\mathrm{remove},\mathrm{unknown}\}\).
Начальное значение --- \(\mathrm{unknown}\). Последнее однозначное
наблюдаемое \code{add} задаёт \(\mathrm{add}\), \code{remove}
задаёт \(\mathrm{remove}\). Если последняя временная группа содержит
оба действия, значение снова равно \(\mathrm{unknown}\).
Повторное удаление не создаёт отрицательного количества товара.

Модель использует индикаторы трёх состояний и давность последнего
однозначного события. Фактический состав корзины не восстанавливается.

\section{Запросы, допустимые товары и будущие цели}\label{sec:requests}

\subsection{Начало наблюдения и расписание}

Подходящей считается сессия, в которой встретился \code{detail}
с корректным идентификатором товара и к концу его временного блока
ещё не наблюдалось покупки. Первый такой блок задаёт \(t_{0,s}\).
Правило не требует наличия будущей покупки, нескольких просмотров,
содержательного вектора или непустого ответа выбранного алгоритма.
Сессии, начавшиеся с покупки до первого подходящего просмотра,
не входят в целевую совокупность; их доля публикуется.

Единственный запрос $q_s^{(1)}$ выполняется в \(t_{0,s}\):
$\mathcal Q_s=\{q_s^{(1)}\}$, $Q_{\max}=1$ для всех методов.
Его префикс включает весь блок одновременных событий.
Будущая покупка не выбирает момент запроса. Последующие выдачи
в этом опыте не оцениваются. Однозначность последнего поискового
блока проверяется по разделу~\ref{sec:data-search-canonicalization}.

\begin{remark}
Историческое расписание \(\mathcal Q_s\) следует записанному пути.
Рекомендации могут изменить будущие действия и запросы; воспроизведение
журнала не моделирует это изменение.
\end{remark}

\subsection{Окно зарегистрированных целей}

Пусть $t_{\mathrm{last},s}^{\mathrm{obs}}$ --- время последнего записанного
события объединённого потока поведения и поиска сессии. Окно целей равно
\begin{equation}\label{eq:request-horizon}
\tau_s=\min(t_{0,s}+H,t_{\mathrm{last},s}^{\mathrm{obs}}).
\end{equation}
Последняя запись и ожидание $g=1800$ секунд не подтверждают закрытия
сессии или полноты телеметрии. Календарный запас $H+g$ по
разделу~\ref{sec:eval-maturity} защищает границы периодов, не доказывая
отсутствия пропущенных покупок. Горизонт первой версии $H=86400$ секунд.

\subsection{Каталог и метки}

В основном историческом варианте каталог фиксируется по $A_G$:
\begin{align}
\mathcal I_G&=\{i:\exists e\text{ из допустимых сессий }A_G,\ i(e)=i\},\notag\\
\mathcal I_q&=\mathcal I_G\setminus\mathcal P_q.
\label{eq:request-catalog}
\end{align}
где \(\mathcal P_q\) --- уже купленные в видимом префиксе товары.
Первый запрос допускается до наблюдавшейся покупки,
поэтому $\mathcal P_q=\varnothing$. Исключение оставлено
как защитный контракт для повторно используемого сервиса. Рекомендации
после покупки, повторное приобретение расходников и количества товара
относятся к отдельной задаче: для неё заново фиксируют расписание,
каталог и метки; бинарная цель текущего опыта не моделирует количество.
Метаданные не подтверждают историческую доступность будущих товаров.
Неизвестный снимку товар префикса получает соответствующую отметку;
основная и резервная выдачи остаются в \(\mathcal I_q\).

В реальном магазине множество~\eqref{eq:request-catalog} заменяется
доступным на момент запроса каталогом с ограничениями наличия,
региона и публикации. Качество такой фильтрации на Coveo не измеряется.
Товары с наблюдаемым \code{add} разрешены в основном варианте;
вариант рекомендаций только новых для корзины товаров оценивается
отдельно на том же знаменателе.

Для $s=s(q)$ множество зарегистрированных целей и кандидатные метки равны
\begin{align}
Y_q=Y_s^{\mathrm{obs}}&=\{i:\exists e\text{ в выгрузке},\ s(e)=s,
\ \kappa(e)=\mathrm{purchase},\ t_{0,s}<t(e)\le\tau_s,\ i(e)=i\},
\label{eq:request-target}\\
y_{qi}&=\ind\{i\in Y_q\}.\label{eq:request-label}
\end{align}
Цель содержит все различные SKU покупок в окне, а не только первый
или последний. Повторные записи покупки дают один элемент множества;
несколько товаров в одном временном блоке сохраняются.
Товары из \(Y_q\), отсутствующие в \(\mathcal I_q\), не удаляются
из оценки. Так отделяется ошибка покрытия каталога от ошибки ранжирования.
Пустое $Y_q$ означает отсутствие записи о целевой покупке, а не доказанную
непокупку. Даже при $Y_q\ne\varnothing$ полнота списка покупок неизвестна.
Нулевая кандидатная метка означает только отсутствие товара в $Y_q$;
правило её использования задано в разделе~\ref{sec:model-ranking}.

\section{Основной бизнес-критерий}\label{sec:business}

\subsection{Сессии с покупкой и причинный прирост}

Политика \(\pi\) включает модель, правила выдачи, резервное поведение
и ограничения числа запросов. Обозначим через \(Y_s(\pi)\)
индикатор хотя бы одной покупки после назначения политики
и до завершения её окна результата. При изменении политики
путь сессии и её завершение могут измениться. Результат определяется
одним и тем же правилом окна для каждого возможного пути.

\begin{definition}[Ценность политики]
Для заранее выбранной совокупности подходящих сессий
\begin{equation}\label{eq:business-value}
B(\pi)=1000\,\E[Y_s(\pi)],\qquad
\Delta B(\pi,\pi_0)=1000\,\E[Y_s(\pi)-Y_s(\pi_0)].
\end{equation}
Единица \(\Delta B\) --- дополнительные сессии с покупкой
на 1000 подходящих сессий.
\end{definition}

Под <<подходящими>> далее понимаются допущенные по правилу
раздела~\ref{sec:requests} сессии. Пусть $\mathcal S_{\mathrm{all}}$
--- все зарегистрированные сессии того же источника и календарного окна,
а $\mathcal S_{\mathrm{elig}}$ --- прошедшие это правило. Отдельно публикуют
$r_{\mathrm{elig}}=|\mathcal S_{\mathrm{elig}}|/|\mathcal S_{\mathrm{all}}|$
и числа исключений по причинам; созревание исхода считают следующим,
отдельным этапом. При пустом знаменателе доля недоступна.
Формула~\eqref{eq:business-value} условна по $s\in\mathcal S_{\mathrm{elig}}$.
Она не измеряет конверсию всех посещений. Умножение на $r_{\mathrm{elig}}$
само по себе не даёт общего эффекта: нужны неизменный допуск, отсутствие
эффекта вне допуска и единая исходная совокупность. Для KPI всех посещений
требуется отдельный эксперимент с назначением до отбора и собственным
знаменателем. При назначении только после допуска доля допуска не имеет
разреза по ещё не назначенному варианту.

В A/B со стабильным ключом основная когорта дополнительно ограничена
входными сессиями $\mathcal S_{\mathrm{entry}}\subseteq\mathcal S_{\mathrm{elig}}$
по~\ref{sec:eval-itt}. Тогда в~\eqref{eq:business-value} и бизнес-пороге
явно используют $B_{\mathrm{entry}}(\pi)=1000\E[Y_s(\pi)\mid
s\in\mathcal S_{\mathrm{entry}}]$, а также публикуют долю этой когорты
среди допущенных сессий. Подтверждение $\Delta B_{\mathrm{entry}}$
не переносится на все $\mathcal S_{\mathrm{elig}}$; нуждающийся в таком
переносе продукт требует отдельного estimand и протокола. Обозначение
$\Delta B$ в приёмке всегда сопровождается выбранной популяцией.

Любая покупка подходящей сессии даёт один успех независимо от количества
товаров и клика по рекомендации. Отбор по такому клику зависел бы
от поведения после назначения варианта.

Например, увеличение доли покупающих сессий с 0.030 до 0.032
соответствует \(\Delta B=2\), абсолютному приросту 0.2 процентного
пункта и относительному приросту около 6.7\%. Это арифметический пример,
а не оценка эффекта Coveo или разрабатываемой системы.

Выбор конфигурации для внедрения имеет вид
\begin{equation}\label{eq:business-optimization}
\max_{\pi\in\Pi}\ \Delta B(\pi,\pi_0)
\quad\text{при}\quad
\begin{cases}
\Prob(\ell_q>L_{\max})\le\rho_L,\\
\Prob(\mathrm{error}_q=1)\le\rho_{\mathrm{err}},\\
\E[C_s^{\mathrm{serve}}]\le c_{\max},\\
L_q\text{ удовлетворяет товарным ограничениям}.
\end{cases}
\end{equation}
Здесь \(\ell_q\) --- задержка ответа, \(C_s^{\mathrm{serve}}\) ---
стоимость обслуживания сессии в выбранной денежной единице.
Ограничения задаёт владелец продукта до сравнения систем.
При отсутствии денежных тарифов вместо стоимости фиксируется
измеримый вычислительный бюджет; денежная экономия не заявляется.

\subsection{Почему покупочная метка не решает причинную задачу}

Исторический журнал содержит один исход каждой сессии,
но не оба \(Y_s(\pi)\) и \(Y_s(\pi_0)\). Распознавание будущей покупки
позволяет выбрать модель для испытания, но не определяет
разность~\eqref{eq:business-value}.

\begin{proposition}[Неидентифицируемость эффекта по совпадениям]\label{prop:business-nonidentification}
При $0<p_0=\E[Y_s(\pi_0)]<1$ одинаковый закон наблюдаемых покупок
и одинаковые совпадения рекомендаций совместимы как с нулевым,
так и с положительным либо отрицательным эффектом нового блока.
\end{proposition}
\begin{proof}
Зафиксируем закон исходов при исходной политике. Три модели с
$Y_s(\pi)=Y_s(\pi_0)$, $Y_s(\pi)=1$ и $Y_s(\pi)=0$ дают соответственно
эффекты $0$, $1000(1-p_0)>0$ и $-1000p_0<0$, но совпадают на исходном
журнале, поскольку новая политика там не назначалась.
В общем случае бинарность ограничивает эффект интервалом
$[-1000p_0,\,1000(1-p_0)]$. При $p_0=0$ или $p_0=1$ возможны не все
знаки, однако величина эффекта по-прежнему не определена.
Отсутствие одного класса в конечном журнале не доказывает
соответствующее граничное значение $p_0$.
\end{proof}

\subsection{Денежное расширение при появлении необходимых данных}

Если в реальной системе доступны оплаченные позиции, возвраты,
себестоимость и дополнительные расходы, можно определить
\begin{equation}\label{eq:business-margin}
G_s(\pi)=\sum_{o\in\mathcal O_s(\pi)}
\bigl(R_o-C_o^{\mathrm{goods}}-C_o^{\mathrm{delivery}}
-C_o^{\mathrm{return}}\bigr)-C_s^{\mathrm{serve}},
\end{equation}
где \(\mathcal O_s\) --- заказы, атрибутированные сессии по заранее
заданному правилу. Все слагаемые должны иметь одну денежную единицу
и согласованное определение учёта налогов и возвратов.
Тогда \(1000\E[G_s(\pi)-G_s(\pi_0)]\) --- дополнительный
маржинальный доход на 1000 сессий. Горизонт возвратов задаётся отдельно.
В Coveo этот функционал имеет состояние \code{not_observed}.

Ценовой интервал не заменяет \(R_o\). Перевод \(\Delta B\) в деньги
требует предположений о числе заказов, корзине и марже, которые могут
меняться при внедрении. Основной бизнес-критерий --- \(\Delta B\).

\subsection{Уровни доказательности}

\begin{longtable}{L{0.26\textwidth}L{0.32\textwidth}L{0.32\textwidth}}
\toprule Уровень & Допустимый вывод & Необходимые данные \\\midrule\endhead
Историческая пригодность & Система находит будущие покупки в фиксированном
журнале и укладывается в бюджет. & Coveo и воспроизводимый протокол. \\
Эксплуатационная пригодность & Запросы обслуживаются с требуемой задержкой
и корректными признаками. & Нагрузочные испытания и теневой запуск. \\
Бизнес-превосходство & \(\Delta B\) превышает заданный практический порог. &
Рандомизированное назначение и зрелые исходы. \\
Денежная выгода & Дополнительная маржа положительна. & Эксперимент,
финансовые данные и завершённый учёт возвратов. \\
\bottomrule
\end{longtable}

\section{Граф взаимодействий с товарами}\label{sec:model-graph}

Используются история \(\mathcal H_q\) и допустимый каталог
\(\mathcal I_q\) из разделов~\ref{sec:time} и~\ref{sec:requests}.
Единица времени \(t_{\mathrm{unit}}=1\,\mathrm{s}\) позволяет записывать
безразмерные аргументы логарифмов.
Топология и статистики построены только по $A_G$ к моменту $r_G$.
Для зарегистрированных событий допуск истории задан
в разделе~\ref{sec:eval-observed}; подтверждённая полнота телеметрии
для этого протокола не предполагается.
Для базовой политики требуется \(r_G<t_q\); отдельного обучения
графового кодировщика нет. На втором этапе нейросетевой отбор
обучается на \(A_N\) и публикуется к \(r_N\), поэтому для его запросов
из \(B,C,E,D\) дополнительно требуется \(r_G<r_N<t_q\).

\subsection{Вершины, отношения и вклад одной сессии}

Вершиной графа служит идентификатор товара. Вводятся три отношения:
\begin{equation}\label{eq:model-relations}
\mathcal R_G=
\{\mathrm{detail}\!\to\!\mathrm{detail},\quad
  \mathrm{detail}\!\to\!\mathrm{add},\quad
  \mathrm{detail}\!\to\!\mathrm{purchase}\}.
\end{equation}
Все события, включая покупки, должны предшествовать снимку.

Для исторической сессии $s$ товарные события объединяются в непустые
группы с одинаковым временем; группы упорядочиваются по времени.
Пусть $i_{s\ell},\kappa_{s\ell},t_{s\ell}$ обозначают товар, действие
и время события $\ell$, а $k_s(\ell)$ --- номер его временной группы.
Внутри группы направленные рёбра не создаются: порядок неизвестен.
Для отношения $r=(\kappa_1,\kappa_2)$ множество допустимых пар имеет вид
\begin{align}\label{eq:model-event-pairs}
\mathcal P^r_{sij}(r_0)=\{(\ell,m):\;&i_{s\ell}=i,
\ i_{sm}=j,\ i\ne j,\ \kappa_{s\ell}=\kappa_1,\ \kappa_{sm}=\kappa_2,\nonumber\\
&0<t_{sm}-t_{s\ell}\le\Delta_G,\quad
0<k_s(m)-k_s(\ell)\le B_G,\quad t_{sm}<r_0\}.
\end{align}
Здесь $r_0$ --- время снимка, $\Delta_G>0$ --- временное окно,
$B_G\in\N$ --- максимальное расстояние между временными группами
товарных событий. События без товара не создают таких групп.
Пары не зависят от порядка строк внутри группы; повторения товара
учитываются признаками сессии, без петель $i\to i$.

Вклад сессии ограничен единицей и определяется максимумом:
\begin{equation}\label{eq:model-session-edge}
w^r_{sij}=
\begin{cases}
\displaystyle\max_{(\ell,m)\in\mathcal P^r_{sij}(r_0)}
\exp\!\left(-\frac{t_{sm}-t_{s\ell}}{\tau_G}\right),
&\mathcal P^r_{sij}(r_0)\ne\varnothing,\\[0.4em]
0,&\text{иначе},
\end{cases}
\end{equation}
где $\tau_G>0$ имеет единицы времени. Тогда
\begin{equation}\label{eq:model-edge-weight}
W^r_{ij}=\sum_s w^r_{sij},\qquad
n^r_{ij}=\sum_s\ind\{w^r_{sij}>0\},\qquad
0\le W^r_{ij}\le n^r_{ij}.
\end{equation}
Суммы используют сессии снимка с отсечением по его границе
согласно протоколу подготовки данных.

\subsection{Разреживание и регуляризация строк графа}

Среди рёбер с поддержкой $n^r_{ij}\ge n_G\ge1$ сохраняются
не более $d_G$ соседей с наибольшими $W^r_{ij}$ для каждой пары $(r,i)$.
Равенства разрешаются фиксированным порядком идентификаторов.
Множество сохранённых соседей обозначается $\mathcal N^r_i$.
\begin{equation}\label{eq:model-transition}
P^r_{ij}=
\begin{cases}
\displaystyle\frac{W^r_{ij}}
 {\sum_{k\in\mathcal N^r_i}W^r_{ik}},
&j\in\mathcal N^r_i,\\[0.4em]
0,&\text{иначе}.
\end{cases}
\end{equation}
Пустая строка равна нулю; непустая имеет сумму один.
$P^r_{ij}$ выражает нормированную силу связи, а не вероятность
следующего действия: сессия может породить нескольких соседей.

Для ослабления строк с малым суммарным весом используется
\begin{equation}\label{eq:model-edge-confidence}
d^r_i=\sum_{j\in\mathcal N^r_i}W^r_{ij},\qquad
\rho^r_i=\frac{d^r_i}{d^r_i+\lambda_G},\qquad \lambda_G>0.
\end{equation}
$\lambda_G$ измеряется в тех же условных единицах, что и $d^r_i$.
Коэффициент $\rho^r_i$ ослабляет строки малого веса;
он не оценивает достоверность ребра. Для сохранённого ребра
\begin{equation}\label{eq:model-row-regularization}
\rho_i^rP_{ij}^r=\frac{W_{ij}^r}{d_i^r+\lambda_G},
\qquad
0\le \rho_i^rP_{ij}^r\le P_{ij}^r\le1.
\end{equation}
При фиксированном $W_{ij}^r$ рост веса других соседей ослабляет ребро.
От случайных связей дополнительно защищает порог $n_G$.

Пусть $\mathcal D_q$ --- последние не более $L_H$ доступных
событий \code{detail}. При отсечении истории равенства времени
разрешаются идентификатором товара; это технический отбор,
не восстановление фактического порядка действий.
Нормированные веса истории задаются как
\begin{equation}\label{eq:model-history-weight}
\alpha_{q\ell}=
\frac{\exp(-(t_q-t_{s(q)\ell})/\tau_H)}
 {\sum_{m\in\mathcal D_q}\exp(-(t_q-t_{s(q)m})/\tau_H)},
\qquad \ell\in\mathcal D_q,\quad\tau_H>0.
\end{equation}
При пустой истории суммы с $\alpha_{q\ell}$ равны нулю.
Действия корзины учитываются отдельными признаками.

Графовый балл кандидата $j$ по отношению $r$ равен
\begin{equation}\label{eq:model-graph-score}
g^r_{qj}=\sum_{\ell\in\mathcal D_q}
\alpha_{q\ell}\rho^r_{i_{s(q)\ell}}
P^r_{i_{s(q)\ell},j}.
\end{equation}
Балл каждого отношения служит отдельным признаком:
$0\le g^r_{qj}\le1$, а при пустой истории $g^r_{qj}=0$.
Контрольный вариант может использовать смесь
$g_{qj}=\sum_r\omega_r g^r_{qj}$, где
$\omega_r\ge0$, $\sum_r\omega_r=1$.
В первом запуске такая смесь не требуется:
каждое отношение формирует собственные кандидаты и признаки.

Дополнительное ослабление отдельных рёбер описано в приложении~\ref{sec:model-edge-optional}.
В первом запуске используется только \eqref{eq:model-graph-score}.

\begin{invariant}[Граница исторического знания]
Изменение событий с временем $t\ge r(q)$ не меняет
$W^r,P^r,\rho^r$ данного снимка. Изменение событий с временем
$t>t_q$ не меняет $\alpha_{q\ell}$ и $g^r_{qj}$.
\end{invariant}

\subsection{Стоимость построения и запроса}

Пусть $N_{A_G}$ --- число исторических товарных событий,
$b_G^{\max}$ --- максимальный размер одной временной группы снимка.
После упорядочивания и группировки перебор пар требует порядка
$N_{A_G}B_Gb_G^{\max}$ операций: для события проверяются до $B_G$
следующих групп размера не более $b_G^{\max}$.
Наблюдаемое $b_G^{\max}$ входит в отчёт о стоимости;
произвольное разрезание одновременной группы запрещено.
Число сохранённых рёбер не превосходит
$|\mathcal R_G|\,|\mathcal I_G|d_G$.
Вычисление графовых кандидатов для одного запроса обращается
не более чем к $|\mathcal R_G|L_Hd_G$ рёбрам.
Память промежуточной агрегации измеряется отдельно от готового графа.
Плотная матрица $|\mathcal I_G|^2$ не создаётся; $B_G,L_H,d_G$
ограничивают вычислительную стоимость.

\section{Содержательные признаки и поисковый контекст}\label{sec:model-content}

\subsection{Раздельные пространства представлений}

Для товара $i$ обозначим предоставленные векторы описания и изображения
через $v_i^{\mathrm{text}}\in\R^{d_t}$ и
$v_i^{\mathrm{image}}\in\R^{d_v}$.
В обозначениях каталога $v_i^{\mathrm{text}}=d_i$,
$v_i^{\mathrm{image}}=v_i$.
Координаты двух пространств имеют разный смысл даже при $d_t=d_v$.
Документация Coveo указывает совместимость вектора запроса
с вектором описания товара; совместимость с изображением
не заявлена~\cite{coveo-data}.

Для каждой модальности $m\in\{\mathrm{text},\mathrm{image}\}$
вводится маска $M_i^m$. Она равна единице, если вектор существует,
имеет ожидаемую размерность, конечные координаты и норму больше
фиксированного $\eps_v>0$. Определим
\begin{equation}\label{eq:model-vector-normalization}
\bar v_i^m=
\begin{cases}
v_i^m/\norm{v_i^m},&M_i^m=1,\\
0,&M_i^m=0.
\end{cases}
\end{equation}
Отсутствующий вектор обозначается нулём и маской; нулевой косинус
корректных векторов остаётся наблюдаемым значением.

Представление истории строится отдельно в каждом пространстве:
\begin{align}\label{eq:model-session-vector}
u_q^m&=\sum_{\ell\in\mathcal D_q}
\alpha_{q\ell}M^m_{i_{s(q)\ell}}\bar v^m_{i_{s(q)\ell}},\\
M_q^m&=\ind\{\norm{u_q^m}>\eps_v\},\qquad
\bar u_q^m=\begin{cases}u_q^m/\norm{u_q^m},&M_q^m=1,\\
0,&M_q^m=0.\end{cases}\nonumber
\end{align}
Норма суммы проверяется из-за возможной компенсации направлений.
Содержательный балл и его маска равны
\begin{equation}\label{eq:model-content-score}
c^m_{qi}=(\bar u_q^m)^\top\bar v_i^m,\qquad
M^m_{qi}=M_q^m M_i^m.
\end{equation}
При $M^m_{qi}=1$ балл лежит в $[-1,1]$.
Для отбора допустима взвешенная смесь только доступных модальностей:
\begin{equation}\label{eq:model-content-fusion}
c_{qi}=
\begin{cases}
\displaystyle\frac{\sum_m\gamma_m M^m_{qi}c^m_{qi}}
 {\sum_m\gamma_m M^m_{qi}},&\sum_m\gamma_m M^m_{qi}>0,\\[0.3em]
0,&\text{иначе},
\end{cases}
\qquad \gamma_m\ge0.
\end{equation}
Ранжировщик получает отдельные баллы и маски.
Варианты сравниваются при одинаковых квотах кандидатов.

\subsection{Последний доступный поисковый запрос}

Пусть $b(q)$ --- последний поисковый блок, доступный к $t_q$,
а $t_{b(q)}$ --- его время. Если его нет либо
$t_q-t_{b(q)}>\Delta_S$, поисковый контекст считается отсутствующим.
Непригодность последнего блока по
разделу~\ref{sec:data-search-canonicalization} отключает поиск;
возврата к более раннему блоку нет.
Порог $\Delta_S$ задаётся в секундах. Корректный вектор запроса
нормируется по правилу~\eqref{eq:model-vector-normalization};
результат обозначается $\bar z_q\in\R^{d_t}$,
а маска доступности --- $M_q^S$.
\begin{equation}\label{eq:model-search-score}
s_{qi}^{\mathrm{search}}=\bar z_q^\top\bar v_i^{\mathrm{text}},\qquad
M_{qi}^S=M_q^S M_i^{\mathrm{text}}.
\end{equation}
$s_{qi}^{\mathrm{search}}$ измеряет близость запроса к описанию;
сравнение с графовым баллом требует обучения или нормировки.

Дополнительные признаки поискового контекста имеют вид
\begin{equation}\label{eq:model-search-features}
d_q^S=\begin{cases}
t_q-t_{b(q)},&M_q^S=1,\\
0,&M_q^S=0,
\end{cases}
\qquad M_q^S,\quad
a_q^S=\log\!\left(1+\frac{d_q^S}{t_{\mathrm{unit}}}\right).
\end{equation}
Величина $d_q^S$ измеряется в секундах, $a_q^S$ безразмерна.
При отсутствии поиска давность и маска равны нулю; отсутствующее
$t_{b(q)}$ в нулевой ветви не вычисляется.
Поле \code{product_skus_hash} исключено из признаков и отбора:
кандидаты выбираются по близости к \code{query_vector} из $\mathcal I_q$,
без копирования исторической поисковой выдачи.

\begin{invariant}[Обратная связь после поискового запроса]
Поле \code{clicked_skus_hash} не участвует в признаках текущего
поискового запроса: оно содержит последующие клики.
При отсутствии индивидуальных времён кликов поле не используется
и для восстановления состояния более позднего запроса.
Самостоятельное событие просмотра товара с известным временем
может учитываться по обычному правилу доступности.
\end{invariant}

Предсказание поисковых кликов допускается отдельным экспериментом
со своими метками. Отсутствие клика не означает отрицательной оценки
и не меняет цель покупки из раздела~\ref{sec:business}.

\subsection{Категории и ценовая группа}

Для категорий кандидата и просмотренного товара с длинами путей
$\ell_i^{\mathrm{cat}},\ell_j^{\mathrm{cat}}$ задаётся сходство
\begin{equation}\label{eq:model-category-similarity}
c_{ij}^{\mathrm{cat}}=
\frac{\operatorname{lcp}(i,j)}
 {\max(1,\ell_i^{\mathrm{cat}},\ell_j^{\mathrm{cat}})},
\end{equation}
где $\operatorname{lcp}$ --- число одинаковых начальных категорий.
При отсутствующем пути значение равно нулю и ставится отдельная маска.
Хеш служит категориальной меткой.

Для порядкового признака \code{price_bucket} используются совпадение,
знак и абсолютная разность номеров групп с масками пропусков.

\section{Формирование множества кандидатов}\label{sec:model-candidates}

\subsection{Источники и квоты}

Все источники ограничиваются $\mathcal I_q$ по разделу~\ref{sec:requests}:
уже купленные товары исключаются, просмотренные и наблюдаемые
в корзине сохраняются и оцениваются отдельными срезами.
Обозначим через $\mathcal C_q^{(a)}$ список источника $a$.
Первый запуск использует шесть групп источников:
\begin{enumerate}
\item последние различные товары из доступных просмотров;
\item соседи по каждому из трёх отношений графа;
\item товары с максимальным сходством в пространстве описаний;
\item товары с максимальным сходством в пространстве изображений;
\item товары допустимого каталога, близкие последнему поисковому запросу;
\item исторически популярные товары;
\end{enumerate}
Начальные квоты: по 50 кандидатов для каждого из трёх отношений графа,
40 по описаниям, 30 по изображениям, 50 по поиску,
20 недавних и 10 популярных. Их сумма равна \(M=300\).
Нейросетевой источник добавляется только на втором этапе
(раздел~\ref{sec:modern-retrieval}), источник семантических кодов ---
в отдельном исследовании (пункт~\ref{sec:sid-serving}).
Недоступный источник возвращает пустой список и код причины.
Содержательный источник включается только при корректном представлении
истории; произвольная сортировка всех отсутствующих баллов не применяется.

Квоты $m_a\ge0$ фиксируются до итоговой оценки:
\begin{equation}\label{eq:model-candidate-union}
\mathcal C_q^{\mathrm{pre}}=\bigcup_a\mathcal C_q^{(a)},\qquad
|\mathcal C_q^{(a)}|\le m_a,\qquad
\sum_a m_a\le M.
\end{equation}
Дубликаты объединяются с сохранением всех признаков источников.
Квота нейросетевого источника включает все подзапросы;
новый источник перераспределяет квоты внутри прежнего $M$.
Недостающие после объединения места заполняются по популярности
среди $\mathcal I_q\setminus\mathcal C_q^{\mathrm{pre}}$.
Пусть $\mathcal C_q^{\mathrm{fill}}$ --- первые не более
$M-|\mathcal C_q^{\mathrm{pre}}|$ таких товаров. Итоговый набор равен
\begin{equation}\label{eq:model-candidate-fill}
\mathcal C_q^{\mathrm{main}}=\mathcal C_q^{\mathrm{pre}}\cup\mathcal C_q^{\mathrm{fill}},
\qquad |\mathcal C_q^{\mathrm{main}}|\le M.
\end{equation}
Заполнение имеет отдельный признак и не ограничено начальной квотой
популярности. Оно прекращается при $M$ или исчерпании каталога.
Добавление будущего правильного ответа запрещено.

Это штатное множество кандидатов. До ранжирования отдельно готовится
$\mathcal C_q^{\mathrm{fallback}}$: до $M$ допустимых товаров, сначала
различные товары последних просмотров в порядке убывания времени
(при равенстве --- по идентификатору), затем популярность без повторов.
Резерв не объединяется со штатным набором сверх бюджета.
Нормативное множество фактически выбранной ветви равно
\begin{equation}\label{eq:model-served-candidates}
\mathcal C_q^{\mathrm{served}}=
\begin{cases}
\mathcal C_q^{\mathrm{main}},&\mathrm{branch}_q=\mathrm{main},\\
\mathcal C_q^{\mathrm{fallback}},&\mathrm{branch}_q=\mathrm{fallback},\\
\varnothing,&\mathrm{branch}_q\in\{\mathrm{empty},\mathrm{error}\}.
\end{cases}
\end{equation}
Ветвь фиксируется при выборе ответа с резервом времени на обязательную
запись и возврат. Если выбранная ветвь затем не успела вернуть ответ,
сохраняют её кандидатов, $L_q=\varnothing$ и отдельный код отказа:
потеря относится к выдаче, а не к отбору. Значение \code{error} в формуле
означает отказ до выбора готовой ветви. Во всех случаях
$L_q\subseteq\mathcal C_q^{\mathrm{served}}\subseteq\mathcal I_q$
и $|\mathcal C_q^{\mathrm{served}}|\le M$.
Далее $\mathcal C_q$ без верхнего индекса обозначает
$\mathcal C_q^{\mathrm{served}}$; при штатном обучении ранжировщика на $B$
это $\mathcal C_q^{\mathrm{main}}$. Резерв имеет свой фиксированный порядок
и не обучается на градиентах основного ранжировщика. Незавершённый штатный
набор регистрируется как неизмеренный, а не как пустой диагностический список.

Историческая популярность определяется количеством различных
сессий снимка с просмотром товара:
\begin{align}\label{eq:model-popularity}
V_i^{\mathrm{detail}}&=\sum_{s\in A_G}
\ind\{s\text{ содержит доступный просмотр }i\},\\
\pi_i^{\mathrm{pop}}&=\frac{V_i^{\mathrm{detail}}+\eta_P}
 {\sum_{j\in\mathcal I_G}(V_j^{\mathrm{detail}}+\eta_P)},
\quad\eta_P>0.
\nonumber\end{align}
При запросе популярность ограничивается $\mathcal I_q$.
Неизвестный товар может изменить состояние сессии и содержательное
представление, но не получает исторических рёбер. Допуск новых
товаров подчиняется правилу доступности раздела~\ref{sec:requests}.

\subsection{Поиск соседей в пространстве векторов}

Первый запуск вычисляет сходство со всеми допустимыми векторами:
стоимость модальности размерности $d_m$ имеет порядок $|\mathcal I_q|d_m$.
Если измеренная задержка чрезмерна, отдельно испытывается
приближённый поиск ближайших соседей (ANN).
Его настройки фиксируются до итоговой проверки; на отдельной выборке
запросов измеряется
\begin{equation}\label{eq:model-ann-agreement}
\mathrm{Agreement@}m=
\frac1{N_Q^+}\sum_{q:\,|\mathcal C_{q,\mathrm{exact}}|>0}
\frac{|\mathcal C_{q,\mathrm{approx}}\cap
\mathcal C_{q,\mathrm{exact}}|}
 {|\mathcal C_{q,\mathrm{exact}}|}.
\end{equation}
$N_Q^+$ --- число непустых точных ответов; при $N_Q^+=0$ показатель
не определён. Доля пустых ответов приводится отдельно.
Одинаковый фильтр допустимости позволяет измерить потери приближения,
но согласие с точным поиском не заменяет оценку покупки.

\subsection{Признаки происхождения кандидата}

Для каждого источника сохраняются его маска, исходный балл
и обратный ранг:
\begin{equation}\label{eq:model-source-rank}
o_{qi}^{(a)}=\ind\{i\in\mathcal C_q^{(a)}\},\qquad
u_{qi}^{(a)}=
\begin{cases}
1/(1+\operatorname{rank}_a(q,i)),&o_{qi}^{(a)}=1,\\
0,&o_{qi}^{(a)}=0.
\end{cases}
\end{equation}
Ранг начинается с единицы; пропуски ранга и балла обозначаются нулём
с маской. Измеримый балл вычисляется для всех кандидатов объединения,
даже если товар не вошёл в квоту соответствующего источника.
Преобразование $1/(1+\operatorname{rank})$ намеренно даёт $1/2$
первой позиции; это ограниченный признак, не вероятность.

\section{Признаки пары <<запрос --- товар>>}\label{sec:model-features}

Вектор $x_{qi}\in\R^p$ содержит только величины, доступные при $q$.
Преобразование категорий фиксируется вместе с моделью;
идентификатор товара не используется как вещественное число.
Таблица задаёт признаки первого запуска.

{\small
\begin{longtable}{p{0.23\textwidth}p{0.65\textwidth}}
\toprule
Группа & Содержание \\
\midrule
Граф & $g^r_{qi}$, связь с последней группой товаров, поддержка ребра,
маски отсутствующих соседей \\
Описание и изображение & $c^m_{qi}$, сходство с последней группой товаров,
маски векторов и пустого представления истории \\
Поиск & $s_{qi}^{\mathrm{search}}$, давность запроса,
маски доступности запроса и описания товара \\
История товара & Число доступных просмотров, давность последнего
просмотра, последнее действие добавления или удаления \\
История сессии & Число событий и различных товаров, длительность
доступной истории, наличие поиска, число зарегистрированных добавлений \\
Каталог & Категория, совпадение категорий, ценовая группа
и её отличие от групп последних просмотренных товаров \\
Снимок & Популярность и число сессий с зарегистрированной покупкой товара \\
Отбор & Маски источников, их исходные баллы и обратные ранги \\
\bottomrule
\end{longtable}}

На втором этапе добавляются балл обученного вектора сессии,
максимальный балл намерений, их дисперсия, энтропия весов,
поисковое сходство в обученном пространстве и маски доступности.
В исследовании семантических кодов добавляются логарифм вероятности
кода, число товаров с этим кодом и маска найденного кода.
Вероятность кода не является вероятностью покупки.
При изменении набора признаков ранжировщик обучается заново на \(B\);
нейросетевые признаки не нужны для реализации первого этапа.

Множество $\mathcal I_q^{\mathrm{last}}$ содержит все различные товары
последней временной группы \code{detail}, без выбора одного последнего.
Вес $n_{qj}^{\mathrm{last}}$ равен $1$ в \code{set_v1} и кратности
просмотра $j$ в этой группе в \code{multiset_v1}, после удаления
транспортных повторов. Для попарного признака $\phi(i,j)$ с маской
доступности $M_{ij}^{\phi}$ среднее и его маска равны
\begin{align}\label{eq:model-last-group-mean}
Z_{qi}^{\phi}&=\sum_{j\in\mathcal I_q^{\mathrm{last}}}
 n_{qj}^{\mathrm{last}}M_{ij}^{\phi},\qquad
M_{qi}^{\mathrm{last},\phi}=\ind\{Z_{qi}^{\phi}>0\},\nonumber\\
\overline\phi_{qi}&=\begin{cases}
\displaystyle\frac{\sum_{j\in\mathcal I_q^{\mathrm{last}}:\,M_{ij}^{\phi}=1}
 n_{qj}^{\mathrm{last}}\phi(i,j)}{Z_{qi}^{\phi}},&Z_{qi}^{\phi}>0,\\
0,&Z_{qi}^{\phi}=0.
\end{cases}
\end{align}
Для описания и изображения $\phi(i,j)$ --- косинус корректных векторов,
$M_{ij}^{\phi}=M_i^mM_j^m$. Для ценовых групп усредняются совпадение,
$\operatorname{sign}(b_i^{\mathrm{price}}-b_j^{\mathrm{price}})$
и абсолютная разность; маска требует обоих известных значений.
Недоступная $\phi(i,j)$ не вычисляется.
Графовый максимум $\max_{j\in\mathcal I_q^{\mathrm{last}}}\rho_j^rP^r_{ji}$
и максимальное сходство категорий от кратности не зависят;
пустой допустимый набор даёт ноль и маску $0$.

Количество просмотров товара $i$ и давность определяются как
\begin{align}\label{eq:model-online-features}
n_{qi}^{\mathrm{detail}}&=
\sum_{e\in\mathcal H_q}\ind\{i_e=i,\ \kappa_e=\mathrm{detail}\},\\
\delta_{qi}^{\mathrm{detail}}&=
t_q-\max\{t_e:e\in\mathcal H_q,\ i_e=i,\ \kappa_e=\mathrm{detail}\}.
\nonumber
\end{align}
При отсутствии события давность и маска равны нулю.
Для числовой устойчивости допустимы преобразования счётчика
$\log(1+n)$ и давности $\log(1+\delta/t_{\mathrm{unit}})$.

Состояние корзины и его давность определены в разделе~\ref{sec:time}.

Сглаженная покупочная частота не используется: без подтверждённой
полноты наблюдений её знаменатель не установлен. Этот признак исключён
из схемы, а не заполнен нулём. Счётчики записанных покупок сохраняются.

Статистики используют только $A_G$; целевое кодирование по меткам
периода обучения ранжировщика исключено. Новая категория получает
отдельный код <<неизвестная категория>>.

\section{Обучение ранжированию по будущим покупкам}\label{sec:model-ranking}

\subsection{Метки, группы и веса}

Ранжировщик обучается на первых запросах сессий $\mathcal S_B^+$
с зарегистрированной целью (раздел~\ref{sec:eval-observed}).
Модель вычисляет балл $f(x_{qi})$ при метке
$y_{qi}=\ind\{i\in Y_q\}$, $i\in\mathcal C_q$.
Запрос $q$ образует группу сравнения. Нулевая метка означает
отсутствие наблюдавшейся покупки в окне, не отрицательное отношение.

Положительные и остальные кандидаты образуют множества
\begin{equation}\label{eq:model-pair-sets}
\mathcal C_q^+=\{i\in\mathcal C_q:y_{qi}=1\},\qquad
\mathcal C_q^-=\{i\in\mathcal C_q:y_{qi}=0\},\qquad
\mathcal J_q=\mathcal C_q^+\times\mathcal C_q^-.
\end{equation}
Одна сессия даёт одну группу. Веса равны
\begin{equation}\label{eq:model-ranking-weights}
w_q^{\mathrm{sess}}=1,\qquad
w_q^{\mathrm{pair}}=\frac{w_q^{\mathrm{sess}}}{\max(1,|\mathcal J_q|)}.
\end{equation}
После умножения на $\Delta_{qij}$ сумма весов пар
и норма градиентов зависят от состава кандидатов и текущих рангов;
равенство фактического вклада разных сессий не утверждается.
Группа без положительного товара не даёт парного градиента,
но сессия с зарегистрированной целью сохраняется в знаменателе оценки.
Целевой товар не добавляют в кандидаты. При отсутствии обучающих пар
сообщают недостаточность данных, а не успешное обучение.

\begin{assumption}[Суррогатное сравнение при неявной обратной связи]
\label{ass:model-implicit-negatives}
Каждый наблюдавшийся купленный кандидат обучается превосходить
каждого кандидата без покупки в том же окне. Это допущение о полезном
суррогате исторической задачи, а не утверждение об экспозиции
или предпочтении: некупленный товар мог вообще не быть показан.
Историческая политика, доступность и популярность влияют на метки;
без журнала показов смещение отбора этим обучением не устраняется.
\end{assumption}
Из этого допущения не следует обязательная замена задачи
на positive--unlabeled обучение: его идентифицирующие предположения
также потребовали бы обоснования. Выбор политики проверяется
независимым историческим сравнением и затем A/B-протоколом.

\subsection{Парная логистическая потеря}

Для положительного $i$ и отрицательного $j$ положим
$d_{qij}=f(x_{qi})-f(x_{qj})$.
Базовая функция обучения имеет вид
\begin{equation}\label{eq:model-pairwise-loss}
\mathcal L_{\mathrm{pair}}(f)=
\sum_q w_q^{\mathrm{pair}}\sum_{(i,j)\in\mathcal J_q}
\log\!\left(1+\exp(-\beta_R d_{qij})\right),
\qquad \beta_R>0.
\end{equation}
Параметр $\beta_R$ задаёт масштаб баллов.
Вычисляется устойчивая $\operatorname{softplus}(-\beta_Rd_{qij})$,
а не непосредственная экспонента большого положительного аргумента.
При $d_{qij}=0$ производная по положительному баллу равна
$-w_q^{\mathrm{pair}}\beta_R/2$, поэтому минимизация повышает положительный балл
и понижает отрицательный. Логистическое преобразование разности
описывает порядок пары, а не вероятность покупки товара.

\subsection{Повышенный вес ошибок в верхней части списка}

Для LambdaMART вес пары зависит от того, насколько её перестановка
изменила бы качество верхних позиций~\cite{lambdamart}.
Пусть $r_{qi}$ --- текущая позиция $i$ среди всех кандидатов,
а $K_R$ --- глубина рангового обучения. Ранги начинаются с единицы;
порядок задаётся ключом $(-f(x_{qi}),\operatorname{id}(i))$,
идентификаторы сравниваются лексикографически по байтам UTF-8.
Для бинарных меток используется gain $G_R(0)=0$, $G_R(1)=1$.
Определим
\begin{equation}\label{eq:model-rank-discount}
D_R(k)=\frac{\ind\{k\le K_R\}}{\log_2(1+k)},\qquad
Z_q=\sum_{k=1}^{\min(K_R,|\mathcal C_q^+|)}D_R(k).
\end{equation}
При $Z_q>0$ нормированное качество списка равно
$\sum_{i\in\mathcal C_q}y_{qi}D_R(r_{qi})/Z_q$.
Его изменение при перестановке положительного и отрицательного
товара по модулю равно
\begin{equation}\label{eq:model-lambda-weight}
\Delta_{qij}=
\frac{|D_R(r_{qi})-D_R(r_{qj})|}{Z_q},\qquad
\omega_{qij}=w_q^{\mathrm{pair}}\Delta_{qij}.
\end{equation}
При $Z_q=0$ пары отсутствуют и градиенты равны нулю.
Пара из двух позиций ниже $K_R$ имеет нулевой вес. Если ровно одна
позиция выше границы, используется полный модуль разности скидок;
никакой дополнительной обрезки пары или нормировки по разности
баллов нет.

Веса $\omega_{qij}$ пересчитываются по текущим позициям
и фиксируются при вычислении производных. Это правило псевдоградиентов
не предполагает единой гладкой функции со всеми этими производными.
При $\Delta_{qij}=1$ получается парное логистическое обучение.

Для одной пары обозначим
$p_{qij}=1/(1+\exp(\beta_Rd_{qij}))$.
Вклады в первую производную и диагональную вторую производную:
\begin{align}\label{eq:model-lambda-derivatives}
g_{qi}^{R,(ij)}&=-\omega_{qij}\beta_Rp_{qij},&
g_{qj}^{R,(ij)}&=+\omega_{qij}\beta_Rp_{qij},\\
h_{qi}^{R,(ij)}&=\omega_{qij}\beta_R^2p_{qij}(1-p_{qij}),&
h_{qj}^{R,(ij)}&=h_{qi}^{R,(ij)}.\nonumber
\end{align}
Производные суммируются по парам кандидата. Деревья используют
диагональное приближение, опуская смешанные вторые производные.
Сумма градиентов группы равна нулю: абсолютный уровень баллов не задан.
Для численной реализации положим $u=\beta_Rd_{qij}$,
$v=\exp(-|u|)$. Тогда $p=v/(1+v)$ при $u\ge0$ и $p=1/(1+v)$
при $u<0$, а $p(1-p)$ вычисляется как $v/(1+v)^2$.
Это сохраняет малую кривизну без вычитания близких чисел;
дополнительный положительный floor для $h$ не вводится.

\subsection{Ансамбль деревьев и регуляризация}

Ранжировщик представляет собой сумму деревьев:
\begin{equation}\label{eq:model-boosted-trees}
f_{N_T}(x)=f_0+\nu\sum_{m=1}^{N_T} T_m(x),\qquad 0<\nu\le1.
\end{equation}
Лист дерева выдаёт одно число; $\nu$ ограничивает вклад нового дерева.
Константа $f_0$ не влияет на порядок.

Для листа $\ell$ очередного дерева обозначим
$G_\ell=\sum_{(q,i)\in\ell}g_{qi}^{R}$ и
$H_\ell=\sum_{(q,i)\in\ell}h_{qi}^{R}$.
При штрафе $\lambda_2 w_\ell^2/2$ и $\lambda_2>0$
локально оптимальное значение листа равно
\begin{equation}\label{eq:model-newton-leaf}
w_\ell^*=-\frac{G_\ell}{H_\ell+\lambda_2}.
\end{equation}
Формула получается минимизацией квадратичного выражения
$G_\ell w+(H_\ell+\lambda_2)w^2/2$.
Уменьшение этого приближения при разделении листа на $L,R$:
\begin{equation}\label{eq:model-tree-gain}
\mathrm{Gain}=\frac12\left[
\frac{G_L^2}{H_L+\lambda_2}+
\frac{G_R^2}{H_R+\lambda_2}-
\frac{(G_L+G_R)^2}{H_L+H_R+\lambda_2}\right]-\gamma_T.
\end{equation}
$\gamma_T\ge0$ --- штраф за добавление листа.
Разделение допускается при положительном выигрыше и выполнении
ограничений на глубину, число листьев, число наблюдений
и сумму вторых производных в каждом дочернем листе.
В первом запуске фиксирован бюджет $200$ итераций; фактическое
$N_T\le200$ сохраняется в паспорте, поскольку отсутствие допустимого
разбиения может остановить построение новых деревьев раньше.
Позднейший выбор бюджета выполняется на \(C_{\mathrm{select}}\) по заранее
заданному плану; после выбора модель замораживается до итоговой проверки.

\subsection{Исполнимый контракт собственных производных}
\label{sec:model-ranker-implementation}

Используется Python API \code{lightgbm.train} версии
\code{lightgbm==4.6.0} с вызываемым
\code{objective=session\_lambdas\_v1}~\cite{lightgbm}.
Функция принимает сырые баллы и \code{Dataset}, возвращает
суммы $g_{qi}^R,h_{qi}^R$ из~\eqref{eq:model-lambda-derivatives}.
При каждой итерации заново вычисляются полные ранги, $Z_q$
и все пары из~\eqref{eq:model-pair-sets}.
Это собственная LambdaMART objective; эквивалентность встроенной
\code{objective=lambdarank} не заявляется. Параметры
\code{label\_gain}, \code{lambdarank\_norm},
\code{lambdarank\_truncation\_level} и \code{sigmoid}
встроенной objective здесь не задаются и не управляют callback:
gain $[0,1]$, $K_R=20$, $\beta_R=1$ реализованы в нём явно.

Строки группируют по запросу, запросы упорядочивают по
$(s,t_q,\operatorname{id}(q))$, внутри группы --- по идентификатору
товара; все порядки входят в схему данных. Размеры непустых групп
передают как \code{Dataset.group}. Пустая группа не создаёт строк,
но остаётся в $|\mathcal Q_s|$ для весов. Веса
$w_q^{\mathrm{pair}}$ применяются один раз внутри callback;
\code{Dataset.weight=None}, дополнительных row/query weights нет.
Группы без пар получают нулевые массивы производных.
Суммы считают в FP64, затем callback возвращает непрерывные массивы
FP32: такой тип принимает C API выбранной версии. Не конечные
баллы, градиенты или гессианы останавливают обучение с ошибкой,
а не заменяются нулями.

Полная конфигурация первого ранжировщика ниже дополняется только
фиксированной схемой признаков и данными; неуказанные параметры
имеют значения по умолчанию именно версии 4.6.0. Сериализованный паспорт
содержит значения по умолчанию, хеш функции производных, файл версий
с хешами пакетов и сведения о платформе/компиляторе. Новая версия
библиотеки требует повторной проверки контракта.
\begin{verbatim}
objective=session_lambdas_v1 (callable); metric="None"
boosting="gbdt"; data_sample_strategy="bagging"
device_type="cpu"; tree_learner="serial"; num_threads=1
deterministic=true; force_col_wise=true; force_row_wise=false
learning_rate=0.05; num_leaves=31; max_depth=8
lambda_l1=0; lambda_l2=1; min_gain_to_split=0
min_data_in_leaf=100; min_sum_hessian_in_leaf=0.001
bagging_fraction=1; bagging_freq=0; bagging_by_query=false
feature_fraction=1; feature_fraction_bynode=1; extra_trees=false
max_bin=255; min_data_in_bin=3; bin_construct_sample_cnt=200000
feature_pre_filter=false; enable_bundle=false
use_missing=true; zero_as_missing=false; use_quantized_grad=false
boost_from_average=false; linear_tree=false
path_smooth=0; max_delta_step=0
cat_l2=0; cat_smooth=1; max_cat_threshold=32
max_cat_to_onehot=4; min_data_per_group=100
seed=17; data_random_seed=17; feature_fraction_seed=17
bagging_seed=17; extra_seed=17; verbosity=1
train: num_boost_round=200; valid_sets=None; feval=None
       init_model=None; callbacks=None
Dataset: weight=None; init_score=None
predict: raw_score=true; num_iteration=N_T
\end{verbatim}
Исходный балл $f_0=0$. Категория --- единственный нативный
категориальный столбец первого запуска: его имя явно передаётся
в \code{categorical\_feature}; словарь строится по $A_G$ в порядке
UTF-8, неизвестное значение имеет код $0$, известные --- $1,2,\ldots$.
Ценовая группа является порядковым числом с отдельной маской.
Набор, порядок и типы остальных числовых признаков сохраняются
в схеме; идентификатор товара не является признаком.
Монотонные, interaction и forced-split ограничения отсутствуют.

LightGBM строит гистограммные разбиения, поэтому перебор всех
вещественных порогов не предполагается. Его \code{min\_data\_in\_leaf}
проверяется приближённо через гессианы: это параметр строк,
не запросов и не строгая гарантия $100$ фактических строк.
По маршрутизации обучающих строк публикуются фактические размеры
листьев; требование строгого минимума потребовало бы отдельного
алгоритма и нового эксперимента. При отключённых L1, сглаживании
листа \code{path_smooth} и ограничении \code{max_delta_step}
выигрыш LightGBM имеет масштаб
$2(\mathrm{Gain}+\gamma_T)$ до вычитания порога, поэтому
\code{min\_gain\_to\_split}$=2\gamma_T$; в первом запуске оба равны
нулю. \code{cat_l2=0} исключает дополнительный категориальный штраф
листа, отсутствующий в~\eqref{eq:model-newton-leaf}.
Положительное проектное \code{cat_smooth=1} действует иначе:
при группировке категорий оно ограничивает допуск bins и сглаживает
ключ сортировки $G_{\mathrm{bin}}/(H_{\mathrm{bin}}+1)$.
Поэтому разрешённый $G_{\mathrm{bin}}=H_{\mathrm{bin}}=0$ не даёт
$0/0$ в ключе сортировки; значения callback и формула листа не меняются.
Этот параметр нельзя выключать вместе со штрафом листа.
Малые внутренние численные константы
библиотеки учитываются допуском теста, а не равенством битов.

Приёмочный тест сравнивает все элементы $g,h$ с независимым FP64
расчётом: для callback до приведения допускается
$10^{-12}+10^{-10}|v_{\mathrm{ref}}|$, после FP32 ---
$10^{-7}+10^{-6}|v_{\mathrm{ref}}|$. Проверяются несколько групп,
две положительные метки, равные баллы, $K_R<|\mathcal C_q|$,
нулевые группы, разное число кандидатов в группе и логиты $\pm10^4$.
Отдельный одношаговый тест на числовом признаке с однозначным
разбиением сверяет листья и удвоенный gain с допуском $10^{-6}$;
для него ограничения размера ослабляются явно только в fixture.
Перестановка входных строк после канонизации должна сохранять
производные. Умножение на веса второй раз должно обнаруживаться
тестом, а встроенная LambdaRank не используется как эталон.

\subsection{Отрицательные примеры и разреженная покупка}

Основная версия использует всех кандидатов без выборки отрицательных.
Для проверки чувствительности заранее регистрируют отдельный
вариант с перевзвешенными некупленными кандидатами. Для каждого
$i\in\mathcal C_q^+$ вместо равномерного веса $1/|\mathcal C_q^-|$
он использует смесь $\tfrac12u_q(j)+\tfrac14v_q(j)+\tfrac14h_{qi}(j)$:
$u_q$ равномерно по $\mathcal C_q^-$, $v_q$ пропорционально
$1+N_j^{A_G}$ (числу сессий $A_G$ с товаром), а $h_{qi}$ равномерно
по первым $\min(20,|\mathcal C_q^-|)$ некупленным товарам по
фиксированному сходству категории и описания с $i$ (сначала
совпадение категории, затем cosine, затем идентификатор; недоступное
сходство равно нулю). Полный вес пары такого варианта равен
$w_q^{\mathrm{sess}}\Delta_{qij}$, умноженному на эту смесь
и делённому на $|\mathcal C_q^+|$. Это другая суррогатная objective,
а не исправление экспозиции. Вариант обучают на $B$ и сравнивают
на $C_{\mathrm{select}}$ в пределах зарегистрированного бюджета.
Для обоих вариантов публикуются Hit@20 и качество ранжирования
по head/tail, определённым только популярностью $A_G$
(head --- первые $\lceil0.1|\mathcal I_G|\rceil$ товаров
по убыванию $N_j^{A_G}$, при равенстве по ID,
tail --- остальные товары). Сравнивают качество политик,
а не численные значения разных функций потерь.
При вынужденном сокращении данных пара $(i,j)$ включается
с известной положительной вероятностью $\xi_{qij}$.
Для фиксированного текущего порядка несмещённый оцениватель
суммы локальных градиентов использует вес
\begin{equation}\label{eq:model-pair-sampling}
\widetilde\omega_{qij}=
\frac{\ind\{(i,j)\text{ отобрана}\}}{\xi_{qij}}
\omega_{qij}.
\end{equation}
$\xi_{qij}$ относится к паре; ранги и $Z_q$ берутся по полной группе.
Их пересчёт по выборке не восстанавливает исходный LambdaMART.
Малые $\xi_{qij}$ повышают дисперсию; отсечение обратных весов
вносит смещение и требует отдельного обозначения.

Отсутствующие покупки не восстанавливаются вымышленными метками.
Баллы вспомогательных задач следующего просмотра и добавления
допустимы как признаки при обучении без временной утечки.

Отдельное будущее исследование может предсказывать покупку в сессии
по доступной истории и кандидатам. Один скалярный прогноз для запроса
сам по себе не сравнивает товары или разные списки: прибавление
общей величины к баллам сохраняет их порядок. Такой прогноз пригоден
для диагностики или вспомогательного обучения, но не оценивает
причинную полезность выдачи и не входит в первый запуск.

\subsection{Ограничения интерпретации баллов}\label{sec:model-calibration}

Калибровка отключена, баллы не являются вероятностями покупки.
Период $E$ остаётся зарезервированным; он не дополняет обучение
или выбор модели. Вероятностные способы формирования выдачи
требуют отдельной постановки и подходящих данных.

\section{Формирование окончательной выдачи}\label{sec:model-slate}

\subsection{Основной способ выдачи}

Базовая политика и нейросетевое расширение сортируют $\mathcal C_q$ по $f(x_{qi})$
и возвращают первые $K_q=\min(K,|\mathcal C_q|)$ различных товаров.
Калибровка с $a_{\mathrm{cal}}>0$ этот порядок не меняет.
При равенстве баллов используется зафиксированный порядок
идентификаторов. В первом запуске \(K=20\); во всех сравнениях
одинаковы верхний предел $K$ и правило длины, а не фактическое $K_q$.
Короткие и пустые ответы учитываются без фиктивного дополнения.

В первом запуске и начальном нейросетевом варианте \(\lambda_D=0\):
штраф за сходство товаров не применяется; калибратор отключён.
Изменение разнообразия списка и вероятностные ограничения рассмотрены в приложении~\ref{sec:model-diversity-optional}.

\subsection{Резервный ответ}

Если ранжировщик недоступен, используется отдельная резервная
политика: последние допустимые товары и историческая популярность
в порядке $\mathcal C_q^{\mathrm{fallback}}$ из
\eqref{eq:model-served-candidates}; выдаются первые $K$ товаров.
Если допустимых товаров нет,
возвращается пустой список. Все ограничения допустимости
сохраняются при отказе; код причины входит в журнал запроса.
Калибратор основной модели не применяется к резервным баллам.
Воздержание от показа и изменение $K$ образуют отдельные политики;
их причинную полезность проверяют по разделу~\ref{sec:eval-online},
при необходимости с контролем без блока.
Объяснения выдачи определены в разделе~\ref{sec:explanation}.


\section{Временной протокол обучения и проверки}\label{sec:eval-temporal}

Протокол един для основной системы и нейросетевого расширения:
вход содержит доступную историю, цели --- последующее поведение.
В первом запуске период обучения нейросети остаётся зарезервированным.

\subsection{Периоды и их назначение}\label{sec:eval-periods}

Пусть $A_G,A_N,B,C,E,D$ --- непересекающиеся календарные интервалы,
расположенные в указанном порядке, причём $A=A_G\cup A_N$:
\begin{equation}\label{eq:eval-period-order}
\sup A_G<\inf A_N,\quad \sup A_N<\inf B,\quad
\sup B<\inf C,\quad \sup C<\inf E,\quad \sup E<\inf D.
\end{equation}
Разрывы и потери у границ определяются горизонтом целей и запасом времени,
а не наблюдаемым значением целевой переменной.
Период $C$ заранее делят на два непересекающихся подинтервала:
\begin{equation}\label{eq:eval-validation-split}
C=C_{\mathrm{smoke}}\cup C_{\mathrm{select}},\qquad
C_{\mathrm{smoke}}\cap C_{\mathrm{select}}=\varnothing,\qquad
\sup C_{\mathrm{smoke}}\le\inf C_{\mathrm{select}}.
\end{equation}
Шесть крупных периодов после разделения $C$ образуют семь частей,
границы которых фиксируются в паспорте.
Календарный допуск~\eqref{eq:eval-admission} применяют к каждому
подинтервалу отдельно: цели отладки не пересекают границу выбора.
Во всём документе выбор параметров <<на $C$>> означает только
$C_{\mathrm{select}}$, а проверка реализации --- $C_{\mathrm{smoke}}$.

{\small
\begin{longtable}{L{0.10\textwidth}L{0.30\textwidth}L{0.47\textwidth}}
\toprule Период & Назначение & Допустимое использование \\\midrule\endhead
$A_G$ & Топология и справочники &
Граф, пути, каталог $\mathcal I_G$, популярности, счётчики записанных покупок,
масштабы исторических признаков. \\
$A_N$ & Обучение нейросетевого отбора &
Префиксы, зарегистрированные будущие покупки, вспомогательные следующие просмотры;
обучение графового кодировщика, содержательных проекций, намерений и поиска.
В основной системе не используется. \\
$B$ & Обучение ранжирования &
Кандидаты и признаки уже зафиксированного отбора;
подгонка деревьев по последующим покупкам. \\
$C$ & Отладка и независимый от неё выбор &
$C_{\mathrm{smoke}}$: проверка реализации и исправления;
$C_{\mathrm{select}}$: однократный выбор из замороженного списка вариантов
(размеры, эпохи, веса потерь, квоты, число деревьев). \\
$E$ & Резерв & Не используется в текущей постановке. \\
$D$ & Итоговая проверка &
Одна независимая оценка замороженной конфигурации и объявленных сравнений. \\
\bottomrule
\end{longtable}}

Начало и конец каждого периода, включая оба подинтервала $C$,
фиксируют в паспорте запуска до обучения. В первом запуске $A_N$ и $E$
не присоединяют к другим периодам и не используют для справочников.

\subsection{Порядок первого запуска и итоговой проверки}
\label{sec:eval-stages}

\begin{enumerate}[leftmargin=*]
\item \textbf{Проверка реализации.}
Запускают граф, содержательные признаки, поиск и LambdaMART
с бюджетом $200$ итераций без подбора параметров. На $C_{\mathrm{smoke}}$ проверяют кандидатов,
признаки, ранжирование и метрики; измеряют время ответа и память,
исправляют ошибки. $C_{\mathrm{select}},E,D$ остаются закрытыми.
Это локальные измерения; нагрузочная приёмка сервиса не требуется.
Полный вариант без метаданных не блокирует первый
$C_{\mathrm{smoke}}$; его обязательные сроки заданы ниже.
После отладки фиксируют код, признаки, метрики, список конфигураций,
случайные начальные состояния и правило разрешения равенств.
\item \textbf{Выбор окончательного варианта.}
Затем на $C_{\mathrm{select}}$ сравнивают конфигурации по замороженному плану
и вычислительному бюджету. Результат не порождает новые варианты.
Изменение реализации, признаков, метрик или списка после просмотра
$C_{\mathrm{select}}$ переводит его в данные разработки; повторный
подтверждающий выбор требует нового, ещё не просмотренного периода.
При отсутствии такого периода результат остаётся исследовательским.
Пример плана приведён
в приложении~\ref{app:ablation}; он не входит в первый запуск.
\item \textbf{Независимая оценка.}
До открытия $D$ отдельно фиксируют весь план итоговой оценки,
проверяют его связь с выбранной моделью и однократно открывают $D$.
$E$ остаётся закрытым.
\end{enumerate}

Нейросетевое исследование допускается после проверки данных, формул,
признаков и затрат базы; доказанный прирост покупок для этого не нужен.
Обучение каждого варианта соблюдает роли периодов из таблицы выше;
$C_{\mathrm{select}}$ выбирает результаты и сохранённые состояния,
но не предоставляет градиенты.

Фоновый каталог обучения, нормировок и реконструкции ограничен
$\mathcal I_G$. Признаки неизвестного товара могут участвовать в обучении
проекций лишь тогда, когда этот товар уже наблюдается в префиксе $A_N$.
Его будущие взаимодействия не расширяют каталог кандидатов.
Для всех предоставленных векторов сохраняется оговорка
из предположения~\ref{ass:static-content}.
После $A_N$ фиксируются веса и индекс векторов. Изменение любого
из них требует повторного построения признаков $B$ и обучения ранжировщика.

Для семантических идентификаторов сначала фиксируют кодировщик,
обученный на $A_N$, затем на том же периоде обучают дополнительный
модуль. Результат на $A_N$ не является независимой оценкой.

При смене отбора деревья заново обучают на $B$.
Новые события последующих периодов меняют только
состояние текущей сессии, не справочники.

Изменение метода по результатам $D$
переводит дальнейшие оценки на $D$ в исследовательские.

\subsection{Календарный допуск и полнота чтения}\label{sec:eval-maturity}

Горизонт $H$ и запас $g$ заданы в разделе~\ref{sec:requests}.
Для любого блока $J=[a_J,b_J)$ применяется единое правило:
\begin{equation}\label{eq:eval-admission}
t_{0,s}\in J,\qquad t_{0,s}+H+g<b_J.
\end{equation}
Правило одинаково для сессий с записанной покупкой и без неё;
метка следующего просмотра также не выходит за границу.
После допуска выполняют отбор ID по разделу~\ref{sec:data}.
Все опубликованные строки поведения и поиска выбранной сессии
читают целиком; части сессии не переходят между периодами.
Пропуск опубликованных строк --- ошибка подготовки.
Полнота чтения выгрузки не подтверждает полноты телеметрии.

В $A_G$ дополнительно требуется
$t_{\mathrm{last},s}^{\mathrm{obs}}+g<b_{A_G}$:
сессию с поздним хвостом исключают целиком независимо от типа события.
История включает выбранные сессии как с покупкой, так и без неё.
Каталог и статистики не расширяются событиями последующих периодов.
Ограничение по времени не доказывает фактического закрытия сессии.

Пусть $N_J^{\mathrm{eligible}}$ --- число сессий с подходящим первым
запросом в $J$ до календарных исключений и отбора ID, а
$N_J^{\mathrm{obs}}$ --- число сохранённых. Публикуют оба числа,
причины исключений и долю сохранённых сессий:
\begin{equation}\label{eq:eval-retention}
\operatorname{Retention}(J)=
\frac{N_J^{\mathrm{obs}}}{N_J^{\mathrm{eligible}}}.
\end{equation}
При пустом знаменателе отношение не определено.

\begin{invariant}[Временная доступность при фиксированной модели]
\label{inv:eval-no-leakage}
В запросах $B,C,D$ признаки и кандидаты
зависят только от доступного префикса, опубликованного до запроса
пакета отбора и оговорённой статической информации.
Изменение событий после запроса при фиксированном пакете
не меняет эти величины. $C_{\mathrm{smoke}}$ служит отладке,
$C_{\mathrm{select}}$ --- выбору; независимую итоговую
оценку выбранной модели проводят только на $D$.
При обучении на $A_N$ и $B$ будущие события разрешены
в целевой переменной, но не во входе того же примера.
\end{invariant}

Обученные веса не считаются доступными до запросов $A_N$.
Нормировки, распределения отрицательных примеров и контрастивные пары
также имеют объявленный период происхождения.

\subsection{Версии модели и окончательная фиксация}
\label{sec:eval-snapshots}

Основная система использует один фиксированный пакет. В отложенном
режиме обновления графа (приложение~\ref{sec:freshness}) пакеты
публикуют в моменты $r_1<r_2<\cdots$; запрос использует последний:
\begin{equation}\label{eq:eval-snapshot}
r(q)=\max\{r_j:r_j<t_q\}.
\end{equation}
Сохраняют времена отсечения графа, обучения кодировщика и публикации.
Признаки $B$ воспроизводят предусмотренный порядок версий;
новый кодировщик или индекс нельзя присвоить старому запросу.

До готовности обновления используют предыдущий пакет, регистрируя
его возраст. В подтверждающем эксперименте обновления запрещены
либо полностью заданы его протоколом.

До открытия $D$ фиксируются данные, интервалы, правило запросов,
$H,Q_{\max},K,M$, обработка отказов, все сравниваемые конфигурации,
основная метрика и построение интервала. Бюджет выбора включает
число конфигураций и случайных инициализаций, время обучения
и пиковую память. Одинаковое число кандидатов не означает
одинаковую вычислительную стоимость.

\subsection{Одно основное итоговое сравнение}
\label{sec:eval-primary-comparison}

Если исследуется улучшение основной системы, до открытия $D$
выбирают на $C_{\mathrm{select}}$ одну конфигурацию $\pi_1$ и фиксируют
контроль $\pi_0$: граф, признаки товаров, поиск и ранжировщик LambdaMART.
Основной исторический результат --- парная разность
$\widehat\Delta_{\mathrm{Hit}}(\pi_1,\pi_0)$ на $D$ и её интервал.
Каждый вариант имеет собственный ранжировщик, обученный на $B$;
совпадают допустимые данные, расписание запросов и бюджет кандидатов.
Выбор на $C_{\mathrm{select}}$ не продолжается после просмотра $D$.

Остальные варианты и разрезы описательны. Для нескольких подтверждающих
утверждений заранее задают $m_{\mathrm{hyp}}$ нулевых гипотез,
корректные способы вычисления $p$-значений и общий уровень $\alpha$.
Например, правило Холма \cite{holm1979} упорядочивает
$p_{(1)}\le\cdots\le p_{(m_{\mathrm{hyp}})}$ и последовательно отвергает
гипотезы, пока
\begin{equation}\label{eq:eval-holm}
 p_{(j)}\le\frac{\alpha}{m_{\mathrm{hyp}}-j+1};
\end{equation}
при первом нарушении остальные гипотезы не отвергаются.
Правило допускает зависимость корректных $p$-значений, но не исправляет
ошибки их расчёта, включая игнорирование зависимых сессий пользователя.
Если реализована только основная система, на $D$ оценивают её охват
и затраты без утверждения о превосходстве.


\section{Оценка на историческом журнале}\label{sec:eval-offline}

Исторический журнал позволяет установить, какие последующие покупки
попали бы в рассчитанные списки. Он не сообщает, изменились бы покупки
после реального показа этих списков. Поэтому основная историческая
метрика измеряет охват наблюдавшихся покупок, а причинная бизнес-метрика
определена отдельно в разделе~\ref{sec:business}.

\subsection{Совокупность сессий и основная метрика}
\label{sec:eval-observed}\label{sec:eval-schedule}

Протокол \code{observed_purchases_v1} оценивает единственную выдачу
$q=q_s^{(1)}$ с $K=20$ по зарегистрированной цели
$Y_q=Y_s^{\mathrm{obs}}$ из~\eqref{eq:request-target}.
Пусть $\mathcal S_J^{\mathrm{obs}}$ --- сессии периода $J$, прошедшие
календарный допуск и отбор по разделу~\ref{sec:eval-maturity}.
Их число $N_J^{\mathrm{obs}}=|\mathcal S_J^{\mathrm{obs}}|$,
границы и размеры подвыборок фиксируют до анализа будущих покупок;
недобор не восполняют покупающими сессиями. Для обучения покупочной
цели и оценки определим
\begin{equation}\label{eq:eval-observed-cohort}
 \mathcal S_J^+=\{s\in\mathcal S_J^{\mathrm{obs}}:
 Y_s^{\mathrm{obs}}\ne\varnothing\},\qquad N_J^+=|\mathcal S_J^+|.
\end{equation}
Отбор $\mathcal S_J^+$ не передаётся в кандидаты, признаки или решение
о выдаче. Исторические справочники строят по всем сохранённым сессиям
$A_G$, не только по покупающим. Расход ресурсов и отказы измеряют
также на сессиях без зарегистрированной цели: при запросе исход неизвестен.

Для фиксированной политики $\pi$:
\begin{align}
 \operatorname{Hit}_s(\pi)
 &=\ind\{L_{q_s^{(1)}}(\pi)\cap Y_s^{\mathrm{obs}}\ne\varnothing\},
 \label{eq:eval-session-hit}\\
 \operatorname{Hit@20}_J(\pi)
 &=\frac{1}{N_J^+}\sum_{s\in\mathcal S_J^+}\operatorname{Hit}_s(\pi),
 \qquad N_J^+>0.
 \label{eq:eval-observed-hit}
\end{align}
Это доля сессий с зарегистрированной будущей покупкой, в которых
первая выдача содержит хотя бы один целевой товар; показатель лежит
в $[0,1]$. Число покупок и повторных записей не увеличивает вклад сессии.
Изменение длины выдачи или числа запросов требует нового протокола.

Для всех методов состав $\mathcal S_J^+$ одинаков. При отказе,
запоздании или пустом результате действует предусмотренный резерв.
Если допустимый список не доставлен в срок, $L_q=\varnothing$;
сессия остаётся в знаменателе. Попадание резерва учитывается,
его частота публикуется отдельно. Цели вне каталога не удаляются.
При $N_J^+=0$ метрика недоступна, а не равна нулю.

Полнота телеметрии остаётся неизвестной даже для положительных сессий.
Сессии с $Y_s^{\mathrm{obs}}=\varnothing$ не считаются доказанно
непокупающими. Оценка условна на регистрации будущей покупки:
она не измеряет частоту покупок среди всех сессий, ненаблюдаемые покупки
или причинный эффект рекомендаций. Ни точечная оценка, ни её интервал
не устраняют смещение регистрации событий и исторической политики.

Контроли первой оценки --- популярность $A_G$ и последние просмотренные
товары с заполнением популярностью; порядок и фильтры заданы
разделом~\ref{sec:model-slate}. С ними публикуют парные разности Hit@20.
Чувствительность к метаданным проверяют по~\ref{sec:metadata-sensitivity}.
При выборе на $C_{\mathrm{select}}$ максимизируют
\eqref{eq:eval-observed-hit}; бюджет и правило равенств фиксируют до
$C_{\mathrm{select}}$, выбранную основную пару --- до $D$.
Интервалы и условия достаточности
заданы в разделе~\ref{sec:eval-uncertainty}, состав отчёта ---
в~\ref{app:diagnostic-output}. Подбор на $D$ запрещён.

\subsection{Точная верхняя граница отбора кандидатов}
\label{sec:eval-candidate-bound}

Ниже для периода $J$ все суммы по сессиям берутся по $\mathcal S_J^+$,
$N=N_J^+$, $q=q_s^{(1)}$. Определим
\begin{align}
 U_s(\pi)&=\ind\{\mathcal C_q^{\mathrm{served}}(\pi)\cap Y_q\ne\varnothing\},
 \label{eq:eval-candidate-hit}\\
 \operatorname{Coverage}_{\mathrm{cand},J}(\pi)
 &=\frac1N\sum_s U_s(\pi).
 \label{eq:eval-candidate-por}
\end{align}
Множество $\mathcal C_q^{\mathrm{served}}(\pi)$ определено
в~\eqref{eq:model-served-candidates}. При переключении ветви учитывают
резерв, при отказе уже выбранной выдачи сохраняют её кандидатов.
Охват завершённого $\mathcal C_q^{\mathrm{main}}$ приводят отдельно
с долей неизмеренных наборов; он не подменяет границу рабочей политики.
Совокупность сессий не ограничивается удачными ответами.

\begin{proposition}[Предел ранжирования при данном отборе]
\label{prop:eval-bound}
Если $L_q(\pi)\subseteq\mathcal C_q^{\mathrm{served}}(\pi)$, то
\begin{equation}\label{eq:eval-candidate-bound}
 \operatorname{Hit@20}_J(\pi)
 \le\operatorname{Coverage}_{\mathrm{cand},J}(\pi).
\end{equation}
При фиксированных ветвях и кандидатах, $K\ge1$, свободном выборе
поднабора размера не более $K$ и отсутствии отказов доставки граница
достижима ранжировщиком, знающим будущие покупки. Фиксированный порядок
резерва или дополнительные ограничения списка могут исключить достижимость;
само неравенство сохраняется.
\end{proposition}

\begin{proof}
Из вложения списков следует $\operatorname{Hit}_s\le U_s$.
При $U_s=1$ идеальный ранжировщик включает целевой товар в выдачу;
при $U_s=0$ попадание невозможно.
\end{proof}

\subsection{Обязательное разложение потерь охвата}
\label{sec:eval-loss-diagnosis}

Наличие допустимой цели равно
\begin{equation}\label{eq:eval-available-hit}
V_s=\ind\{\mathcal I_q\cap Y_q\ne\varnothing\}.
\end{equation}
Для сессий $\mathcal S_J^+$ выполнено
$1\ge V_s\ge U_s(\pi)\ge\operatorname{Hit}_s(\pi)$, поэтому
\begin{equation}\label{eq:eval-loss-decomposition}
1-\operatorname{Hit}_s(\pi)
=(1-V_s)+(V_s-U_s(\pi))+(U_s(\pi)-\operatorname{Hit}_s(\pi)).
\end{equation}
Три непересекающихся случая: нет допустимых целей; отбор не нашёл
допустимые цели; отобранные цели не доставлены в выдаче.
Средние слагаемых дают точное разложение $1-\operatorname{Hit@20}_J$.
Уже на $C_{\mathrm{smoke}}$ публикуют $N_J^{\mathrm{obs}}$, $N_J^+$,
$\sum_sV_s$, $\sum_sU_s$, $\sum_s\operatorname{Hit}_s$ и три потери.

Условный кандидатный охват при наличии допустимой цели равен
\begin{equation}\label{eq:eval-candidate-conditional}
\operatorname{Coverage}_{\mathrm{cand,avail}}
=\frac{\sum_sU_s}{\sum_sV_s},\qquad \sum_sV_s>0.
\end{equation}
Нулевой знаменатель означает недоступность показателя.
Вместе с ним приводят $\sum_sV_s/N$, отделяя покрытие каталога
от качества отбора.

В таблице покупок различают отсутствие товара в $\mathcal I_G$
и недопустимость при запросе. Её строки не складывают как сессионные
потери: в одной сессии возможны покупки разных типов.
Отказы источников, пропуски признаков и превышения срока ответа
дают отдельные разрезы, не четвёртое слагаемое потерь.
Попадание $i\in L_q\cap Y_q$ всегда означает $\operatorname{Hit}_s=1$.

\subsection{Вклад и пересечение источников}
\label{sec:eval-source-contribution}

Первый отчёт использует уже рассчитанные результаты источников.
Через $\mathcal R_q^{(a)}$ обозначим доступную часть упорядоченного
списка допустимых товаров источника $a$ до квоты $m_a$;
глубину и усечение этого списка записывают. Если источник возвращает
только ответ после квоты, охват до неё помечают неизмеренным.
Будущие цели не участвуют в поиске; диагностические списки
не подменяют фактические $\mathcal C_q$.
Для непустого объединения определим
\begin{align}\label{eq:eval-source-overlap}
O_q^{ab}&=\frac{|\mathcal R_q^{(a)}\cap\mathcal R_q^{(b)}|}
 {|\mathcal R_q^{(a)}\cup\mathcal R_q^{(b)}|},\\
\mathcal R_q^{(a),\mathrm{uniq}}
&=\mathcal R_q^{(a)}\setminus\bigcup_{b\ne a}\mathcal R_q^{(b)}.
\nonumber
\end{align}
При пустом объединении пересечение не определено. Для каждого
источника публикуют долю пустых ответов, размеры списков до квоты
и после неё, охват будущих целей до квоты, после квоты и в итоговом
объединении. Уникальные товары и пересечения показывают отдельно.
Охват любого семейства измеренных списков $\mathcal D=(\mathcal D_q)_q$
считают в тех же сессионных единицах:
\begin{equation}\label{eq:eval-source-coverage}
\mathsf H(\mathcal D)
=\frac1N\sum_s\ind\{\mathcal D_{q_s^{(1)}}\cap Y_s^{\mathrm{obs}}\ne\varnothing\}.
\end{equation}
Неполное измерение не заменяют пустыми списками; полный охват
до квоты в таком случае не определён. Пересечения усредняют
по определённым значениям с равными весами сессий;
отдельно сообщают долю пропусков.
Сумма охватов источников не является охватом их объединения.
Дополнительный поиск на большую глубину необязателен; его стоимость
отделяют от задержки рабочей конфигурации.

Для вклада при уже зафиксированных кандидатах обозначим через
$\mathcal O_q(i)$ множество их источников, включая резервное заполнение.
Удалим только товары, полученные исключительно из источника $a$:
\begin{align}\label{eq:eval-source-removal}
\mathcal C_q^{-a}
&=\{i\in\mathcal C_q:\mathcal O_q(i)\setminus\{a\}\ne\varnothing\},\\
\Delta_a^{\mathrm{fixed}}
&=\frac1N\sum_s(U_s-U_s^{-a})\ge0.
\nonumber
\end{align}
Здесь $U_s^{-a}$ вычисляется по~\eqref{eq:eval-candidate-hit}
с $\mathcal C_q^{-a}$; освободившиеся места не заполняются.
Уникальный товар запроса не обязательно даёт уникальный вклад сессии:
другой источник мог найти другой целевой товар той же сессии.
Это потеря кандидатного охвата при удалении источника из готовых
результатов, а не оценка новой работающей системы.
Последняя проверяется отдельным заранее выбранным опытом:
перераспределяют квоты, выполняют резервное заполнение, заново обучают
ранжировщик на $B$ и сравнивают политики по приложению~\ref{app:ablation}.
Такая разность может иметь любой знак. Переобучение для каждого
источника не входит в обязательный первый запуск.

\subsection{Выбор следующего изменения}
\label{sec:eval-next-step}

Исследовательскую гипотезу формируют на $C_{\mathrm{smoke}}$ после
проверки реализации и таблицы потерь; список сравниваемых вариантов
фиксируют до открытия $C_{\mathrm{select}}$.
Большое $\sum_s(1-V_s)$ требует анализа каталога и допуска целей;
большое $\sum_s(V_s-U_s)$ --- источников и квот;
большое $\sum_s(U_s-\operatorname{Hit}_s)$ --- ранжирования
и своевременного возврата списка. Сначала проверяют технические причины.
Малый кандидатный охват сам по себе не доказывает слабость отбора:
учитываются доступность целей и~\eqref{eq:eval-candidate-conditional}.

Кандидатная граница использует знание будущего. Разрыв с выдачей
не гарантирует достижимого улучшения ранжировщика и не оценивает
потерю денег или конверсии. Для сравнения только ранжирования фиксируют
одни $\mathcal C_q$: сокращение разрыва из-за ухудшения отбора
не считается улучшением. Новую гипотезу, бюджет и критерий успеха
фиксируют до её проверки; после просмотра $D$ нужен новый итоговый период.

\subsection{Вспомогательные показатели и разрезы}
\label{sec:eval-diagnostics}

Для запроса с $|Y_q|>0$ диагностический охват равен
\begin{equation}\label{eq:eval-request-recall}
 R_q@K=\frac{|L_q\cap Y_q|}{|Y_q|}.
\end{equation}
Знаменатель среднего по запросам с покупками публикуют отдельно.
Для диагностики порядка используют
\begin{align}
 \operatorname{DCG}_q@K
 &=\sum_{j=1}^{|L_q|}
       \frac{\ind\{L_{q,j}\in Y_q\}}{\log_2(j+1)},
 \label{eq:eval-dcg}\\
 \operatorname{NDCG}_q@K
 &=\frac{\operatorname{DCG}_q@K}
 {\displaystyle\sum_{j=1}^{\min(K,|Y_q|)}1/\log_2(j+1)},
 \quad |Y_q|>0.
 \label{eq:eval-ndcg}
\end{align}
Веса DCG не являются измеренными вероятностями просмотра позиций.
Средние Recall и NDCG считают на той же $\mathcal S_J^+$.

Отдельно показывают охват покупок товаров из наблюдаемой корзины
и вне неё. Состояние корзины определяется только по префиксу;
оно может быть неполным. Охваты групп не складывают: одна сессия
может содержать покупки обеих групп.

Обязательные разрезы задаются до просмотра итоговых результатов:
длина доступной истории, наличие предшествующего поиска,
наличие корзины, историческая популярность целевого товара,
наличие содержательных векторов и известность товара снимку.
Границы групп популярности вычисляют по $A_G$, а не по частотам периода $D$.

Покупки неизвестных или недопустимых товаров сохраняют в целях
и знаменателях, отдельно сообщая их долю. Первое наблюдение в журнале
не доказывает новизну товара: дата появления в каталоге неизвестна.

\section{Неопределённость исторического сравнения}
\label{sec:eval-uncertainty}

Оцениваются первые выдачи на $\mathcal S_D^+$ --- сессиях с зарегистрированной
будущей покупкой по разделу~\ref{sec:eval-observed}. Далее $N=N_D^+=|\mathcal S_D^+|$,
а суммы по $s$ берутся только по этой совокупности. Интервалы условны на попадание
покупки в журнал; они не устраняют смещение из-за неполной телеметрии.

\subsection{Парное сравнение сессий}\label{sec:eval-paired}

Политики $\pi_1$ и $\pi_0$ оцениваются на одних и тех же сессиях.
При $N>0$ разность долей первых выдач с попаданием равна
\begin{equation}\label{eq:eval-paired-diff}
 d_s=\operatorname{Hit}_s(\pi_1)-\operatorname{Hit}_s(\pi_0)\in\{-1,0,1\},\qquad
 \widehat\Delta_{\mathrm{Hit}}=
          \frac{1}{N}\sum_s d_s.
\end{equation}
При повторной выборке результаты обеих политик переносят вместе,
сохраняя информацию об их расхождениях на каждой сессии.

Для одной фиксированной политики оценивают сам охват:
\begin{equation}\label{eq:eval-single-policy}
x_s=\operatorname{Hit}_s(\pi),\qquad
\widehat H=\operatorname{Hit@20}_D(\pi)=\frac{1}{N}\sum_s x_s.
\end{equation}
Его интервал строят тем же календарным методом, заменяя $d_s$ на $x_s$
в~\eqref{eq:eval-block-stat}--\eqref{eq:eval-block-bootstrap}.
Порог полуширины $h_{\mathrm{single}}>0$ задают отдельно от
$h_{\mathrm{pair}}$ парного сравнения; оба порога выражены в долях.
При $N=0$ оценки недоступны, результат --- \code{insufficient}.

При $N>1$ и приближённой независимости сессий вторичная,
условная на это предположение стандартная ошибка равна
\begin{equation}\label{eq:eval-paired-se}
 \widehat{\operatorname{se}}(\widehat\Delta_{\mathrm{Hit}})
 =\sqrt{\frac{1}{N(N-1)}
       \sum_s(d_s-\bar d)^2},\qquad
 \bar d=N^{-1}\sum_s d_s.
\end{equation}
Основной интервал Coveo строят временным бутстрэпом последовательных
блоков по разделу~\ref{sec:eval-time-block}, при условиях
раздела~\ref{sec:eval-information}. Независимая повторная выборка сессий
\cite{bootstrap} служит лишь проверкой чувствительности и не заменяет
непригодный основной интервал.

Для повторения $b$ выбирают $N$ индексов сессий независимо и равновероятно
из $\{1,\ldots,N\}$ и вычисляют
\begin{equation}\label{eq:eval-bootstrap}
 \widehat\Delta_{\mathrm{Hit}}^{*(b)}
 =\frac{1}{N}\sum_{j=1}^{N}d_{I_j^{(b)}},
 \qquad b=1,\ldots,B_{\mathrm{boot}}.
\end{equation}
Парные исходы первой выдачи переносят вместе. Число повторений
$B_{\mathrm{boot}}=2000$ и начальное состояние генератора фиксируют
заранее. Квантили определены в разделе~\ref{sec:eval-time-block}.

\subsection{Планирование точности периода \texorpdfstring{$D$}{D}}
\label{sec:eval-precision}

До открытия $D$ задают полуширину $h_{\mathrm{pair}}>0$ интервала парной
разности в долях и плановую дисперсию
$v_{\mathrm{plan}}\approx\Var(d_s)$. Её обосновывают по разрешённому
пилоту на $C_{\mathrm{smoke}}$ после обучения, без $D$ и утечек будущего.
Порог, дисперсия, допущения, множитель временной зависимости и происхождение
пилота входят в неизменяемый план оценки до $D$; их отсутствие не блокирует
\code{preflight} и первое обучение. При независимых сессиях
\begin{equation}\label{eq:eval-precision-size}
 N_D^{+,\mathrm{iid}}=
 \left\lceil
 \frac{z_{1-\alpha/2}^{\,2}\,v_{\mathrm{plan}}}
      {h_{\mathrm{pair}}^2}
 \right\rceil,
\end{equation}
где $z_u$ --- квантиль уровня $u$ стандартного нормального распределения.
Это следствие приближения полуширины
$z_{1-\alpha/2}\sqrt{v_{\mathrm{plan}}/N_D^+}$; мощность эксперимента
по причинному эффекту рассчитывается отдельно.

Для оценки одной политики используют ту же схему с заменами
$d_s\mapsto\operatorname{Hit}_s(\pi)$,
$h_{\mathrm{pair}}\mapsto h_{\mathrm{single}}$ и
$v_{\mathrm{plan}}\approx\Var(\operatorname{Hit}_s(\pi))$.
Тип основной оценки фиксируют в \code{inference.primary_estimate_type}.
Пара политик и порог парной разности обязательны только для
\code{paired}; для \code{single_policy} обязательны одна политика,
$h_{\mathrm{single}}$ и плановая дисперсия её исхода.

Дисперсия зависит от расхождений политик, а не только числа покупок.
Нулевая дисперсия малого пилота не обосновывает $v_{\mathrm{plan}}=0$.
При недостатке данных задают запас либо грубую границу $\Var(d_s)\le1$;
происхождение и допущения расчёта включают в протокол.

Для временной зависимости задают множитель
$\kappa_{\mathrm{dep}}\ge1$ и планируют
\begin{equation}\label{eq:eval-precision-dependent}
 N_D^{+,\mathrm{plan}}=
 \left\lceil\kappa_{\mathrm{dep}}N_D^{+,\mathrm{iid}}\right\rceil.
\end{equation}
На пилоте делят дисперсию при повторной выборке временных блоков
на дисперсию при независимых сессиях, ограничивая результат снизу единицей.
При нулевом знаменателе отношение не определено. Длительность блоков
и запас фиксируют заранее. Множитель не гарантирует точности при смене
режима или малом числе блоков; он не корректирует итоговый интервал
для независимых сессий.
Для основного интервала сохраняются условия зависимости
из раздела~\ref{sec:eval-time-block}.

Границы периодов учитывают прогноз числа сессий с зарегистрированной целью,
календарные циклы и исключение последних $H+g$ из окна стартов.
Правило отбора сессий не зависит от наличия цели. После открытия $D$
границы и правило набора не меняют ради удобного результата.
Публикуют фактический интервал, полуширину и выполнение требования
$h_{\mathrm{pair}}$ либо $h_{\mathrm{single}}$.
Дополнительный набор требует нового протокола.
Для фиксированного открытого журнала допустим результат
\texttt{insufficient}: публикуют доступную оценку и причину недостаточной
точности, не сдвигая периоды и не требуя недоступных дополнительных данных.
Выбранный порог $h_{\mathrm{pair}}$ либо $h_{\mathrm{single}}$,
параметры $v_{\mathrm{plan}},\kappa_{\mathrm{dep}}$,
их происхождение и допущения становятся обязательными полями паспорта
на \code{pre_D}; фактические границы периодов закреплены до обучения.
Незаполненное поле не подменяют вымышленным измерением.

\subsection{Временная зависимость и границы интерпретации}
\label{sec:eval-time-block}

Общие акции, ассортимент и спрос создают временную зависимость сессий.
До открытия $D$ фиксируют плотную календарную сетку с шагом
$\delta_{\mathrm{grid}}>0$ и началом, привязанным к границе $D$.
Для $D=[a_D,b_D)$ её $J$ ячеек покрывают только допустимое окно
стартов $[a_D,b_D-H-g)$, без завершающего календарного запаса $H+g$.
Этот запас не подтверждает полноту наблюдений.
Требуют целое число ячеек: границы и шаг согласуют до открытия $D$;
при несогласованных границах паспорт недействителен.
Сессию целиком относят к ячейке по $t_{0,s}$. Все ячейки внутри окна
сохраняют, в том числе без сессий с зарегистрированной целью;
искусственный хвост не добавляют. Пусть $\mathcal S_j^+$ содержит сессии
из $\mathcal S_D^+$ в ячейке $j$. Записывают
\begin{equation}\label{eq:eval-block-stat}
 n_j=|\mathcal S_j^+|,\qquad D_j=\sum_{s\in\mathcal S_j^+}d_s.
\end{equation}
Бутстрэп последовательных блоков (moving-block bootstrap)
для стационарных зависимых наблюдений
\cite{kunsch1989} выбирает длину
$\ell_{\mathrm{boot}}=m\delta_{\mathrm{grid}}$, $m\in\{1,\ldots,J\}$,
по разрешённому пилоту на $C_{\mathrm{smoke}}$ и предметным ограничениям
без использования $D$; длину блока фиксируют до открытия $D$.
Для каждого повторения независимо выбирают
$\lceil J/m\rceil$ начал равновероятно из $\{1,\ldots,J-m+1\}$;
после каждого начала берут $m$ последовательных ячеек, склеивают
полученные цепочки и оставляют первые $J$ индексов
$J_1^*,\ldots,J_J^*$. Таким образом, соседство сохраняется внутри
выбранного блока, а две политики всегда переносятся вместе.
Используют отношение сумм
\begin{equation}\label{eq:eval-block-bootstrap}
 \widehat\Delta_{\mathrm{Hit}}^*
 =\frac{\sum_{\ell=1}^{J}D_{J_\ell^*}}
              {\sum_{\ell=1}^{J}n_{J_\ell^*}}.
\end{equation}
Числитель и знаменатель пересчитывают в каждом повторении.
Невзвешенное среднее дневных метрик при разном числе сессий имеет иной смысл.
Основной интервал --- эмпирические квантили $\alpha/2$, $1-\alpha/2$
из $B_{\mathrm{boot}}$ повторений~\eqref{eq:eval-block-bootstrap}.
Для всех бутстрэп-интервалов этой спецификации квантиль уровня $p$
равен элементу с номером
$\lceil pB_{\mathrm{boot}}\rceil$ отсортированной выборки.
Нулевой знаменатель делает повторение неопределённым: такие случаи
регистрируют и присваивают расчёту интервала \texttt{insufficient},
а не отбрасывают незаметно и не заменяют нулём.

\begin{assumption}[Допустимость временного бутстрэпа]
\label{ass:eval-block-dependence}
Для каждого оцениваемого функционала задаётся свой вектор блочных сумм.
Для парной разности это $(D_j,n_j)$, для одного охвата ---
$(\sum_{s\in\mathcal S_j^+}\operatorname{Hit}_s(\pi),n_j)$.
Соответствующий совместный процесс достаточно близок к стационарному
на окне допустимых стартов, имеет конечные вторые моменты,
положительные средние знаменатели и ослабевающую с лагом зависимость.
Длина блока покрывает существенный горизонт зависимости; отдалённые
блоки не связаны доминирующим общим шоком. Асимптотическое обоснование
требует роста длины блока при её малости относительно длины наблюдения.
Требуются совместная центральная предельная теорема для средних
компонент и состоятельность выбранного бутстрэпа для того же предела;
качественная стационарность сама по себе этих свойств не доказывает.
Оцениваемый функционал должен быть регулярным в предельных средних
и иметь ненулевую предельную вариацию.
\end{assumption}
Устойчивого межсессионного ID в Coveo нет, поэтому календарные блоки
не доказывают устранение зависимости повторных пользователей.
Если в другом журнале такой ID доступен и пользовательская зависимость
доминирует, до открытия теста выбирают кластерный анализ по пользователям;
общие календарные шоки требуют дополнительного учёта временной зависимости.

В паспорте до $D$ фиксируют сетку, длины блоков, две дополнительные
длины для чувствительности, генератор и число повторений.
Необходимый технический порог ---
$J_{\mathrm{eff}}=\lfloor J/m\rfloor\ge20$ для основного расчёта;
это число доступных полных блоков, не оценка числа независимых наблюдений.
Для фиксированного разбиения $\{1,\ldots,m\},\{m+1,\ldots,2m\},\ldots$
из $J_{\mathrm{eff}}$ полных блоков число блоков с $\sum n_j>0$
должно быть не меньше 20; перекрывающиеся бутстрэп-блоки здесь
не пересчитывают как новые наблюдения. Пустые промежутки не создают информацию.
Для каждого функционала и разреза учитывают только блоки с его
наблюдениями. Невыполнение порога или предположения
\ref{ass:eval-block-dependence} даёт \code{insufficient} для соответствующего
интервала. Пригодность процесса парной разности не доказывает
пригодности процесса охвата одной политики и наоборот.

Публикуют временной график $D_j/n_j$ при $n_j>0$, число сессий $n_j$,
границы акций и изменений ассортимента, основной и два заранее заданных
интервала чувствительности. Пустые ячейки отмечают как отсутствующую
метрику. Тренд или существенная зависимость вывода от длины блока
ограничивает интерпретацию; выбирать наиболее удобный интервал запрещено.
Данные $D$ не используют для настройки длины, удаления проблемных дней
или замены основного метода. Недостаточное число блоков не исправляют
простым увеличением числа бутстрэп-повторений.
Для границ интервала указывают погрешность случайного моделирования
по заранее выбранному способу; число повторений не определяет
статистическую достаточность исходного потока.

Оба способа оценивают уже обученные модели на данном потоке.
Неопределённость повторного обучения и смены сезона требует нескольких
временных проверок с повторением обучения; перекрывающиеся проверки
не являются независимыми.

\subsection{Достаточность информации и вырожденные интервалы}
\label{sec:eval-information}

Число непустых блоков не защищает от редких расхождений политик.
Для проверяемого охвата положим $c_s=d_s$ при сравнении пары и
$c_s=\operatorname{Hit}_s(\pi)$ при оценке одной политики.
Пусть $\mathcal B_b$ содержит сессии из $\mathcal S_D^+$, относящиеся к полному
непересекающемуся блоку из $m$ ячеек по разделу~\ref{sec:eval-time-block}.
Для каждой категории $a$ ($-1,0,1$ для пары; $0,1$ для одной политики)
считают
\begin{equation}\label{eq:eval-information-gate}
N_a=\sum_s\ind\{c_s=a\},\qquad
J_a=\sum_{b=1}^{J_{\mathrm{eff}}}
 \ind\{\exists s\in\mathcal B_b:c_s=a\},\qquad
\left|\{a:N_a\ge100,\ J_a\ge20\}\right|\ge2.
\end{equation}
Это начальные защитные пороги: для одной политики нужны как минимум
100 попаданий и 100 промахов, каждый исход --- не менее чем в 20 блоках.
Для пары требуются две достаточно представленные категории; следовательно,
не менее 100 расхождений и 20 блоков с ними. Редкие категории
не объединяют после просмотра $D$, чтобы пройти порог.

Нарушение~\eqref{eq:eval-information-gate}, нулевая эмпирическая вариация,
неопределённый повтор или совпадение границ основного бутстрэп-интервала
дают \code{insufficient}. Точечную оценку и все счётчики публикуют,
но нулевая ширина не подтверждает точность и не разрешает
$S_{\mathrm{offline}}=\mathrm{passed}$. Доказанное для всех допустимых входов
тождество политик сообщают отдельно как алгебраический результат
с доказательством и хешами артефактов, а не как измеренную точность.

Для пригодного интервала $[l,u]$ дополнительно проверяют
$(u-l)/2\le h_{\mathrm{pair}}$ либо $h_{\mathrm{single}}$.
Пороги числа наблюдений, блоков и полуширины фиксируют до $D$.
Они отсекают явно неинформативные оценки, но не гарантируют покрытия,
мощности или отсутствия долговременной зависимости. Интервал для редких
исходов с иными предпосылками требует отдельного предварительного плана;
его не выбирают вместо основного после просмотра результата.


\section{Хранение и обновление признаков}\label{sec:operations-features}

\subsection{Согласование обучения и обслуживания}

Признак \(j\) определяется контрактом, общим для обучения и обслуживания:
\begin{equation}\label{eq:ops-feature-contract}
\mathsf F_j=(\mathrm{entity},\phi_j,\mathrm{sources},
\mathrm{available\_at},\mathrm{version},\mathrm{missing\_rule}).
\end{equation}
Версия признака принадлежит закреплённому пакету. Пакет выбирают среди
совместимых версий с временем публикации строго меньше $t_q$, даже если
$a_q>t_q$; внутри этого множества берут последнюю по порядку публикации
из паспорта. Версия, опубликованная между $t_q$ и $a_q$, запрещена.
При отсутствии подходящего пакета фиксируют \code{version_unavailable};
новый пакет не используется даже для резерва.

Для воспроизведённого запроса проверяется
\begin{equation}\label{eq:ops-feature-parity}
\max_j\frac{\abs{x^{\mathrm{replay}}_{qj}-x^{\mathrm{serve}}_{qj}}}
{\eps_j+\abs{x^{\mathrm{replay}}_{qj}}}\le\delta_{\mathrm{feat}}.
\end{equation}
Категории, маски, кандидаты, снимки и порядок при равных баллах
совпадают точно; относительный допуск относится к числам с плавающей
точкой. Параметры, тип чисел и индекс одинаковы, обучающие преобразования,
случайное удаление компонент и шум кодировщика отключены.
Для FP16 и BF16 допуски устанавливаются отдельно от двойной точности.
Проверка включает исходные упорядоченные списки каждого источника,
маски доступности, квоты, устранение повторов и выбранную ветвь;
совпадения только итогового вектора признаков недостаточно.

\subsection{Численная устойчивость}\label{sec:ops-numerics}

Эквивалентные математические выражения реализуются устойчивыми
примитивами. Для конечного $x$ используются
\begin{align}
\operatorname{softplus}(x)&=\max(x,0)+\log1p(e^{-|x|}),\\
\sigmoid(x)&=\begin{cases}
1/(1+e^{-x}),&x\ge0,\\
e^x/(1+e^x),&x<0.
\end{cases}
\end{align}
Парная логистическая потеря вычисляется как
$\operatorname{softplus}(-\beta_R d)$; бинарная кросс-энтропия ---
через логиты, без вычисления $\log(\sigmoid(x))$ и
$\log(1-\sigmoid(x))$. Суммы экспонент используют
\begin{equation}\label{eq:ops-stable-lse}
\operatorname{LSE}(z)=z_*+\log\sum_j e^{z_j-z_*},
\qquad z_*=\max_j z_j,
\end{equation}
либо \code{logaddexp} для блочного накопления. Правило действует
для внимания, softmax, весов намерений, контрастивной потери, SID
и нормировок каталога. Маски применяются до максимума. Пустая допустимая область
обрабатывается отдельной ветвью: источник отсутствует, а не вычисляется
$-\infty-(-\infty)$. Значение $-\infty$ допустимо только как явно
заданный маркер исключённого логита; оно не является баллом товара.
Энтропия и KL используют сохранённые логарифмы вероятностей
по~\ref{sec:modern-log-probabilities}: логарифм округлённого нуля
не вычисляют; проверяются также производные.
Точные биномиальные границы вычисляются через квантили beta-распределения
с аналитическими граничными случаями для нуля и всех успехов;
хвосты --- через log-CDF, без суммирования факториалов.

Первый запуск использует FP64 для признаков, баллов,
метрик и статистических интервалов. Исключение: вычисленные в FP64
производные заданной функции передаются в LightGBM через непрерывные
FP32-массивы по контракту~\ref{sec:model-ranker-implementation};
их допуски проверяются отдельно. Идентификаторы и счётчики не
хранятся в FP16/BF16; счётчики используют целые 64-битные значения
с проверкой переполнения. В нейросетевых вариантах FP16/BF16
допустимы для матричных операций только с накоплением не ниже FP32;
нормы, softmax/LSE, потери и производные накапливаются не ниже FP32,
а итоговые метрики и статистика --- в FP64. Паспорт фиксирует
типы чисел каждого этапа, библиотеки, режимы округления и допуски.
Начальные допуски проверки~\eqref{eq:ops-feature-parity} для FP64:
$\eps_j=10^{-8}$ и $\delta_{\mathrm{feat}}=10^{-6}$.
Другой тип чисел требует отдельной проверки против FP64 до выбора модели.

Норму конечного вектора вычисляют с масштабированием: при
$v_*=\max_j|v_j|>0$ сначала находят $w=v/v_*$ и $r=\sqrt{\sum_jw_j^2}$,
а единичное направление --- $w/r$. Проверку $\norm v>\eps$
выполняют как $v_*>\eps/r$, без потенциально переполняющегося
произведения $v_*r$; при $v_*=0$ сразу выбирают нулевую ветвь.
Конечность отдельных координат не оправдывает прямую сумму их квадратов.

Внешняя модальность и вычислительная ошибка различаются.
Отсутствующий внешний вектор получает маску 0 и причину
\code{vector_absent}; неверная размерность, нечисловые/неконечные
координаты --- маску 0 и \code{vector_invalid}; норма не больше
$\eps$ --- маску 0 и \code{vector_zero}. Во всех трёх случаях
используется нулевой вектор ожидаемой размерности. Это одинаковое
преобразование подготовки, обучения и обслуживания; остальные пригодные
модальности сохраняются. Ошибка обязательного ID, времени или структуры
конверта даёт \code{invalid_input}, а не отсутствующую модальность.
NaN/Inf, неожиданно возникшие после такой подготовки в признаках,
потере, производных или баллах, останавливают обучение; в обслуживании
они делают расчёт недействительным и включают объявленный резерв.
Подмена внутренней ошибки нулевой маской запрещена.
Тесты охватывают логиты $\pm100,\pm10^4$, пустые маски,
нулевые нормы и максимально допустимые группы. Эталон в FP64,
конечность разрешённых результатов и неизменность масок обязательны;
для FP64 абсолютный/относительный допуски равны $10^{-10}/10^{-8}$,
для FP32 --- $10^{-6}/10^{-5}$. Изменение порядка почти равных
баллов не скрывается числовым допуском: действует общий порядок разрешения равенств.

\subsection{Обновление состояния сессии}

Состояние \(\mathsf U_s(t)\) содержит ограниченную последнюю историю,
счётчики действий, время последнего события, наблюдаемое состояние
корзины и поисковый контекст. Для нового доступного блока \(E_{s,t}\)
\begin{equation}\label{eq:ops-state-update}
\mathsf U_s(t)=\mathcal T\bigl(\mathsf U_s(t^-),E_{s,t};\nu_{\mathrm{features}}\bigr).
\end{equation}
\(\mathcal T\) детерминировано и зависит только от доступной истории.
Повторная доставка обрабатывается по контракту
раздела~\ref{sec:ops-state-storage}; совпадение строк Coveo
не заменяет транспортный идентификатор события.

Для затухающего агрегата одного типа действия удобно использовать
рекуррентную формулу
\begin{equation}\label{eq:ops-decay-update}
x_{si}(t)=e^{-(t-t')/\tau_x}x_{si}(t')+
\sum_{e\in E_{s,t}}\ind\{i(e)=i,\ \kappa(e)=\kappa_x\}.
\end{equation}
При последовательной обработке это сумма затухающих прошлых событий.
Точное скользящее окно требует очереди событий.

Хранятся последние \(L_{\mathrm{hist}}\) общих блоков для воспроизведения,
отдельная очередь последних $L_H$ событий \code{detail} для графа,
последний допустимый поисковый контекст и необходимые агрегаты
всего видимого префикса. Общий хвост не заменяет очередь просмотров:
частые поиски не должны вытеснять ещё нужные $L_H$ событий.
Память всех очередей, агрегатов и отметок доставки учитывается отдельно.
Удаление блока не обнуляет
агрегаты всей видимой сессии. Оперативное состояние удаляется
по неактивности после необходимых записей об исходе; сроки хранения
состояния и исходов задаются раздельно.

\subsection{Контракт хранилища состояния}
\label{sec:ops-state-storage}

Для воспроизведения достаточно памяти процесса; обслуживание требует
хранилища по ключу с атомарным обновлением записи. Продукт не предписан;
сроки хранения и пределы доставки фиксируются в описании испытания.
\begin{enumerate}
\item \textbf{Ключ и состав записи.}
Ключ $(\nu_{\mathrm{source}},s)$ включает пространство идентификаторов
источника. Запись содержит $\mathsf U_s$, номер изменения $v_s$,
версию $\nu_{\mathrm{features}}$, время последнего поступления,
признак полноты и отметки принятых событий.
Покупки и журнал запросов хранятся независимо от оперативной записи.

\item \textbf{Одновременные обращения.}
Проверка неизменности прочитанного $v_s$ и запись нового состояния,
$v_s+1$ и отметок событий выполняются одним атомарным действием.
При несовпадении номера запись перечитывается и переход пересчитывается
в пределах общего срока ответа.
Запрос использует одну подтверждённую версию состояния и один пакет
$\mathsf V$ без смешения полей. Срез включает только доступные запросу
события (раздел~\ref{sec:time}); при наличии в записи более поздних
событий он восстанавливается из журнала.
Версия преобразований совпадает с пакетом; её смена требует
объявленного преобразования состояния или восстановления.

\item \textbf{Повторная доставка.}
Квитанция блока, проверка ID и изменение счётчиков/полноты атомарны.
Квитанции и отметки событий хранятся не менее $\Delta_{\mathrm{retry}}$;
подтверждённый повтор не меняет логическое состояние, включая полноту.
После истечения этого срока ID проверяют по журналу либо отвечают
\code{retry_unverifiable}, не считая доставку новым событием.
Порядок транспортных проверок задан в~\ref{sec:ops-transport-block}.

\item \textbf{Срок хранения и запоздание.}
Раздельно задаются время жизни после последнего поступления
$T_{\mathrm{state}}$, предел запоздания $\delta_{\mathrm{late}}$
и срок хранения исходов. Хвост ограничен числом блоков и временем.
После удаления записи позднее событие требует восстановления
по журналу либо отметки неполноты.

\item \textbf{Отказ и восстановление.}
Недоступная, несовместимая или неполная запись включает заданный резерв
с указанием причины; непрочитанная история не считается пустой.
Неготовность допустимого резерва или обязательной записи к сроку
означает технический отказ.
Восстановление среза, проверка версии и повтор атомарной записи
входят в бюджет состояния, а не выполняются после его истечения.
Начальные ограничения: один повтор атомарного обновления,
не более 1000 событий восстановления и не более 1 МиБ журнала
на запрос. Превышение любого ограничения даёт
\code{state_replay_limit}; фоновое восстановление не задерживает ответ.

\item \textbf{Сохранённые результаты вычислений.}
Ключ повторного использования зависит от видимого префикса,
времени запроса и пакета преобразований:
\[
(\nu_{\mathcal H_q},t_q,\mathsf V).
\]
Идентификатор $\nu_{\mathcal H_q}$ определяет сессию, видимые события
и поиск. Веса истории и намерения зависят также от времени и модели;
одного идентификатора сессии для повторного использования недостаточно.
\end{enumerate}

\subsection{Что требуется для воспроизведения запроса}
\label{sec:ops-replay-trace}

Паспорт задаёт один из двух режимов. В \code{ideal_deterministic}
используются $a=t$, замороженные данные и штатные источники; этот режим
не воспроизводит наблюдавшиеся эксплуатационные отказы.
Для него оценки на $C_{\mathrm{smoke}},C_{\mathrm{select}},E,D$
не требуют реализации сервиса или этапа \code{post_load}.
Режим и его предположения указывают в отчёте; локальные время и память
не подтверждают эксплуатационные гарантии.
В \code{execution_trace} дополнительно сохраняется трасса оценки:
$t_q,a_q$, пакет, ссылка на неизменяемый видимый префикс/снимок состояния,
выбранная ветвь и статус доставки, маски и причины недоступности
источников, решения об усечении, seed и расписание внесённых отказов.
Ссылка включает хеш, но сам хеш не заменяет сохраняемые входы.

Трасса фиксирует фактическое число шагов/узлов при досрочной остановке
поиска; если оно не воспроизводит список однозначно, сохраняется список
этого источника. Полнота трассы и наличие входов проверяются до расчёта
$\mathcal C_q^{\mathrm{served}}$ и $U_s$. Здоровый повторный запуск
не подставляется вместо утраченной трассы: такой набор кандидатов
помечается неизмеренным. Доступ и хранение трасс подчиняются
разделу~\ref{sec:ops-security}; истечение срока входных данных прекращает
возможность точного воспроизведения, но не отменяет сохранённые агрегированные оценки.

\subsection{Время поступления и запоздавшие события}
\label{sec:ops-event-time-validation}

Входное время --- целое число миллисекунд UTC от начала эпохи Unix
в знаковом 64-битном диапазоне. До рабочего запуска в паспорте задают
допустимый календарный диапазон $[T_{\min}^{\mathrm{in}},T_{\max}^{\mathrm{in}})$;
он должен покрывать сроки приёма и разрешённого запаздывания.
Неверные тип, единица или диапазон дают \code{invalid_input};
значение не округляют и не исправляют по предположению об единицах.
Проверенный повтор принятого блока обрабатывают до повторной проверки
временных границ и не меняют его исходные отметки.
Для нового блока до изменения истории проверяют
\begin{equation}\label{eq:ops-event-time-window}
-\delta_{\mathrm{late}}\le t(e)-a_{\mathrm{block}}
\le\delta_{\mathrm{future}}.
\end{equation}
$a_{\mathrm{block}}$ --- серверное время атомарного решения.
В начальном профиле $\delta_{\mathrm{future}}=0$: буфера ожидания нет.
Нарушение правой границы даёт карантин с причиной \code{timestamp_future},
левой --- \code{timestamp_late}. Блок не меняет историю,
$t_{\mathrm{last},s}$ и закрытые временные группы. Для нового доверенного
блока отмечают неполноту по~\ref{sec:ops-transport-block};
до восстановления действует резерв. Ошибка недоверенного отправителя
не меняет чужую сессию. Положительный допуск опережения в этом профиле запрещён.

Принятое позднее событие меняет только ответы после поступления.
Состояние восстанавливают из контрольной точки перед хвостом
и событий хвоста в порядке времён. Точка содержит все агрегаты и очереди;
без неё нужен полный журнал в пределах бюджета, иначе --- неполнота и резерв.
Независимый журнал покупок продолжает учёт исходов; прошлые ответы не меняются.
Задержки, повторы и нарушение порядка проверяются отдельно
от основного исторического режима \(a=t\).

\subsection{Частичная доступность данных}\label{sec:ops-partial-data}

Отсутствующий признак, отключённый источник и отказ чтения имеют
разные коды. Отсутствующая у товара модальность использует её штатную
маску, если такой режим был включён в обучение. Сбой хранилища
не кодируется как отсутствие модальности: источник отключается только
по заранее зарегистрированному правилу деградации; незарегистрированное
сочетание недоступных источников включает резерв. Частичный список
источника допустим лишь с явным признаком полноты и предписанной квотой;
произвольная часть списка на момент истечения срока не считается полным результатом.

Каталог допустимости является обязательным: товар выдаётся только
после подтверждения его допустимости в версии запроса. Старая копия
разрешена лишь при действующем сроке актуальности, закреплённом в
паспорте. Неизвестная допустимость исключает товар из основной и
резервной ветвей; невозможность подтвердить хотя бы резерв даёт
технический отказ. Непрочитанная история не трактуется как новая сессия.
Для каждой маски деградации публикуются частота, длина ответа,
качество и число технических отказов.

\subsection{Границы доверия, ограничения входа и хранение трасс}
\label{sec:ops-security}

Модель угроз включает поддельные/массовые события, Sybil-сессии,
рост популярности SKU, дорогое восстановление истории и подмену артефактов.
Аутентификация источника подтверждает происхождение, а не достоверность
поведения. До вычислений проверяют обязательную схему и размеры;
непригодные внешние векторы обрабатывают по~\ref{sec:ops-numerics}.

\paragraph{Закрытый транспортный блок.}\label{sec:ops-transport-block}
Начальный профиль рабочего сервиса --- \code{sealed_block_v1}.
За каждое пространство сессий отвечает один зарегистрированный
агрегатор: он объединяет поведение и поиск до отправки блока.
Реестр владельцев фиксируется до запуска; полномочия проверяют по
аутентифицированному отправителю, а не присланному имени пространства.
Один неизменяемый конверт содержит пространство источника, ID сессии,
время событий, уникальный \code{block_id}, число событий, их устойчивые
ID, сами события, SHA-256 нагрузки и \code{closed=true}.
Все события имеют одни $s,t$; число событий равно длинам нагрузки
и списка ID, а \code{event_ids[k]=events[k].event_id} для каждого $k$.
ID внутри конверта уникальны. Хеш вычисляется сервисом по точным байтам неизменяемой
нагрузки. Квитанция связывает пространство, сессию, время, ID блока,
число событий и этот хеш; совпадения одного ID недостаточно.
Иная сериализация не является идентичным транспортным повтором.
Сервис принимает только целый конверт и не объединяет фрагменты
из разных запросов. Неполный конверт не добавляет событий;
его отклонение для установленной доверенной сессии учитывается
по правилам полноты ниже.
Получение полного конверта и ожидание очереди входят в общий срок от первого
байта по разделу~\ref{sec:ops-deadlines}. Отдельный предел 10 мс
для демонстрационного сервиса не применяется. Сборка у агрегатора отражается
в $a(e)-t(e)$ и отдельно измеряемой задержке от события до ответа.
Сервис фиксирует $a_{\mathrm{block}}$ при атомарном решении о проверенном
блоке, в том числе при карантине. Для принятых событий $a(e)=a_{\mathrm{block}}$;
отклонённые события в доступную историю не включаются.
Время отправителя не заменяет эти отметки.
При совмещённом приёме и запросе отсечка $a_q$
фиксируется после принятия блока; срок ответа всё равно отсчитывается
от первого байта. Восстановление повторов не меняет исходное $a(e)$.

Пределы: 1 МиБ на конверт, 256 байт UTF-8 на SKU,
от 1 до 100 сырых событий в конверте, 100 логических событий в блоке,
1000 логических событий в сессии, 1000 элементов поискового списка.
Логические счётчики определяются после удаления повторных event ID
и выбранной канонизации. Перед приёмом всего блока проверяется
сумма прежнего счётчика и его новых логических событий; пересечение
предела отклоняет весь новый блок, без частичного усечения.
Неиспользуемые списки поиска удаляются после проверки входа.

Уже закрытые $(\nu_{\mathrm{source}},s,t)$ не дополняются и не открываются
повторно. Тот же \code{block_id} с другим содержимым либо новые ID
для закрытого времени --- \code{closed_block_conflict}, а не допустимое
запаздывание. Новая неизвестная часть карантинируется, принятая история
сохраняется. Обычная поздняя доставка допускается только для ещё
не принятого блока и подчиняется~\eqref{eq:ops-event-time-window}.
Адаптер исторических данных группирует весь известный временной блок
до создания конверта; рабочий сервис не предполагает знания будущих пакетов.

Предел скорости --- token bucket 10 сырых событий/с с burst 100;
стоимость полной попытки, включая повторы, проверяется атомарно.
Стоимость определяется фактическим числом сырых событий, не заголовком;
для пустой, отклонённой или неразбираемой попытки она не меньше 1.
До установления сессии действуют ограничения соединения и отправителя.
Далее применяются раздельные переходы:
\begin{enumerate}
\item Подтверждённый полный повтор принятого блока не меняет счётчики,
полноту или ответы. При недостатке токенов его можно отклонить как
\code{duplicate_rate_rejected}; это не разрыв истории. Правило действует
и после $\delta_{\mathrm{late}}$, пока квитанция проверяема.
\item Конверт, содержащий только известные события, не создаёт нового
блока. Частичный повтор с новыми ID не обходится поэлементно:
для уже закрытого времени это конфликт, для нового блока --- нарушение
соответствия event ID его исходным полям. Новые неподтверждённые события
не считаются принятыми; старые записи не меняются.
\item Отклонённый новый блок от доверенного агрегатора оставляет
отметку разрыва и неполноты; дальнейшие запросы используют резерв
до восстановления. Непроверяемый повтор получает
\code{retry_unverifiable}, без повторного начисления событий;
восстановление или проверка журнала определяют полноту.
\end{enumerate}
Проверка квитанции ограничена одним индексным чтением и размером
конверта. При недоступном хранилище действует резервный контракт
состояния; ошибка аутентификации не меняет чужую сессию.
Карантин новых событий воспроизводится offline с отметкой разрыва:
переходы графа через неизвестный участок запрещены. Входные ограничения
не используют будущую длину сессии, не меняют назначение, не удаляют
допущенную сессию из оценки/ITT и не отключают независимый журнал покупок.
Они не заменяют $Q_{\max}$ числа рекомендационных запросов.

Лимиты арендатора на события/с, байты/с, активные сессии и суточный вклад
одного SKU обязательно численно заполняются до запуска. Достижение
лимита вызывает ограничение приёма или карантин с кодом причины.
Для графа сохраняется не более одного вклада сессии в поддержку
каждой пары, для популярности --- одного вклада сессии в SKU;
дополнительные ограничения SKU/источника регистрируются как часть
преобразования данных. Всплески поддержки рёбер и популярности проверяются по
предыдущим окнам; автоматическое включение подозрительных событий
в новый пакет запрещено до проверки карантина. Sybil-нагрузка,
атака на один SKU, огромный блок и массовое восстановление истории входят в обязательные
испытания. Эти меры ограничивают воздействие, но не гарантируют
распознавание всех ложных пользователей. Неизвестные эксплуатационные
лимиты означают неготовность рабочего профиля, а не неограниченный приём.

Минимальная обязательная запись ответа содержит псевдонимный ключ
сессии, версии, выбранную ветвь, идентификаторы выданных товаров,
исход и технический статус. Полные кандидаты, баллы и подробные
причины сохраняются только в ограниченном диагностическом журнале;
сырые запросы, векторы и история в обычный журнал не попадают.
Псевдонимный ключ строится ключевым преобразованием с разделением
пространств; хеш идентификатора не считается анонимизацией.
Для связности A/B ключ назначения стабилен на всём горизонте
эксперимента, а ключи диагностических журналов вращаются независимо.

Начальная политика хранения: состояние удаляется через 24 часа
неактивности; $\Delta_{\mathrm{retry}}=24$ часа,
$\delta_{\mathrm{late}}=30$ минут; журнал для воспроизведения хранится 48 часов,
подробные трассы --- 7 суток. Минимальные записи назначения и исхода
хранятся до заранее установленной даты закрытия анализа плюс 30 суток,
после чего удаляются либо заменяются агрегатами без сессионных ключей.
Метки повторной доставки живут независимо от срока хранения состояния и не меньше
срока повторной доставки. Сроки и владельцы хранилищ включаются
в паспорт; эти значения являются проектным профилем, а не
утверждением о правовых сроках. Доступ разделён по ролям,
передача и хранение шифруются, чтение диагностических трасс аудируется.
Процедура удаления охватывает состояние, журнал воспроизведения, диагностические копии
и резервные копии по объявленному сроку истечения; восстановление
не возвращает удалённые записи без повторного применения реестра удалений.
Публичные отчёты и пользовательские объяснения не раскрывают
идентификаторы сессий, сырые запросы или последовательность действий.

\section{Вычислительный бюджет и обслуживание}\label{sec:operations-serving}

\paragraph{Объём демонстрационной версии.}
Для портфолио достаточно одного локального процесса с проверенным пакетом,
состоянием, повторяемыми ответами, резервом, объяснениями и журналом.
Временные границы, допустимость товаров, входные ограничения и запрет
поздней публикации сохраняются. Промышленная нагрузка
(раздел~\ref{sec:ops-load-profile}), RTO/RPO и восстановление нескольких
экземпляров (раздел~\ref{sec:ops-recovery}), рабочие подписи и мониторинг
не входят в её обязательную приёмку и требуют отдельной задачи.
Поля рабочего профиля \code{serving}/\code{operations} паспорта,
\code{post_load} и $S_{\mathrm{serve}}$ относятся к этой отложенной
эксплуатационной приёмке, а не к демонстрации.

\subsection{Подготовка, обучение и ответ на запрос}
\label{sec:ops-compute-stages}

Отдельно измеряются подготовка версии, обучение и обслуживание.
Подготовка первого запуска создаёт граф, содержательные векторы
и справочники; расширения добавляют пути, обучаемые векторы,
индексы приближённого поиска соседей (ANN) и дерево SID.
Обучающие преобразования и шум кодов используются только при обучении.
Запрос использует готовые пути и товарные векторы, вычисляя состояние
сессии, кандидатов, признаки и ранжирование. Поиск по признакам товаров
в первом запуске точный; кратчайшие пути заново не вычисляются.
Авторегрессивное декодирование SID относится к бюджету обслуживания
отдельно включённого исследовательского варианта.

\subsection{Состав версии}

Публикуется согласованный пакет
\begin{equation}\label{eq:ops-package}
\mathsf V=(\nu_D,\nu_{\mathcal I},\nu_G,\nu_{\mathrm{vec}},
\nu_{\mathrm{features}},\nu_f,\nu_{\mathrm{policy}},
\nu_N,\nu_{\mathrm{SID}}).
\end{equation}
Граф, векторы и ранжировщик заменяются атомарно.
Совместимость векторных пространств проверяется по происхождению,
не только размерности.

Версия $\nu_N$ объединяет топологию непрямых связей, веса кодировщика,
таблицу обученных товарных векторов и её поисковый индекс.
Версия $\nu_{\mathrm{SID}}$ содержит кодировщик, словари кодов,
декодер, соответствие кодов товарам и дерево допустимых префиксов.
Отключённый модуль получает версию \code{not_used};
в первом запуске это относится к $\nu_N$ и $\nu_{\mathrm{SID}}$.
Замена словаря без обновления соответствующего дерева
и декодера считается несовместимым пакетом.

Паспорт пакета содержит SHA-256 каждого артефакта, точные версии
библиотек и lockfile, commit сборки, cutoff данных и время публикации,
идентификатор сборщика и зависимости преобразований. В рабочем сервисе
паспорт подписывается доверенным издателем (начальный профиль:
Ed25519); ключ доверия поступает из отдельного доверенного хранилища,
а не из проверяемого пакета. До активации проверяются подпись,
все хеши, совместимость схем и разрешённость версии. Проверка одного
хеша без аутентификации издателя недостаточна. Артефакты читаются
декларативным загрузчиком без исполнения сериализованного кода;
пути вне пакета, неподдерживаемые типы и превышение заявленных
размеров отклоняются. Локальный исследовательский пакет допускает
профиль \code{local_checksum_only}, явно не подтверждающий
подлинность рабочего пакета. Не прошедший проверку пакет не активируется.

Запрос выполняет последовательность:
\begin{enumerate}
\item проверяет структуру входа и фиксирует опубликованный пакет;
\item получает допустимый каталог, проверяет версию преобразований
состояния и атомарно обновляет состояние;
\item для включённого нейросетевого варианта один раз вычисляет
представление сессии и поисковые веса намерений;
\item формирует кандидатов и их происхождение;
\item вычисляет признаки и баллы;
\item применяет ограничения итогового списка;
\item записывает ответ, задержку и состояние исполнения.
\end{enumerate}
Резерв входит в \(\pi\) с фактически наблюдаемой частотой использования.

\subsection{Память и число операций}

Пусть \(n_I\) --- число товаров, \(|\mathcal R_G|\) --- число отношений графа,
\(d_G\) --- сохраняемое число соседей, \(d_T,d_V\) --- размерности
содержательных векторов, \(p\) --- число признаков, \(T_f\) --- число деревьев,
\(D_f\) --- их предельная глубина. Порядки основных затрат:
\begin{align}
\mathrm{Mem}_{G}&=O(|\mathcal R_G|n_Id_G),\label{eq:ops-graph-memory}\\
\mathrm{Mem}_{\mathrm{content}}&=O(n_I(d_T+d_V)),\label{eq:ops-content-memory}\\
\mathrm{Work}_{\mathrm{graph},q}&=O(|\mathcal R_G|L_Hd_G),\label{eq:ops-graph-work}\\
\mathrm{Work}_{\mathrm{rank},q}&=O(Mp+MT_fD_f).\label{eq:ops-rank-work}
\end{align}
Точный содержательный поиск по каталогу требует
\(O(n_I(d_T+d_V))\) скалярных операций на запрос.
При необходимости приближённый индекс сравнивается с точным
до открытия $D$. Полный граф на \(n_I^2\) пар не материализуется.

\paragraph{Затраты необязательных расширений.}
Товарные векторы размерности $d$ требуют $O(n_Id)$ памяти
и пересчитываются до публикации индекса. Представления сессии
вычисляются однократно для источников и признаков ранжирования.

Подготовка непрямых рёбер ограничена длиной пути, числом расширений
и числом сохраняемых соседей. Детерминированное досрочное прекращение
является частью модели; его частота по вершинам входит в отчёт.

Точная нормировка по каталогу требует $O(n_Id)$ операций на запрос;
ANN-индекс её не ускоряет. Обучение хранит также промежуточные значения
для производных. Блочное накопление логарифма суммы экспонент сохраняет
точную потерю. Выборка отрицательных товаров задаёт другую процедуру
(приложение~\ref{sec:sampled-overview}) и не отменяет пересчёт товарных
представлений или глобальные регуляризаторы.
Измеряются полное время обучения и пиковая память RAM и ускорителя.
Исчерпание бюджета считается вычислительным отказом;
число эпох не меняется без объявления.

\paragraph{Контракт приёмки ANN.}\label{sec:ops-ann-contract}
Для каждого приближённого источника измеряется его собственное качество
\begin{equation}\label{eq:ops-ann-recall}
\operatorname{IndexRecall}@m=\operatorname{Agreement}@m
\end{equation}
из~\eqref{eq:model-ann-agreement} на фиксированной probe-выборке.
Это совпадение результатов поиска соседей, а не показатель покупки.
Точный и приближённый поиск используют одинаковые векторы, фильтр,
квоту $m$ и порядок равных баллов. Приближённые результаты повторно
оцениваются точным скалярным произведением, устраняются повторы
и проверяется допустимость; исходные доли недопустимых и повторных
результатов публикуются, в окончательном списке обе равны нулю.
Детерминированное расширение поиска после фильтрации имеет заданный
предел; исчерпание предела не разрешает снять фильтр или увеличить $M$.

Дополнительно вычисляются set precision с числителем
$|C_q^{\mathrm{approx}}\cap C_q^{\mathrm{exact}}|$ и знаменателем
$|C_q^{\mathrm{approx}}|$, совпадение exact top-1 и разрыв баллов
\[
\operatorname{Regret}_q=
\max(0,r_q(i_q^{\mathrm{exact}})-r_q(i_q^{\mathrm{approx}})).
\]
Здесь $r_q(i)$ --- скалярное произведение исходных представлений
запроса и товара данного источника до деления на температуру и
иных масштабных множителей; их нормы не превышают 1. Для содержательного
поиска это косинус, для нейросетевого --- числитель его формулы балла.
Паспорт фиксирует \code{regret_score=unscaled_bounded_dot_v1} и ссылки
на обе формулы представлений. При непустых ответах regret лежит
в $[0,2]$ и не меняется от положительного масштаба ранжирующего балла.
При непустом точном и пустом приближённом ответе precision, recall
и top-1 совпадение равны нулю, regret равен $+\infty$;
поэтому на контрольных запросах такие пустые ответы не допускаются.
Пустой точный ответ учитывается отдельно и требует пустого ответа
ANN при тех же масках и фильтре. Стратификация фиксируется
по маскам доступных модальностей и популярности точного top-1:
популярная группа --- первые $\lceil0.1|\mathcal I_G|\rceil$ товаров
по числу сессий в $A_G$ по убыванию, при равенстве по ID;
остальные образуют вторую группу. Граница фиксируется до проверки ANN.

Начальные инженерные пороги: средние IndexRecall и precision
не ниже $0.98$, top-1 agreement не ниже $0.99$, средний regret
не выше $0.01$. В каждом заранее объявленном непустом сегменте
recall и precision не ниже $0.95$, top-1 agreement не ниже $0.95$,
regret не выше $0.02$. Требуются не менее 10\,000 непустых точных
ответов в целом и 1000 на сегмент; меньшая выборка означает
\code{insufficient}, а не пройденный сегмент. На десяти повторениях
одних входов списки после точного пересчёта совпадают точно;
начальное состояние генератора, индекс, число потоков
и порядок равных баллов входят в паспорт.
Это эмпирические пороги на фиксированной контрольной выборке из
$C_{\mathrm{smoke}}$, а не гарантии для генеральной совокупности.
Параметры индекса выбираются до $C_{\mathrm{select}}$.

ANN также проходит общий профиль P95/P99, отказов и резервных ответов
и общий предел памяти раздела~\ref{sec:ops-load-profile}, включая индекс.
На $C_{\mathrm{select}}$ изменение Hit@20 относительно точного поиска
публикуется на общей $\mathcal S_{C_{\mathrm{select}}}^+$.
Допустимое снижение $\delta_{\mathrm{ANN}}$ задают в долях Hit@20
до сравнения; без заданного допуска ANN не участвует в выборе.
Этот фильтр выбора не заменяет итоговый интервал на $D$.
Если ANN отключён, его статус --- \code{not_used}.

\paragraph{Таблица обучающих примеров.}
Размер таблицы кандидатов, до \(N_{\mathrm{train}}Q_{\max}M\) строк,
оценивается до запуска. Для выборки запросов или отрицательных кандидатов
фиксируются вероятности включения и потеря; оценочный каталог
ради сокращения вычислений не уменьшается.

\subsection{Задержка и резервные ветви}
\label{sec:ops-deadlines}

Пусть \(\ell_{qj}\) --- время последовательного этапа,
\(\ell_q^{\mathrm{wait}}\) --- ожидание обслуживания. Тогда
\begin{equation}\label{eq:ops-latency}
\ell_q=\ell_q^{\mathrm{wait}}+\sum_{j=1}^{J}\ell_{qj}.
\end{equation}
Для порогов \(b_j\) и \(b_0\) без предположения независимости
\begin{equation}\label{eq:ops-latency-bound}
\Prob\!\left(\ell_q>b_0+\sum_j b_j\right)
\le \Prob(\ell_q^{\mathrm{wait}}>b_0)+\sum_j\Prob(\ell_{qj}>b_j).
\end{equation}
Процентили этапов не складываются в процентиль общей задержки.
Фиксируются частота запросов, параллельность, каталог и доля прогретых состояний.

Резервный список готовится в начале общей процедуры кандидатов.
При исчерпании бюджета основной ветви до общего срока ответа
используется уже подготовленная допустимая ветвь с резервом времени
на запись и возврат; при её отсутствии возвращается пустой список
с кодом причины. Достижение общего срока не запускает новые вычисления.
Оценка не удаляет такой запрос. Резерв образует отдельную ветвь
$\mathcal C_q^{\mathrm{fallback}}$ и не расширяет
$\mathcal C_q^{\mathrm{main}}$. При выборе готовой ветви с запасом
на запись и возврат фиксируются \code{branch} и
$\mathcal C_q^{\mathrm{served}}$; смешение двух множеств
для улучшения кандидатной границы запрещено. Всегда
$L_q\subseteq\mathcal C_q^{\mathrm{served}}\subseteq\mathcal I_q$
и $|\mathcal C_q^{\mathrm{served}}|\le M$. Если готовая ветвь
не выбрана, активное множество пусто. Если после выбора не выполнены
срок или обязательная запись, $L_q=\varnothing$ и
\code{delivery_status=error}, но зафиксированное множество кандидатов
сохраняется: на уровне сессии потеря доставки остаётся в
$U_s(\pi)-\operatorname{Hit}_s(\pi)$.

Для демонстрации начальный общий тайм-аут
$b_{\mathrm{deadline}}=1000$ мс фиксируется до запуска и охватывает
весь ответ, включая обязательную запись. Это предел ожидания,
а не гарантия производительности под нагрузкой. Жёстких сроков отдельных
этапов нет; запас для записи и возврата резерва --- параметр исполнения,
не самостоятельный порог приёмки. Промышленный профиль имеет отдельные
условия раздела~\ref{sec:ops-load-profile}.
Монотонные часы контролируют общий срок с запасом на возврат резерва.
Непрерываемый оператор ограничивается
заранее или изолируется; асинхронная отмена не гарантирует завершения.
Каждая операция использует остаток общего срока с учётом ещё обязательной
записи и возврата; нового независимого срока этапа не возникает.
Повтор обращения входит в тот же бюджет. После переключения ветви
поздний результат отменённой операции не меняет ответ и не перезаписывает
состояние; подтверждённая запись состояния остаётся валидной даже
при отказе ответа. Незавершённые операции освобождаются в изолированном
пуле с фиксированным пределом, новые вызовы при его заполнении
не принимаются. Отдельные коды: \code{invalid_input},
\code{overload_rejected}, \code{backend_timeout}, \code{cancelled},
\code{state_replay_limit}, \code{logging_failure}; один агрегат
\code{error} не заменяет эти причины.

SID запускается лишь при достаточном остатке общего бюджета поиска
кандидатов. Предел числа шагов и ширина поиска фиксированы.
При исчерпании времени действуют объявленные квоты и резервные правила,
без скрытого увеличения $M$. Частота пропуска SID и изменение качества
из-за ограничения времени публикуются отдельно.
Если резерв ещё не готов или обязательная запись не завершена,
возвращается код технического отказа; такой запрос остаётся
в знаменателях оценки.

Демонстрационный журнал связывает запрос, пакет, кандидатов, публичный
ответ и исход. Двухфазная фиксация и RPO=0 не обязательны; ошибка записи
или неизвестная доставка не объявляются успешным завершением.
Приёмка проверяет штатную основную выдачу, повтор без пересчёта,
достигнутый вычислительный тайм-аут с резервом и отказ при неполном состоянии.
HTTP-ответ сверяется с журналом. Полная задержка, загрузка и память
публикуются как измерения; P95/P99 и доля резерва не имеют обязательных
порогов для демонстрации. Обычное распределение процессорных ядер сохраняется.

\subsection{Воспроизводимый нагрузочный профиль}
\label{sec:ops-load-profile}

Отложенный промышленный профиль: применяется только по отдельной задаче,
не блокирует демонстрационный API. Его требования ещё не подтверждены испытанием.
\begin{center}
\begin{tabular}{p{0.42\linewidth}p{0.51\linewidth}}
\toprule
Параметр & Начальное значение\\
\midrule
Вычислительный узел & 8 виртуальных ядер CPU, 32 ГиБ RAM,
без GPU при обслуживании\\
Предлагаемый поток & 100 запросов в секунду\\
Генератор нагрузки & Независимое от ответов расписание;
не более 32 незавершённых запросов генератора\\
Прогрев & 5 минут перед измерением\\
Измерение & 30 минут, три отдельных повторения\\
Цель по задержке сервиса & P95 $\le50$ мс, P99 $\le100$ мс\\
Общая пиковая RAM процессов сервиса & Не более 28 ГиБ\\
Предельный срок ответа & 120 мс от поступления в сервис\\
Допустимые технические отказы & Не более $0.1\%$ всех запланированных запросов\\
Допустимая доля резервных ответов & Не более $1\%$\\
\bottomrule
\end{tabular}
\end{center}
Отчёт указывает модель CPU, процессы и потоки, версии библиотек,
тип чисел, каталог, параметры индекса и память. Одинаковое число ядер
не гарантирует одинаковой производительности; GPU проверяется отдельно.

\paragraph{Общий предел памяти.}\label{sec:ops-memory-limit}
Ограничение RAM относится ко всем вариантам, включая точный поиск,
индексы, локальное состояние, очереди и одновременно загруженные пакеты
при обновлении/откате. Начальная метрика --- максимум суммы PSS процессов
сервиса: общие страницы учитываются пропорционально их разделению.
Если PSS недоступен, допускается консервативная сумма пиковых RSS
всех процессов испытания; повторно учтённые общие страницы не вычитаются
по неподтверждённой оценке. Выбранный измеритель, период измерений и
состав процессов фиксируются до испытания и одинаковы у сравниваемых
вариантов. Превышение табличного предела означает \code{failed} независимо от ANN.

Распределение длины префиксов, наличия поиска, пропусков и повторных
обращений берут из входов $C_{\mathrm{smoke}}$ либо из заранее утверждённого
производственного профиля. Исходы $D$ для выбора нагрузки не используются.
Состояния сессий воспроизводят объявленные переходы;
повтор одного и того же прогретого запроса не заменяет поток.
Отдельно измеряют запуск с пустыми кэшами и восстановление после отказа.

Генератор планирует запросы каждые 10 мс независимо от ответов;
входы и начальное состояние одинаковы для всех вариантов.
Профиль всплесков объявляется отдельно. Если предел 32 мешает выдержать
100 запросов в секунду, учитываются недоставленные запросы и фактический
поток; замедленное испытание не подтверждает заявленную нагрузку.
Параллельность и очередь сервера задаются отдельно от предела генератора.
Для начального узла сервер допускает не более 16 активных запросов
и 16 ожидающих; приём входа и пребывание в очереди вместе ограничены 10 мс.
При заполнении очередь не растёт: возвращается
\code{overload_rejected} без дорогостоящего чтения состояния.
Пределы глобальны для узла, а не повторяются независимо в каждом
процессе. Отклонение и вызванный им пустой ответ учитываются
в операционных знаменателях; назначенная сессия остаётся в ITT.
Повтор рекомендационного запроса разрешён лишь с тем же ключом идемпотентности
и в пределах исходного срока; ответ сервера не запускает
неограниченный автоматический цикл повторов. Доставка событий подчиняется
разделу~\ref{sec:ops-transport-block} и не возобновляет срок старого ответа.

Задержка сервиса измеряется от первого байта запроса на его входе
до готовности полного ответа, включая ожидание, получение признаков
и обязательную запись. Внешняя сеть и обработка в клиенте входят
в отдельную сквозную задержку. Для планового момента отправки,
фактической отправки, поступления в сервис и завершения сохраняют
раздельные отметки времени; задержка генератора не скрывается
внутри неполного измерения.

Запросы отбираются по плановому времени отправки и наблюдаются
до завершения или срока, включая незавершённые на конце окна.
Для приёмочных процентилей отказ, недоставка или превышение срока
получают задержку $+\infty$; фактические времена до ошибки сохраняются.
Резервный ответ учитывается с измеренной задержкой и отметкой ветви.
Процентили только успешных основных ответов служат дополнительной диагностикой.
Эмпирический квантиль уровня $p\in(0,1]$ для $n>0$ значений
определён как $x_{(\lceil pn\rceil)}$ после сортировки, включая $+\infty$.
При пустом окне доли и квантили не определены: \code{insufficient},
не ноль. Журнал статусов сверяют с независимым счётчиком входа и
планом генератора; пропавшие записи учитывают как отказы, а не удаляют
из знаменателя. Это правило действует и для операционных окон $W$.

Каждое из трёх повторений должно удовлетворять целям по задержке,
отказам и резервным ответам; объединение повторений не скрывает
неудачный запуск. Отдельно публикуют пропускную способность,
максимальную очередь, фактическую параллельность и использование CPU/RAM.
Нагрузочная приёмка предшествует A/B-испытанию, а в самом эксперименте
задержки и отказы продолжают контролироваться как защитные показатели.

\subsection{Восстановление и откат}\label{sec:ops-recovery}

До публикации сохраняются последняя проверенная версия пакета
и её схема состояния. Начальное требование RTO отката --- 60 с
от решения оператора/защитного автомата до готовности проверенной
версии на всех экземплярах; RPO подтверждённых записей назначения
и исхода --- ноль. Это цели, проверяемые имитацией отказа,
а не заявленные измерения. Каждый запрос закреплён за одной версией;
откат не смешивает артефакты и не меняет ранее сохранённые ответы.
Несовместимое состояние восстанавливается или мигрирует по
проверенному правилу; до готовности действует допустимый резерв.
Возврат старой версии с просроченным каталогом запрещён.
Двойная загрузка версий входит в пиковую память; если она невозможна,
переключение выполняется поэтапно с той же проверкой RTO и доступности.
Испытание проверяет потерю процесса, недоступность состояния,
ошибку обязательной записи, повреждение пакета и откат с другой
версией признаков. Измеренные RTO/RPO и коды неуспешных запросов
публикуются отдельно от нагрузки в установившемся режиме.

\section{Наблюдение за качеством и изменениями данных}\label{sec:operations-monitoring}

\subsection{Немедленно наблюдаемые показатели}

Для окна запросов \(W\) определяются
\begin{align}
\widehat e_W&=\frac1{|W|}\sum_{q\in W}\ind\{\mathrm{error}_q=1\},\\
\widehat f_W&=\frac1{|W|}\sum_{q\in W}\ind\{\mathrm{fallback}_q=1\},\\
\widehat u_W&=\frac1{|W|}\sum_{q\in W}\ind\{\text{в префиксе есть неизвестный товар}\},\\
\widehat m_{W,j}&=\frac1{|W|}\sum_{q\in W}\ind\{x_{qj}\text{ отсутствует}\}.
\end{align}
Отдельно измеряются длина выдачи, доля пустых ответов,
задержки по процентилям, распределение числа кандидатов,
возраст снимка \(t_q-r(q)\) и доля несовместимых версий.
Для нейросетевой части контролируются доли отключённых модальностей,
нормы векторов, распределение весов намерений и расхождение ANN
с точным поиском. Для семантического модуля дополнительно учитываются
число используемых кодов, частоты совпадения кодов разных товаров,
усечение крупных групп и разрыв качества между мягким обучением
и окончательной дискретной выдачей.
Уменьшение задержки за счёт роста числа пустых ответов
не считается улучшением без учёта остальных ограничений.

\subsection{Изменение распределений}

Для числового признака границы интервалов определяются на обучающих
данных; пропуск образует отдельную категорию. Пусть \(p_b\) и \(q_b\)
--- сглаженные доли в исходном и текущем окнах, \(b=1,\ldots,B_h\):
\begin{equation}\label{eq:ops-hist}
p_b=\frac{n_b+\eps_h}{n+B_h\eps_h},\qquad
q_b=\frac{m_b+\eps_h}{m+B_h\eps_h},\qquad h_b=\frac{p_b+q_b}{2}.
\end{equation}
Расхождение Йенсена--Шеннона задаётся формулой
\begin{equation}\label{eq:ops-js}
\operatorname{JS}(p,q)=\frac12\sum_b p_b\log\frac{p_b}{h_b}
+\frac12\sum_b q_b\log\frac{q_b}{h_b}.
\end{equation}
Оно равно нулю для одинаковых распределений и не превосходит
\(\log2\). Это индикатор изменения входа; он не доказывает
ухудшение рекомендаций. Пороги устанавливаются по последовательным
контрольным окнам, без подбора после просмотра проблемного периода.

\paragraph{Реестр и правило решения.}\label{sec:ops-drift-rule}
Реестр включает входные признаки и пропуски, неизвестные товары,
маски источников, их квоты и размеры списков, долю уникальных кандидатов,
баллы и длину итогового списка, выбранную ветвь и причины отказа.
Категории и числовые границы замораживаются по ранним данным;
сглаживание~\eqref{eq:ops-hist} использует $\eps_h=1/2$.
Базовое окно содержит не менее 10\,000 запросов; текущие окна ---
непересекающиеся часы UTC с минимумом 1000 запросов и 200 сессий.
При меньшем объёме возвращается \code{insufficient}, окно не объявляется
стабильным. Отдельный счётчик отсутствия достаточной телеметрии
поднимает уведомление через три последовательных таких окна.

Пороги фиксируются совместно для всего реестра. До эксплуатации
на не менее 100 контрольных часовых окнах, предшествующих
$C_{\mathrm{select}}$ и не пересекающихся с базовым окном, вычисляют
$A_w=\max_j\operatorname{JS}_{w,j}/0.05$. Для окна $w$ используется
только ранее доступная база. Порог
$h_{\mathrm{drift}}=\max(1,Q_{0.99}(A_w))$ фиксируется вместе
с базой, набором признаков и правилом эмпирического квантиля nearest-rank.
Контрольные окна из $C_{\mathrm{smoke}}$ допустимы; до их открытия
заданы реестр и правило калибровки. Нехватка контрольных окон
означает непроверенный мониторинг, а не разрешение подобрать пороги
по $D$. На окнах настройки раздельно регистрируются доля превышений
$A_w>h_{\mathrm{drift}}$ и частота тревог полного правила ниже;
это внутривыборочные показатели, а не оценка ложных тревог.
Для такой оценки нужны отдельные окна и обоснованное предположение
отсутствия реального сдвига. Время суток, день недели и зависимость
окон учитываются в диагностике; будущая частота тревог не гарантируется.

Тревога возникает при $A_w>h_{\mathrm{drift}}$ в двух последовательных
достаточных окнах; снимается после трёх окон с
$A_w<0.8h_{\mathrm{drift}}$. Повторное уведомление по той же причине
отправляется не чаще раза в час, но новые отказы регистрируются сразу.
Максимум по всему реестру учитывает множественность в эмпирической
настройке; отдельные JS не интерпретируются как $p$-значения
и к ним не применяется FDR-коррекция как к тестам значимости.
Замена реестра, границ или базы требует новой калибровки на ранних окнах.
Адресат уведомлений и владелец разбора указываются в паспорте.

Тревога изменения распределений запускает проверку полноты, версий, источников
и выходных распределений, без автоматического переобучения.
Отдельные оперативные окна длиной 5 минут с минимумом 1000 запросов
сравниваются с начальными пределами технических отказов $0.1\%$
и резерва $1\%$. Превышение в двух последовательных окнах
включает заранее проверенный резерв/откат и приостанавливает
набор новых экспериментальных единиц; уже назначенные остаются в ITT.
Начальные защитные требования, независимо от объёма окна:
ту же защиту включают 20 последовательных
технических отказов валидных допущенных запросов либо три неуспешные
внутренние проверки работоспособности подряд. Проверки выполняются
каждые 5 с с отдельными диагностическими ID и сроком 120 мс; они не
изменяют пользовательское состояние и не входят в бизнес-знаменатели.
Для последовательности запросов используют порядок поступления;
незавершённый запрос получает исход по истечении своего срока.
Успех обнуляет соответствующий счётчик. Ошибки недоверенных или
некорректных входов не запускают общий откат исправной версии.
Неподлинный пакет блокируется до активации; единичная попытка его
загрузки не останавливает действующий сервис. Обнаруженное нарушение
допустимости фактической выдачи блокирует её немедленно.
Во время A/B действуют правила раздела~\ref{sec:eval-online}:
вывод обычного анализа с фиксированным сроком после защитной
остановки не считается подтверждающим.

Даты Coveo сдвинуты на целое число недель. Сохраняется порядок
дней недели. Связь с реальными праздниками не устанавливается
без исходных календарных дат.

\subsection{Показатели с задержанным исходом}

В исторической оценке обновляют Hit@20 для зарегистрированных целей
по общему календарному допуску. Неизвестное покрытие остаётся неизвестным.
Конверсию рабочего сервиса рассчитывают только при подтверждённом
полном учёте исходов для зрелых когорт, определённых временем назначения.
Учёт ранних положительных исходов при незавершённых отрицательных запрещён.

В эксплуатации можно наблюдать текущую конверсию
\(\widehat p_W=\sum_{s\in W}Y_s/|W|\).
Её изменение относительно прошлой недели смешивает эффект модели,
состав трафика и внешние обстоятельства. Причинный эффект
оценивается по одновременно работающей контрольной группе.
Сегменты подтверждающего анализа реальных показов определяются до назначения
варианта. Например, наличие поиска после назначения нельзя
использовать как основание для заявления о причинном эффекте
<<среди воспользовавшихся поиском>>.

\subsection{Обновление модели}

Изменение данных запускает диагностику: полнота журналов,
совместимость версий, наличие целей, изменение источников кандидатов.
Новая модель обучается только на истории, прошедшей календарный допуск
и проверку полноты чтения, проходит те же
проверки и публикуется новой версией. Сам факт срабатывания
\(\operatorname{JS}\)-порога не разрешает автоматически заменить модель.

Во время подтверждающего A/B-испытания поведенческие справочники
и модель по умолчанию заморожены. Обучение общей модели на откликах
обоих вариантов создаёт межгрупповое влияние; такой режим требует
отдельного экспериментального дизайна. Операционные исправления,
не меняющие предписанную политику, документируются.

\section{Объяснение результата и границы интерпретации}\label{sec:explanation}

Для каждого рекомендованного товара сохраняются источники попадания
в кандидаты, их баллы и версия снимка в диагностической трассе
с доступом и сроками раздела~\ref{sec:ops-security}.
Связанные события заменяются минимальными кодами причин;
их сырая последовательность не включается в публичное объяснение.
Графовое объяснение указывает отношение, товар префикса с наибольшим
вкладом и поддержку $n^r_{ij}$. Содержательное объяснение указывает использованную
модальность и измеренную близость. Поисковое объяснение ссылается
на доступный запрос и совместимое пространство векторов.

Объяснение описывает расчёт, а не причину покупки. Оно ограничено полями
Coveo: обезличенный вектор не раскрывает название, назначение,
свойства товара или совместимость с конкретной моделью телефона.

Первый вариант использует шаблоны по трассе расчёта.
Языковая модель не расширяет доступные факты; ранжирование на её основе
требует каталога с текстами и отдельной проверки качества, стоимости
и доступности признаков.

\section{Полнота результата и условия приёмки}\label{sec:acceptance}

Историческую оценку принимают по протоколу \code{observed_purchases_v1}
раздела~\ref{sec:eval-observed}. Её успех не доказывает рост покупок
и сам по себе не разрешает A/B-испытание.

\subsection{Состояния результата}

Состояние представляется кортежем
\begin{equation}\label{eq:acceptance-state}
S=(S_{\mathrm{data}},S_{\mathrm{math}},S_{\mathrm{offline}},
S_{\mathrm{serve}},S_{\mathrm{business}}).
\end{equation}
Каждая компонента принимает значение
\(\{\mathrm{not\_run},\mathrm{passed},\mathrm{failed},
\mathrm{insufficient}\}\). Для бизнес-эффекта дополнительно
используется \(\mathrm{not\_observed}\), если рандомизированного
испытания не было. Ошибка и недостаточность данных --- разные состояния.

\begin{longtable}{L{0.16\textwidth}L{0.73\textwidth}}
\toprule Компонента & Необходимые условия \code{passed} \\\midrule\endhead
$S_{\mathrm{data}}$ & Паспорт заполнен для текущего этапа; хеши и схема
совпали; правило канонизации зафиксировано; пересечение сессий периодов
пусто; выбранные сессии полностью прочитаны из поставки, календарный
допуск и зарегистрированные цели рассчитаны верно. Неизвестная полнота
телеметрии сохранена, включая сессии с покупкой; исключения перечислены. \\
$S_{\mathrm{math}}$ & Все тесты включённых компонентов, обязательные
на текущем этапе по~\ref{app:tests}, пройдены в заданных допусках, включая производные,
полноту масок, границы времени и тождество потерь охвата. \\
$S_{\mathrm{offline}}$ & Замороженный протокол исполнен на $D$, имеется
$N=N_D^+>0$ и независимый полный отчёт по $\operatorname{Hit@20}_D$,
включая сравнение чувствительности
без метаданных. Режим воспроизведения и полнота трассы отвечают
разделу~\ref{sec:ops-replay-trace}. Для каждой заявленной оценки
выполнены предположения, пороги информации и невырожденности
из~\ref{sec:eval-information}, минимум блоков и требуемая полуширина:
$h_{\mathrm{pair}}$ для пары, $h_{\mathrm{single}}$ для одной политики,
оба в долях.
Знак разности сообщается отдельно: \code{passed} означает пригодность
оценки, а не превосходство модели. \\
$S_{\mathrm{serve}}$ & Каждый из трёх нагрузочных прогонов удовлетворяет
всем порогам~\ref{sec:ops-load-profile}, включая общий предел и учёт
пиковой памяти по~\ref{sec:ops-memory-limit}; пройдены проверки
совпадения признаков, отказа, восстановления, целостности и входных ограничений.
Частичный отказ источника допустим лишь внутри зарегистрированной
резервной политики и её пределов. \\
$S_{\mathrm{business}}$ & Полный разрешённый анализ A/B выполнен без
нарушения назначения и учёта, нижняя граница выше заранее заданного
бизнес-порога и выполнены защитные ограничения. \\
\bottomrule
\end{longtable}
При обнаруженном нарушении контракта или измеренном превышении
технического порога ставят \code{failed}; при корректной процедуре,
но недостатке сессий с зарегистрированной целью, блоков или точности --- \code{insufficient}.
Для бизнес-компоненты при пригодном интервале $[\ell_B,u_B]$ и выполненных
остальных условиях статус равен \code{passed} при
$\ell_B>\delta_{\mathrm{accept}}$, \code{failed} при
$u_B\le\delta_{\mathrm{accept}}$ и \code{insufficient} при
$\ell_B\le\delta_{\mathrm{accept}}<u_B$;
порог определён в~\ref{sec:eval-acceptance}. Если одновременно известны
нарушение и недостаток данных, \code{failed} имеет приоритет;
все причины сохраняются.
\code{not_run} означает, что этап ещё не выполнялся.

Этапы запуска определены в~\ref{sec:eval-stages}. Пока $D$ закрыт,
$S_{\mathrm{offline}}=\mathrm{not\_run}$, независимо от успеха отладки.
Для локального \code{ideal_deterministic} состояние
$S_{\mathrm{offline}}$ не зависит от $S_{\mathrm{serve}}$;
невыполненная приёмка сервиса сохраняет $S_{\mathrm{serve}}=\mathrm{not\_run}$.
Отчёт включает разложение потерь~\ref{sec:eval-loss-diagnosis}
и вклад источников~\ref{sec:eval-source-contribution}; технические сбои
образуют отдельный разрез, не дополнительное слагаемое потерь охвата.

Для пилотного A/B-испытания дополнительно требуется проверенная
полнота учёта исходов эксперимента, обозначенная $O_{\mathrm{AB}}=1$:
\begin{equation}\label{eq:acceptance-pilot}
\mathsf{Ready}_{\mathrm{pilot}}=
\ind\{O_{\mathrm{AB}}=1\}\,
\ind\{S_{\mathrm{data}}=S_{\mathrm{math}}=
S_{\mathrm{offline}}=S_{\mathrm{serve}}=\mathrm{passed}\}.
\end{equation}
Для этой формулы $S_{\mathrm{serve}}$ должен быть получен именно
в эксплуатационном профиле: ограничения арендатора и товаров заполнены,
подпись пакета проверена, действуют правила хранения и доступа,
мониторинг откалиброван на достаточном числе окон. Успех локального
профиля \code{local_checksum_only} либо \code{insufficient} мониторинга
не разрешает пилот. В отчёте явно указывают профиль приёмки.
Доказанное бизнес-превосходство дополнительно требует положительного
решения по \(\Delta B\) согласно пункту~\ref{sec:eval-acceptance}
и всем защитным ограничениям. Мощность относится к заранее заданному
плану, а не к вычисляемой после эксперимента характеристике результата.
Полнота этой спецификации не означает, что система уже прошла испытания.

\subsection{Результат одного запроса}

\begin{equation}\label{eq:acceptance-request-result}
\mathsf R_q=(q,s,t_q,\mathsf V,L_q,\mathcal C_q,
\mathrm{branch},\mathrm{delivery\_status},
\mathrm{scores},\mathrm{sources},\mathrm{status},\mathrm{latency}).
\end{equation}
Это логическая диагностическая структура, не требование сохранять
полные кандидаты и баллы каждого запроса рабочего сервиса бессрочно.
Минимальная обязательная запись и ограниченный диагностический журнал
разделены по~\ref{sec:ops-security}. Зарегистрированные цели присоединяют
только при оценке; исходы рабочего эксперимента --- после созревания.
Кандидаты воспроизводятся по зафиксированному режиму
и трассе из~\ref{sec:ops-replay-trace}. Без полной трассы фактические
деградации и $\mathcal C_q^{\mathrm{served}}$ не считаются измеренными;
исправный повторный расчёт не восстанавливает их автоматически.
Исторический отчёт сохраняет набор сессий, зарегистрированные цели,
неизвестный статус полноты телеметрии, параметры отбора,
интервалы, отказы и стоимость. Идентификатор запуска объединяет версии
данных, конфигурации, кода и зависимостей. Смена начального состояния
генератора не создаёт независимую выборку посетителей.

\subsection{Замкнутая вычислительная постановка}

Для конфигурации $\xi$ последовательно строятся исторические артефакты
$\mathsf A_G$, параметры отбора $\theta_\xi$ и ранжировщик $f_\xi$.
Обозначим процедуры их подгонки через
$\mathcal T_G,\mathcal T_N,\mathcal T_R$:
\begin{align}
\mathsf A_G&=\mathcal T_G(D_{A_G};\xi),\\
\theta_\xi&=
\begin{cases}
\varnothing,&\text{нейросетевое расширение отключено},\\
\mathcal T_N(D_{A_N};\mathsf A_G,\xi),&\text{расширение включено},
\end{cases}\\
\mathcal C_q^\xi&=\mathcal G_\xi(\mathcal H_q,\mathsf A_G,\theta_\xi,\mathcal I_q),\\
x_{qi}^\xi&=\Phi_\xi(\mathcal H_q,i,\mathsf A_G,\theta_\xi),\\
f_\xi&=\mathcal T_R(\{(x_{qi}^\xi,y_{qi}):s\in\mathcal S_B^+,
q=q_s^{(1)},\ i\in\mathcal C_q^\xi\}),\\
\widehat\xi&\in\argmax_{\xi\in\Xi_{\mathrm{feasible}}}
\operatorname{Hit@20}_{C_{\mathrm{select}}}(\pi_\xi).
\end{align}
История для $\mathcal T_G$, первые запросы сессий с зарегистрированной
целью для $\mathcal T_R$ и их цели заданы в разделе~\ref{sec:eval-observed}.
В первом запуске $\Xi_{\mathrm{feasible}}=\{\xi_0\}$:
выполняется одна заранее заданная конфигурация, а не подбор
по результатам $C_{\mathrm{select}}$. Период $A_N$ остаётся зарезервированным;
его данные не дополняют граф и справочники базовой системы.
Конфигурации последующих сравнений и правило выбора фиксируются
до открытия $C_{\mathrm{select}}$.
В $\Xi_{\mathrm{feasible}}$ входят только варианты, прошедшие объявленные
ограничения. При точном равенстве Hit@20 выбирают меньшую медиану P99
трёх одинаковых нагрузочных прогонов, затем меньшую пиковую RAM,
меньшее число обучаемых параметров и лексикографически меньший
\code{config_id}. Допуск равенства Hit@20 не подбирается по результатам:
на общих сессиях сравниваются целые суммы попаданий.
Если допустимых вариантов нет, выбор имеет статус \code{insufficient}
или \code{failed} по установленной причине и победитель не назначается.
Итоговый список сортируется по $f_{\widehat\xi}$, если заранее не выбрана
отдельная процедура разнообразия. Период $D$ не входит ни в одну
процедуру подгонки или выбора.
В варианте SID в $\theta_\xi$ включены дополнительный кодировщик,
декодер и соответствие кодов товарам, обученные до $B$.

Бизнес-функционал~\eqref{eq:business-optimization} проверяется
с реальными показами по приложению~\ref{sec:eval-online}.

\clearpage\part{Нейросетевое расширение}\label{part:neural}
\section{Основания и порядок нейросетевого расширения}\label{sec:research}

Часть используется после базового запуска на с.~\pageref{sec:first-run}.
Она добавляет один обучаемый отбор; несколько намерений и товарные
векторы являются его компонентами, а не отдельными обязательными этапами.

\subsection{Связь с работами 2024--2026 годов}

В таблице указаны заимствованные принципы. Формулы ниже задают
собственную адаптацию для Coveo. Статус источников проверен
по состоянию на 19 сентября 2026 года.

\begin{longtable}{L{0.26\textwidth}L{0.64\textwidth}}
\toprule Работа & Принцип и роль в проекте \\\midrule\endhead
SPARE, ACM TORS, 2025 \cite{spare2025} &
Кратчайшие пути и компактное графовое распространение.
Типы связей, их стоимость и покупочная цель адаптированы. \\
MiaSRec, SIGIR 2024 \cite{miasrec2024} &
Несколько намерений сессии. Зависимость их весов от поиска
добавлена в настоящей модели. \\
MACL, Pattern Recognition, 2026 \cite{macl2025} &
Контрастивное согласование модальностей; в Coveo доступны готовые
векторы. \\
DIGER, SIGIR 2026 \cite{diger2026}; GRID, 2025 \cite{grid2025} &
Обучаемые семантические коды и сравнение способов кодирования.
Используются только в отдельном исследовании. \\
\bottomrule
\end{longtable}

Результаты этих работ относятся к их данным, продуктам и целям;
они не задают ожидаемого прироста покупок на Coveo.

\subsection{Конфигурация и хронология обучения}

Нейросетевая политика $\pi_{\mathrm{main}}$ расширяет
$\pi_{\mathrm{base}}$ обучаемыми графовыми и содержательными векторами,
намерениями сессии, поисковыми весами и новым источником кандидатов.
Исследовательская $\pi_{\mathrm{SID}}$ дополнительно использует
семантические коды. В обоих случаях общий предел $M$ сохраняется.

Порядок обучения и календарный допуск целей заданы в~\ref{sec:modern-time}.
Граф публикуется в $r_G$ до $A_N$, обученный отбор --- в $r_N$ до $B$.
Подгонка на $A_N$ не является независимой оценкой на нём.
Перед обучением ранжировщика на $B$ кодировщик и товарные векторы
замораживаются; градиент деревьев в них не проходит.

\subsection{Порядок реализации и контролируемые расширения}
\label{sec:research-stages}

Первое обучение использует конфигурацию~\ref{app:neural-config}.
Сравнение одного и двух намерений проводится по плану~\ref{app:ablation}.
При превышении бюджета точной полнокаталожной потери допускается
выборочный режим~\ref{sec:sampled-overview}; его качество проверяется
отдельно. Он не получает дополнительной квоты кандидатов.

Для контроля обработки истории приложение~\ref{sec:sequence-overview}
сохраняет товарный граф и входные сведения, заменяя кодировщик сессии.
Вариант с дополнительными действиями корзины или последовательностью
поисковых запросов проверяется отдельно: он меняет также состав входов.
Польза самого графа проверяется отключением графовой ветви.

Основание для следующего опыта задаёт диагностика~\ref{sec:eval-next-step}.
Во всех сравнениях фиксируются контроль, изменяемый компонент,
число запусков, итоговый охват, неопределённость разности и стоимость.
Меньший разрыв до кандидатной границы сам по себе не является
критерием улучшения.

\subsection{Гипотезы после диагностики}\label{sec:research-hypotheses}

Следующие варианты не входят в обязательный рецепт и не имеют
заранее установленного выигрыша.
\begin{description}
\item[История поиска.]
При потерях в сессиях с несколькими запросами проверить число
доступных запросов, изменение их векторов и согласованность с просмотрами.
Все признаки вычисляются по доступной истории; смена интереса
не считается ошибкой поискового сигнала. Длина исходного текста
не восстанавливается из предоставленного вектора.
\item[Трудные отрицательные примеры.]
Исследовать изменение распределения выборки с условиями
приложения~\ref{sec:sampled-training}; сравнивать итоговые политики,
а не численные значения разных приближённых потерь.
\item[Категорийное сглаживание.]
При потерях на редких товарах проверить смесь строки графа
с нормированным категорийным распределением. Для взвешенных $W_{ij}$
байесовская интерпретация требует отдельной модели наблюдений.
Полный приор нельзя разворачивать в плотную матрицу каталога:
нужна ограниченная опора или отдельный источник кандидатов.
\end{description}

Нейросетевой ранжировщик рассматривается после диагностики потерь
при фиксированном отборе (раздел~\ref{sec:model-ranking}).
Прогноз вероятности покупки и его калибровка требуют другого набора
наблюдений (раздел~\ref{sec:model-calibration}).
Расширение каталога требует отдельного правила допуска новых товаров
из приложения~\ref{sec:freshness}; метаданные сами не задают дату наличия.

\section{Обучаемая модель графа, содержания и намерений сессии}
\label{sec:modern-model}

Правила данных и обучения заданы в~\ref{sec:modern-time}, квоты
и настройки --- в~\ref{app:neural-config}.
Выборочный знаменатель~\ref{sec:sampled-overview}
не меняет архитектуру представлений.

\subsection{Какие величины обучаются и когда}
\label{sec:modern-time}

Каталог $\mathcal I_G$, графы $P^r$, пути и статистики строятся по $A_G$
и фиксируются. Параметры $\theta_N$ обучаются на $A_N$ по первым запросам
с календарным допуском~\eqref{eq:eval-admission}. Обучение ранжировщика,
выбор конфигурации и итоговая оценка следуют общему порядку:
\begin{equation}\label{eq:modern-periods}
A_G\prec A_N\prec B\prec C\prec E\prec D.
\end{equation}
Символ $\prec$ означает строгое временное предшествование с календарным
запасом $H+g$. Период $E$ зарезервирован и не используется.
Одна сессия не входит в два периода; полнота телеметрии остаётся неизвестной.

Соседи определяются только по $A_G$: граф из всего $A_N$ раскрывал бы
будущие связи его ранним префиксам. Число эпох, размерность и коэффициенты
потерь первого обучения фиксированы. Последующие сравнения на $C_{\mathrm{select}}$
следуют заранее заданному плану; каждый вариант обучается только на $A_N$
без дообучения на $C$ или $E$.

\subsection{Кратчайшие пути в фиксированном графе}
\label{sec:modern-paths}

Для поиска непрямых связей объединим известные отношения:
\begin{equation}\label{eq:modern-path-base}
\bar P_{ij}=\sum_{r\in\mathcal R_G}\zeta_rP^r_{ij},\qquad
\zeta_r\ge0,\quad\sum_r\zeta_r=1.
\end{equation}
Коэффициенты $\zeta_r$ входят в конфигурацию; суммы строк могут быть
меньше единицы из-за пустых отношений. При $\bar P_{ij}>0$
зададим безразмерную длину ребра
\begin{equation}\label{eq:modern-path-edge}
\ell_{ij}^{G}=c_G-\log\bar P_{ij},\qquad c_G>0.
\end{equation}
Сильная связь уменьшает длину, $c_G$ штрафует переход.
Эта геометрия не задаёт вероятности последовательностей покупок.

Для направленного пути $p=(i=i_0,i_1,\ldots,i_h=j)$ его длина
равна сумме длин рёбер. Ограничим число переходов числом $H_P$:
\begin{equation}\label{eq:modern-distance}
D_{ij}^{G}=
\min_{\substack{p:i\leadsto j\\1\le |p|\le H_P}}
\sum_{(u,v)\in p}\ell_{uv}^{G}.
\end{equation}
При отсутствии пути минимум равен $+\infty$; $i=j$ не рассматривается.
Все длины положительны, поэтому кратчайший путь не содержит цикла.
Расстояние направленное: обычно $D_{ij}^{G}\ne D_{ji}^{G}$.
Для каждой вершины сохраняются не более $d_P$ товаров
с наименьшим $D_{ij}^{G}\le R_P$, $R_P>0$.
Полученное множество $\mathcal N_i^{P}$ и расстояния замораживаются.
Равные расстояния разрешаются порядком идентификаторов.

Ограниченный поиск учитывает состояние <<вершина, число переходов>>.
Обычный поиск с единственной меткой на вершину может пропустить
допустимый путь при ограничении $H_P$.
В приближённом варианте раскрываются не более $B_P$ состояний
на исходную вершину, по возрастанию длины с фиксированным разрешением
равенств. При исчерпании бюджета сохраняются найденные допустимые соседи
и признак усечения. Их расстояния $\widetilde D_{ij}^{G}$ --- длины
найденных путей, не обязательно точные минимумы.
Далее $D_{ij}^{G}$ обозначает вариант, заданный конфигурацией:
точный или приближённый. Он общий для версии модели и вычисляется заранее.

Учебный пример: $\bar P_{ab}=0.8$, $\bar P_{bc}=0.7$,
$\bar P_{ac}=0.1$, $c_G=0.1$. Длина ребра $a\to c$ равна примерно
$2.403$, пути $a\to b\to c$ --- $0.780$; при $H_P\ge2$
связь с $c$ может сохраняться через $b$.

\subsection{Обучаемые связи и раздельные представления товара}
\label{sec:modern-item}

Для $i\in\mathcal I_G$ обучаемый вектор идентификатора $b_i^N\in\R^d$
инициализируется по разделу~\ref{sec:modern-training-recipe}. Введём нормировку
\begin{equation}\label{eq:modern-normalize}
\mathsf n_\eps(v)=\frac{v}{\max(\eps_N,\norm{v})},\qquad\eps_N>0.
\end{equation}
Она определена и в нуле; норма результата не превышает единицы.
Маска пригодности направления равна $\ind\{\norm{v}>\eps_N\}$.

Историческое распределение по сохранённым путям и его обучаемая
поправка задаются следующим образом:
\begin{align}\label{eq:modern-edge-attention}
\nu_{ij}^{P}
&=\frac{\exp(-D_{ij}^{G}/\tau_P)}
{\sum_{v\in\mathcal N_i^P}\exp(-D_{iv}^{G}/\tau_P)},\\
A_{ij}^{N}
&=\frac{\nu_{ij}^{P}\exp(\eta_N\langle
\mathsf n_\eps(b_i^N),U_N\mathsf n_\eps(b_j^N)\rangle)}
{\sum_{v\in\mathcal N_i^P}\nu_{iv}^{P}\exp(\eta_N\langle
\mathsf n_\eps(b_i^N),U_N\mathsf n_\eps(b_v^N)\rangle)}.
\nonumber
\end{align}
Здесь $\tau_P>0$, $\eta_N\ge0$ --- параметры конфигурации,
$U_N\in\R^{d\times d}$ обучается на $A_N$.
Формулы применяются только к непустому $\mathcal N_i^P$;
в остальных случаях все $A_{ij}^{N}$ равны нулю.
Обучаются веса фиксированных связей.

\paragraph{Логарифмы вероятностей.}\label{sec:modern-log-probabilities}
На непустой допустимой опоре сначала находят
$\widetilde z_j=z_j-\max_v z_v$, затем
$\operatorname{logsoftmax}(z)_j=\widetilde z_j-\log\sum_v e^{\widetilde z_v}$.
Все разности $z_j-\operatorname{LSE}(z)$ ниже вычисляют так,
не через логарифм округлённого softmax.
Для~\eqref{eq:modern-edge-attention} сохраняют
\begin{align}\label{eq:modern-edge-log-probabilities}
\ell_{ij}^{P}&=-D_{ij}^{G}/\tau_P
-\operatorname{LSE}_{v\in\mathcal N_i^P}(-D_{iv}^{G}/\tau_P),\\
z_{ij}^{A}&=\ell_{ij}^{P}+\eta_N\langle
\mathsf n_\eps(b_i^N),U_N\mathsf n_\eps(b_j^N)\rangle,\nonumber\\
\ell_{ij}^{A}&=z_{ij}^{A}
-\operatorname{LSE}_{v\in\mathcal N_i^P}(z_{iv}^{A}),\qquad
A_{ij}^N=\exp(\ell_{ij}^{A}).\nonumber
\end{align}
Округление экспоненты до нуля допустимо; её логарифм не восстанавливают
из этого нуля. Энтропия и KL ниже используют сохранённые логарифмы
и дифференцируются через тот же расчёт. Маскированные позиции исключают
до операций; пустая опора даёт отдельную нулевую ветвь с нулевым
градиентом. Умножение нуля на бесконечность и замена NaN нулём запрещены.

Представления четырёх ветвей имеют вид
\begin{align}\label{eq:modern-item-branches}
e_i^{\mathrm{id}}&=\mathsf n_\eps(b_i^N),&
e_i^{G}&=\mathsf n_\eps\!\left(
\sum_{j\in\mathcal N_i^P}A_{ij}^{N}B_Nb_j^N\right),\\
e_i^{T}&=\mathsf n_\eps(T_N\bar v_i^{\mathrm{text}}),&
e_i^{V}&=\mathsf n_\eps(V_N\bar v_i^{\mathrm{image}}).
\nonumber
\end{align}
Размеры матриц: $B_N\in\R^{d\times d}$,
$T_N\in\R^{d\times d_t}$, $V_N\in\R^{d\times d_v}$.
Для идентификатора маска равна $\ind\{i\in\mathcal I_G\}$;
для графа --- $\ind\{\mathcal N_i^P\ne\varnothing\}$;
для текста и изображения используются $M_i^{\mathrm{text}}$
и $M_i^{\mathrm{image}}$. Эти четыре маски обозначим $m_i^b$,
$b\in\{\mathrm{id},G,T,V\}$.
Недоступная ветвь даёт нулевой вектор и нулевую маску.

Пусть $\chi_i$ содержит маски модальностей и
$\log(1+V_i^{\mathrm{detail}})$ из $A_G$.
При наличии хотя бы одной ветви зададим
\begin{equation}\label{eq:modern-item-fusion}
\gamma_{ib}^{N}=
\frac{m_i^b\exp(c_b+v_b^\top\chi_i)}
{\sum_a m_i^a\exp(c_a+v_a^\top\chi_i)},\qquad
e_i=\mathsf n_\eps\!\left(\sum_b\gamma_{ib}^{N}e_i^b\right).
\end{equation}
Коэффициенты $c_b,v_b$ обучаются. Если все маски нулевые,
$e_i=0$; положим $m_i^N=\ind\{\norm{e_i}>\eps_N\}$.
Нулевые направления исключаются только из нейросетевого поиска.

Векторы Coveo фиксированы; обучаются проекции и перечисленные параметры.
Неизвестный идентификатор в префиксе может использовать предоставленные
векторы при допущении статических метаданных. Его ветви идентификатора
и графа отсутствуют, историческая частота равна нулю;
каталог кандидатов не расширяется.

Обучаемое представление поиска использует ту же текстовую проекцию:
\begin{equation}\label{eq:modern-query-vector}
z_q^N=\mathsf n_\eps(T_N\bar z_q),\qquad
m_q^{SN}=M_q^S\ind\{\norm{T_N\bar z_q}>\eps_N\}.
\end{equation}
Запрос и изображение сравниваются только в обученном общем пространстве,
не через скалярное произведение исходных векторов.

\subsection{Несколько намерений и выбор по поисковому контексту}
\label{sec:modern-intents}

Из $\mathcal D_q$ оставим события хотя бы с одной доступной входной
ветвью, то есть $\sum_bm_i^b>0$; получим $\mathcal D_q^N$.
Это условие зависит от данных, а не от обучаемой нормы $e_i$.
Пусть $n_q^N$ --- число различных
товаров в этом множестве, $J_I\ge1$ --- верхняя граница числа
намерений, а $J_q=\min(J_I,n_q^N)$.
Каждому $k=1,\ldots,J_I$ соответствует обучаемый вектор
$c_k^N\in\R^d$. Для события $\ell$ определим
\begin{align}\label{eq:modern-history-transform}
r_{q\ell}^{N}&=\tanh\!\left(
F_Ne_{i_{s(q)\ell}}
+a_N^t\log\!\left(1+\frac{t_q-t_{s(q)\ell}}{t_{\mathrm{unit}}}\right)
+a_N^f\log(1+n_{q,i_{s(q)\ell}}^{\mathrm{detail}})
\right),\\
z_{qk\ell}^{H,N}
&=-\frac{t_q-t_{s(q)\ell}}{\tau_H}+(c_k^N)^\top r_{q\ell}^{N},
\qquad k\le J_q,\nonumber\\
\beta_{qk\ell}^{N}
&=\exp\!\left(z_{qk\ell}^{H,N}
-\operatorname{LSE}_{v\in\mathcal D_q^N}(z_{qkv}^{H,N})\right),
\nonumber\\
h_{qk}&=\mathsf n_\eps\!\left(
\sum_{\ell\in\mathcal D_q^N}\beta_{qk\ell}^{N}e_{i_{s(q)\ell}}
\right).
\nonumber
\end{align}
Матрица $F_N$ имеет размер $d\times d$, а $a_N^t,a_N^f$ ---
векторы размерности $d$. Функция $\tanh$ применяется покоординатно.
Это нормирование временных весов~\eqref{eq:model-history-weight}
с поправкой на намерение, вычисленное без промежуточного округления
$\alpha_{q\ell}$: маска $\mathcal D_q^N$ применяется до устойчивого LSE.
Частоты и давности берутся из префикса; равные времена дают
равную давность, перестановки внутри группы не меняют суммы.

Динамический выбор намерений через entmax из полного MiaSRec
не воспроизводится. Верхняя граница $J_I$ ограничивает стоимость;
для единственного товара $J_q=1$. Направления --- скрытые векторы,
не установленные категории целей покупателя.

Обозначим через $\mathcal K_q=\{1,\ldots,J_q\}$ индексы направлений.
Если это множество пусто, но $M_q^S=1$, положим
$\mathcal K_q=\{1\}$ и $h_{q1}=z_q^N$.
Если отсутствует и входной поисковый вектор, множество остаётся пустым.
Нулевое направление сохраняет обучающий запрос с нулевыми баллами:
уменьшение норм не должно исключать примеры из потери.
При непустом $\mathcal K_q$ веса намерений равны
\begin{align}\label{eq:modern-intent-weights}
\delta_{qk}^{N}&=a_k^I+(v_k^I)^\top h_{qk}
+\beta_Sm_q^{SN}e^{-d_q^S/\tau_{SN}}
\langle z_q^N,h_{qk}\rangle,\\
\omega_{qk}&=\frac{\exp(\delta_{qk}^{N})}
{\sum_{v\in\mathcal K_q}\exp(\delta_{qv}^{N})},\qquad
u_q^I=\sum_{k\in\mathcal K_q}\omega_{qk}h_{qk}.
\nonumber
\end{align}
Параметры $a_k^I,v_k^I$ обучаются;
для энтропии сохраняют $\ell_{qk}^{I}=\delta_{qk}^N-
\operatorname{LSE}_{v\in\mathcal K_q}(\delta_{qv}^N)$ и
$\omega_{qk}=\exp(\ell_{qk}^{I})$.
$\beta_S=\log(1+e^{b_S})\ge0$ задаётся обучаемым $b_S$,
$\tau_{SN}>0$ измеряется в секундах.
При $m_q^{SN}=0$ всё поисковое слагаемое равно нулю;
$d_q^S$ полностью определена в~\eqref{eq:model-search-features}.
Вес $\omega_{qk}$ задаёт относительную значимость направления,
не вероятность покупки. При единственном направлении его вес всегда равен
единице, поэтому добавим прямую примесь поиска с ограниченным коэффициентом:
\begin{align}\label{eq:modern-search-gate}
\ell_q^g&=\begin{cases}
b_g+(v_g)^\top u_q^I+(w_g)^\top z_q^N+a_g^t a_q^S,
&m_q^{SN}=1,\\
0,&m_q^{SN}=0,
\end{cases}\\
g_q^{SN}&=m_q^{SN}\bar g_N\sigmoid(\ell_q^g),
\qquad 0\le\bar g_N\le1,\nonumber\\
u_q^N&=(1-g_q^{SN})u_q^I+g_q^{SN}z_q^N.
\nonumber
\end{align}
Здесь $\sigmoid(x)=(1+e^{-x})^{-1}$;
$b_g,a_g^t\in\R$ и $v_g,w_g\in\R^d$ входят в $\theta_N$.
При отсутствии корректного направления поиска $m_q^{SN}=0$
получаем $g_q^{SN}=0$, в том числе при $M_q^S=1$ и нулевой проекции.
Вырожденный поиск не уменьшает масштаб пригодного исторического вектора.
Предел $\bar g_N$ задаётся конфигурацией.
При пустом $\mathcal K_q$ определим $u_q^I=0$.
Допуск обучающего примера использует исходную маску $M_q^S$, а не
обучаемую $m_q^{SN}$; уменьшение нормы проекции не исключает пример
из потери. Порог направления задаёт кусочно-дифференцируемый расчёт:
маска при вычислении градиента фиксирована, в точке порога производная
не объявляется непрерывной.
При наличии только поиска $h_{q1}=z_q^N$ и формула сохраняет
$u_q^N=z_q^N$. Без поиска получаем $u_q^N=u_q^I$.
Выпуклая комбинация сохраняет $\norm{u_q^N}\le1$ и может изменить
направление даже при $\omega_{q1}=1$. При $\bar g_N=0$ поиск влияет
через веса намерений и базовый источник; для истории только с поиском
по-прежнему действует резервное правило $h_{q1}=z_q^N$.

Учебный пример: для ортогональных направлений с равными исходными
весами поиск, совпадающий с первым, при $\beta_S=\log 3$
и нулевой давности даёт веса $3/4$ и $1/4$.

\subsection{Отбор кандидатов и признаки для ранжирования}
\label{sec:modern-retrieval}

При $\tau_N>0$ нейросетевой балл задаётся формулой
\begin{align}\label{eq:modern-score}
s_{qi}^{N}
&=\frac{1-g_q^{SN}}{\tau_N}
\sum_{k\in\mathcal K_q}\omega_{qk}\langle h_{qk},e_i\rangle
+\frac{g_q^{SN}}{\tau_N}\langle z_q^N,e_i\rangle\\
&=\frac{\langle u_q^N,e_i\rangle}{\tau_N}.
\nonumber
\end{align}
Температура $\tau_N$ меняет масштаб потери, сохраняя порядок товаров.
Из ограничений норм следует $|s_{qi}^{N}|\le1/\tau_N$.

Основной новый источник возвращает первые $m_N^{\mathrm{mix}}$ товаров
по $s_{qi}^N$ среди $i\in\mathcal I_q$ с $m_i^N=1$.
Матрица товарных векторов фиксируется после обучения.
При практически нулевом $u_q^N$ источник возвращает пустой список
с маской, исключая произвольный выбор по одинаковым нулевым баллам.

Дополнительный источник выбирает до $m_{Nk}$ товаров
по $\langle h_{qk},e_i\rangle$ для каждого активного намерения.
Так сохраняются товары направления с малым $\omega_{qk}$.
Списки считаются отдельными источниками, но используют одну матрицу
$e_i$: для смеси $u_q^N$ и всех $h_{qk}$ достаточно разных запросов
к одному общему индексу. Отдельный индекс на намерение не требуется.
Общая нейросетевая квота $m_N$ и базовые квоты удовлетворяют
\begin{equation}\label{eq:modern-quota}
m_N^{\mathrm{mix}}+\sum_{k=1}^{J_I}m_{Nk}\le m_N,\qquad
m_N+\sum_{a\in\mathcal A_{\mathrm{base}}}m_a\le M.
\end{equation}
Неактивные направления и направления с нормой не больше $\eps_N$
имеют пустой список при отборе кандидатов.
Общий бюджет не увеличивается; удаление повторов, заполнение
популярностью и допустимость следуют разделу~\ref{sec:model-candidates}.

Для каждого кандидата объединения ранжировщик дополнительно получает
$s_{qi}^{N}$, $\max_{k\in\mathcal K_q}\langle h_{qk},e_i\rangle$,
сходство $\langle z_q^N,e_i\rangle$ с маской $m_q^{SN}m_i^N$,
коэффициент $g_q^{SN}$,
число направлений с нормой больше $\eps_N$ и следующие две характеристики:
\begin{align}\label{eq:modern-ranking-features}
H_q^I&=-\sum_{k\in\mathcal K_q}\exp(\ell_{qk}^{I})\ell_{qk}^{I},\\
V_{qi}^{I}&=\sum_{k\in\mathcal K_q}\omega_{qk}
\left(\frac{\langle h_{qk},e_i\rangle}{\tau_N}
-\frac{\langle u_q^I,e_i\rangle}{\tau_N}\right)^2.
\nonumber
\end{align}
Энтропия описывает распределённость весов, дисперсия --- различие
товарных оценок. Её центр $u_q^I$ не включает прямую примесь поиска,
которая не является намерением смеси.
При отсутствии пригодных векторов значения равны нулю с масками.
Отдельно сохраняются признаки происхождения и ранга нового источника.
Финальный $f(x_{qi})$ обучается заново на $B$ для каждой сравниваемой
архитектуры. Нейросетевой балл не называется вероятностью конверсии.

\subsection{Основная цель обучения: зарегистрированные будущие покупки}
\label{sec:modern-objective}

Пусть $\mathcal Q_N$ содержит первые запросы выбранных сессий
$\mathcal S_{A_N}^{\mathrm{obs}}$. Основная потеря использует только сессии
$\mathcal S_{A_N}^{+}$ с зарегистрированной будущей покупкой;
остальные доступны вспомогательным задачам ниже. Для запроса определим
$Y_q=Y_{s(q)}^{\mathrm{obs}}$, $Y_q^G=Y_q\cap\mathcal I_q$ и
$a_q^{\mathrm{buy}}=\ind\{|Y_q^G|>0\}$.
При $a_q^{\mathrm{buy}}=1$ целевое распределение равно
$t_{qi}^{\mathrm{buy}}=\ind\{i\in Y_q^G\}/|Y_q^G|$.
Различные товары зарегистрированных покупок имеют равные доли независимо
от неизвестных количества и цены.
Фиксированная маска обучающего входа равна
$m_q^N=\ind\{|\mathcal D_q^N|>0\ \text{или}\ M_q^S=1\}
\ind\{|\mathcal I_q|>0\}$.
Множество $\mathcal D_q^N$ и маска зависят от входных данных,
не от обучаемых норм.
При пустом $\mathcal K_q$ и отсутствии поиска $u_q^N=0$.
Для любого непустого $\mathcal I_q$ нормированная модель задаётся как
\begin{equation}\label{eq:modern-softmax}
p_N(i\given q)=\frac{\exp(s_{qi}^{N})}
{\sum_{j\in\mathcal I_q}\exp(s_{qj}^{N})}.
\end{equation}
Точный знаменатель включает весь непустой каталог запроса,
в том числе нулевые товарные представления. При $u_q^N=0$ распределение
равномерно, пример сохраняется. Запросы без доступного входа
исключаются до обучения; их доля регистрируется.

На каждую сессию приходится один запрос, поэтому $w_q^N=1$;
дополнительно положим $a_q^{\mathrm{buy}}=0$ при $m_q^N=0$.
Основная потеря равна
\begin{equation}\label{eq:modern-buy-loss}
\mathcal L_{\mathrm{buy}}^N=
-\frac{\displaystyle\sum_{q\in\mathcal Q_N}w_q^Na_q^{\mathrm{buy}}
\sum_{i\in\mathcal I_q}t_{qi}^{\mathrm{buy}}\log p_N(i\given q)}
{\displaystyle\sum_{q\in\mathcal Q_N}w_q^Na_q^{\mathrm{buy}}}.
\end{equation}
Слагаемое запроса с $a_q^{\mathrm{buy}}=0$ считается нулевым без
вычисления неопределённого $t_{qi}^{\mathrm{buy}}$.
Если знаменатель равен нулю, обучение покупкам невозможно;
вариант получает статус недостаточности данных.
Потеря учит выбирать товар зарегистрированной покупки, не оценивать
вероятность покупки или причинный эффект. Сессии без зарегистрированной
цели доступны вспомогательным задачам, но не входят в Hit@20 и не становятся
подтверждёнными отрицательными исходами.
Покупки вне $\mathcal I_G$ не расширяют каталог; их доля регистрируется,
и при оценке они остаются возможными промахами.

Для одного запроса производная по баллу имеет простой вид:
\begin{equation}\label{eq:modern-softmax-gradient}
\frac{\partial\left[-\sum_i t_{qi}\log p_N(i\given q)\right]}
{\partial s_{qj}^{N}}=p_N(j\given q)-t_{qj}.
\end{equation}
Производные по параметрам вычисляются по правилу цепочки.
Для устойчивости вычитается максимальный балл; знаменатель можно
считать блоками. Выборка товаров и смещение приближённой потери
заданы в приложении~\ref{sec:sampled-overview}.

\subsection{Вспомогательное предсказание следующего просмотра}
\label{sec:modern-next}

Пусть $t_q^P$ --- время первой зарегистрированной покупки после $t_q$
в той же сессии (при отсутствии записи $+\infty$).
Предсказывается следующий записанный просмотр, без утверждения
о полноте действий между записями. Множество допустимых времён равно
\begin{equation}\label{eq:modern-next-time}
\mathcal T_q^+=\{t_e:\ e\in s(q),\ \kappa_e=\mathrm{detail},\
t_q<t_e\le\tau_s,\ t_e<t_q^P\}.
\end{equation}
При $\mathcal T_q^+\ne\varnothing$ положим $t_q^+=\min\mathcal T_q^+$;
товары всей группы этого времени образуют $D_q^+$.
Последний просмотр и просмотр точно на горизонте допустимы,
одновременный с покупкой --- нет. При пустом $\mathcal T_q^+$
полагается $D_q^+=\varnothing$ без вычисления $t_q^+$.
После пересечения с $\mathcal I_q$ вводятся маска
$a_q^{\mathrm{next}}$ непустой цели при $m_q^N=1$ и равномерное распределение
$t_{qi}^{\mathrm{next}}$ на ней.
Потеря $\mathcal L_{\mathrm{next}}^N$ задаётся
формулой~\eqref{eq:modern-buy-loss} с заменой
$(a^{\mathrm{buy}},t^{\mathrm{buy}})$ на
$(a^{\mathrm{next}},t^{\mathrm{next}})$.
При пустом знаменателе эта вспомогательная потеря равна нулю
с признаком отсутствия обучающих примеров.
Вес $\lambda_{\mathrm{next}}\ge0$ сначала фиксирован,
затем выбирается на $C_{\mathrm{select}}$ по Hit@20, включая нулевой вариант.
Точность следующего просмотра не заменяет этот критерий.

\subsection{Контрастивное согласование двух представлений}
\label{sec:modern-contrastive}

Для одного доступного префикса строятся два представления
путём независимого отключения доступных ветвей
$\mathrm{id},G,T,V$ с заданными вероятностями.
Для каждого товара сохраняется хотя бы одна исходно доступная ветвь;
при полном отключении восстанавливается первая из них
в фиксированном порядке. Отключение меняет маски в
\eqref{eq:modern-item-fusion}, после чего повторяется расчёт сессии.
События, их времена и вектор поиска сохраняются.
Исходные тексты и изображения не требуются.

Обозначим результаты нормировки $\mathsf n_\eps$ через
$\tilde u_q^{(1)},\tilde u_q^{(2)}$.
Одинаковые маски не создают пару; нулевой обученный вектор
сохраняется в потере с нулевыми баллами.
В каждом пакете оставляется не более одного запроса одной сессии.
Сигнатура $\mathsf{sig}(q)$ содержит только входы кодировщика:
товары, относительные времена, частоты и предоставленные векторы.
В пакете запросов $\mathcal B_N$ зададим множество сравнений
\begin{equation}\label{eq:modern-contrastive-negative}
\mathcal V_q=\{v\in\mathcal B_N:
s(v)\ne s(q),\quad \mathsf{sig}(v)\ne\mathsf{sig}(q),\quad
Y_v^G\cap Y_q^G=\varnothing\}.
\end{equation}
В это множество включаются лишь запросы с различающимися масками
случайных преобразований; повторяющиеся сигнатуры среди сравнений
оставляют один раз по фиксированному порядку.
Метки $A_N$ используются только в потере, исключая совпадающие покупки
из противопоставляемых примеров. Разные покупки не доказывают
различия интересов.

При $\mathcal V_q\ne\varnothing$ и $\tau_C>0$ положим
\begin{equation}\label{eq:modern-contrastive-loss}
\ell_q^{1\to2}=-\log
\frac{\exp(\langle\tilde u_q^{(1)},\tilde u_q^{(2)}\rangle/\tau_C)}
{\exp(\langle\tilde u_q^{(1)},\tilde u_q^{(2)}\rangle/\tau_C)
+\sum_{v\in\mathcal V_q}
\exp(\langle\tilde u_q^{(1)},\tilde u_v^{(2)}\rangle/\tau_C)}.
\end{equation}
Потеря $\mathcal L_C^N$ --- взвешенное среднее
$(\ell_q^{1\to2}+\ell_q^{2\to1})/2$ по пригодным запросам
с весами $w_q^N$, нормированными на их сумму.
Без пригодных пар потеря нулевая с отметкой в отчёте.
При случайных пакетах минимизируется ожидание потери;
состав и размер пакета заданы конфигурацией.

\paragraph{Воспроизводимый выбор контрастивных сравнений.}
Правило \code{contrastive_sampler=label_filtered_v1} входит в паспорт
обучения; оно не является обучаемым параметром. Порядок пакетов,
случайных масок и нормировка потерь заданы в
разделе~\ref{sec:modern-training-recipe}. После формирования пакета
применяются только фильтры~\eqref{eq:modern-contrastive-negative};
сохраняются маски и числа исключённых сравнений по причинам.

В заранее зарегистрированном сравнении
\code{contrastive_sampler=label_free_v1} удаляется только условие
$Y_v^G\cap Y_q^G=\varnothing$, остальные правила и бюджет сохраняются.
Так измеряется влияние использования целей при выборе сравнений.
Запрос без зарегистрированной покупки может быть сравнением и без этого
фильтра; отсутствие записи не подтверждает различия интересов. Эти правила
применяются только к обучающему периоду $A_N$; кандидаты и входы
$B,C,D$ не используют будущие покупки при построении представления.

Предполагается устойчивость представления к частичной утрате модальностей.
Она может нарушаться при потере различающего признака, поэтому нужен
вариант без контрастивной части. Сессии в знаменателе не являются
подтверждёнными отрицательными показами; сохранность смысла при отключении
векторов и рост конверсии не гарантируются.

\subsection{Регуляризация и окончательная задача обучения}
\label{sec:modern-training}

Для ограничения отклонения графовых весов от исторической опоры
введём
\begin{equation}\label{eq:modern-graph-regularizer}
\mathcal L_G^N=
\frac1{|\mathcal I_G^P|}
\sum_{i\in\mathcal I_G^P}\sum_{j\in\mathcal N_i^P}
\exp(\ell_{ij}^{A})(\ell_{ij}^{A}-\ell_{ij}^{P}),\qquad
\mathcal I_G^P=\{i:\mathcal N_i^P\ne\varnothing\}.
\end{equation}
При пустом $\mathcal I_G^P$ положим $\mathcal L_G^N=0$.
При общей положительной опоре штраф неотрицателен и обнуляется
при совпадении строк. Он ограничивает изменение исторических связей,
не устанавливая причинности.

Итоговая функция имеет вид
\begin{equation}\label{eq:modern-total-loss}
\min_{\theta_N}\quad
\mathcal L_{\mathrm{buy}}^N
+\lambda_{\mathrm{next}}\mathcal L_{\mathrm{next}}^N
+\lambda_C\mathcal L_C^N
+\lambda_G^N\mathcal L_G^N
+\lambda_\theta\norm{\theta_N}^2,
\qquad\lambda_{\mathrm{next}},\lambda_C,\lambda_G^N,\lambda_\theta\ge0.
\end{equation}
Вес основной потери равен единице. Остальные веса не выражают
денежную стоимость событий: сначала они заданы приложением~\ref{app:parameters},
затем выбираются на $C_{\mathrm{select}}$ по плану сравнения. Нулевой вес отключает задачу.
Квадрат нормы $\theta_N$ --- сумма квадратов скалярных параметров;
размерности фиксируются до обучения.

\subsection{Воспроизводимый рецепт обучения расширений}
\label{sec:modern-training-recipe}

Профиль \code{neural_training_v2} применяется к включённым нейросетевому,
последовательному и SID-модулям. Эталонная среда: NumPy 2.1.3, PyTorch 2.6.0,
CPU, FP32 для параметров, потерь и состояния оптимизатора; один поток
вычислений и межоперационного параллелизма, детерминированные алгоритмы.
AMP, компиляция графа, накопление градиентов, расписание шага и ранняя
остановка выключены. Точные сборки, оборудование и хеши зависимостей
закрепляются в файле версий по разделу~\ref{app:run-manifest};
побитовое совпадение между разными средами не предполагается.

\paragraph{Начальные значения и случайность.}
Используют \code{Generator(PCG64(seed))}. Для модуля $m=0,1,2$
(основная нейросеть, последовательная модель, SID) раздельные
потоки имеют зерно $17+100m+j$: $j=0$ --- инициализация,
$j=1$ --- порядок примеров, $j=2$ --- маски ветвей,
$j=3$ --- шум SID, $j=4$ --- отрицательные товары,
$j=5$ --- отключение поиска, $j=6$ --- случайные контрольные SID-коды.
Отключённый потребитель не делает вызовов; состояния не сбрасываются
между эпохами. Порядок вызовов задан в соответствующих разделах.
Схема параметров содержит имена, размеры, роли и общие ссылки:
тензоры обходятся по UTF-8 именам, координаты --- построчно;
общая текстовая/поисковая матрица инициализируется один раз.
Случайные значения получают в FP64 вызовом \code{uniform} для целого
тензора, затем однократно приводят к FP32.

Матрица линейного преобразования размера $d_{\mathrm{out}}\times
d_{\mathrm{in}}$ получает независимые
$\mathrm{Uniform}[-a,a]$, $a=\sqrt{6/(d_{\mathrm{in}}+d_{\mathrm{out}})}$.
Векторы идентификаторов $b_i^N$, намерений $c_k^N$ и словарей
$c_{\ell b}^{\mathrm{SID}}$ получают
$\mathrm{Uniform}[-\sqrt{3/d_{\mathrm{emb}}},\sqrt{3/d_{\mathrm{emb}}}]$,
где $d_{\mathrm{emb}}$ --- их размерность;
строки упорядочены по ID товара или числовым индексам.
Позиции $p_b$ последовательной модели получают
$\mathrm{Uniform}[-0.02,0.02]$. Остальные векторы и скаляры равны нулю,
кроме множителей $\gamma$ нормировки LN, равных единице,
и $b_g$ из~\eqref{eq:app-search-init}.
Исключение для матриц: $U_N=0$, поэтому начальные $A_{ij}^N=\nu_{ij}^P$.
Эти же правила применяются к собственным параметрам последовательной
модели; отсутствующие в её архитектуре блоки не создаются.
Так задаются в том числе $b_S=0$, смещения декодеров и коэффициенты
времени/частоты; одинаковая нулевая инициализация намерений не используется.

\paragraph{Пакеты и масштаб потери.}
До начала обучения строят планы всех эпох. Для каждой эпохи
отсортированные по UTF-8 сессии переставляются вызовом
\code{permutation}; их первые запросы разбиваются на пакеты размера 64.
Все пригодные запросы встречаются ровно один раз за эпоху;
короткий пакет не отбрасывается. План пакетов и их число
$T_{\mathrm{ep}}$ сохраняются до обучения.
Для контрастивных масок обходят запросы в порядке пакета, представления
$1,2$, уникальные товары входного префикса по UTF-8 ID и ветви
$\mathrm{id},G,T,V$. Каждой ветви соответствует один вызов
\code{random()}, включая отсутствующую; отключение доступной ветви
происходит при значении меньше заданной вероятности. Восстановление
после полного отключения следует разделу~\ref{sec:modern-contrastive}.

Пусть $\ell_q^r$ --- слагаемое запроса в полной потере задачи
$r\in\{\mathrm{buy},\mathrm{next}\}$, уже включающее $w_q^N$,
маску и глобальный знаменатель~\eqref{eq:modern-buy-loss}.
На пакете $\mathcal B$ минимизируется
\begin{equation}\label{eq:modern-batch-loss}
\mathcal L_{\mathcal B}
=T_{\mathrm{ep}}\sum_{q\in\mathcal B}
(\ell_q^{\mathrm{buy}}+\lambda_{\mathrm{next}}\ell_q^{\mathrm{next}})
+\lambda_C\mathcal L_{C,\mathcal B}^N
+\lambda_G^N\mathcal L_G^N+\lambda_\theta\norm{\theta_N}^2.
\end{equation}
Глобальные знаменатели не пересчитываются по пакету.
При фиксированных параметрах среднее первых слагаемых по всем пакетам
равно полной покупочной и вспомогательной потере, включая короткие
пакеты. Контрастивная потеря нормируется внутри пакета по своему
определению; графовый штраф считается по всему фиксированному графу
на каждом шаге. Эта точная стоимость входит в бюджет обучения.

\paragraph{Обновление параметров.}
Используют \code{torch.optim.Adam} с единой группой всех незамороженных
параметров, нулевыми начальными моментами и настройками
\begin{verbatim}
lr=0.001; betas=(0.9,0.999); eps=1e-8; weight_decay=0
amsgrad=false; foreach=false; fused=false; maximize=false
capturable=false; differentiable=false
\end{verbatim}
На шаге вычисляют полный градиент $\mathcal L_{\mathcal B}$, включая
явный L2-штраф; градиенты плотные, отсутствие вклада даёт ноль.
Его общую евклидову норму вычисляют в FP64. При норме больше 5
весь градиент умножается на $5/\norm{g}$ до обновления моментов Adam;
при нулевой норме деление не выполняется. L2 не добавляется повторно
через оптимизатор. Неконечные потери, градиенты или параметры прерывают
обучение. После каждого пакета выполняется ровно один шаг Adam;
его моменты и счётчик продолжаются между эпохами.
Основная и последовательная модели обучаются три эпохи.
Контрольная точка хранит параметры, моменты, номер шага, план пакетов
и состояния всех генераторов.

\subsection{Проверка различимости намерений}
\label{sec:modern-intent-diagnostics}

Несколько направлений могут описывать один интерес. Их совпадение
не исключается функцией потерь. Пусть
$\mathcal K_q^+=\{k\in\mathcal K_q:\norm{h_{qk}}>\eps_N\}$,
$n_q^I=|\mathcal K_q^+|$ и
$\widehat h_{qk}=h_{qk}/\norm{h_{qk}}$ для этих направлений.
При $n_q^I\ge2$ и заранее заданном $0<\delta_I<1$ определим
\begin{align}\label{eq:modern-intent-duplicate}
\operatorname{Sim}_q^I
&=\frac{2}{n_q^I(n_q^I-1)}
\sum_{\substack{k<l\\k,l\in\mathcal K_q^+}}
\langle\widehat h_{qk},\widehat h_{ql}\rangle,\\
\operatorname{Dup}_q^I
&=\frac{2}{n_q^I(n_q^I-1)}
\sum_{\substack{k<l\\k,l\in\mathcal K_q^+}}
\ind\{\langle\widehat h_{qk},\widehat h_{ql}\rangle\ge1-\delta_I\}.
\nonumber
\end{align}
Первый показатель --- среднее сходство, второй --- доля почти совпавших
пар. При $n_q^I<2$ они не определены; публикуются доля таких запросов,
нормы и частота нулевых направлений, чтобы зануление не скрывало вырождение.

Пусть $\mathcal C_{qk}^I$ --- кандидаты отдельного намерения.
Для пар с непустым объединением измеряется
\begin{equation}\label{eq:modern-intent-overlap}
O_{qkl}^I=
\frac{|\mathcal C_{qk}^I\cap\mathcal C_{ql}^I|}
{|\mathcal C_{qk}^I\cup\mathcal C_{ql}^I|}.
\end{equation}
Для запроса берут среднее $O_{qkl}^I$ по допустимым парам $k<l$;
при отсутствии таких пар результат не определён.
Показатели усредняют по первым запросам, где они определены.
Доли пропусков считают по всем первым запросам; пропуски не заменяются нулями.
Отдельно сравниваются $\operatorname{Coverage}_{\mathrm{cand}}$
и окончательный Hit@20 для $J_I=1$ и $J_I=2$
при одинаковой общей квоте нейросетевого источника.
Ранжировщик каждого варианта обучается заново.
Высокая $H_q^I$ возможна и при одинаковых направлениях.

Разрезы --- длина истории и наличие поиска. Штраф
$\sum_{k\ne l}\langle c_k^N,c_l^N\rangle^2$ не входит в стартовую цель:
он минимален и при всех $c_k^N=0$, не гарантируя различия намерений.
Усложнение оправдывает измеренное совпадение направлений
без дополнительного охвата, а не само число просмотренных товаров.

\subsection{Устойчивость поискового сигнала}
\label{sec:modern-search-robustness}

Предел $\bar g_N$ ограничивает прямую примесь $z_q^N$.
Он не ограничивает всё влияние поиска: параметр $\beta_S$
в весах намерений остаётся обучаемым и не ограничен сверху.
Сохраняются распределения $g_q^{SN}$, $\max_k\omega_{qk}$,
$H_q^I$ и покупочного охвата по давности запроса.
Максимум берётся по $\mathcal K_q$; при пустом множестве
он не определён и сопровождается маской отсутствия.
Для сравнения запроса с историей используется вектор без поискового
смешивания
\begin{equation}\label{eq:modern-history-probe}
v_q^{\mathrm{hist}}=
\mathsf n_\eps\!\left(\sum_{\ell\in\mathcal D_q^N}
\alpha_{q\ell}e_{i_{s(q)\ell}}\right).
\end{equation}
Косинус $z_q^N$ и $v_q^{\mathrm{hist}}$ вычисляется только
при обеих нормах выше $\eps_N$. Расхождение может выражать смену интереса. Это историческая диагностика;
поиск после назначения варианта A/B не определяет подгруппу
подтверждающего причинного сравнения.

Если диагностика показывает неустойчивость к поиску, допускается
отдельный вариант случайного отключения этого сигнала при обучении.
Для каждого префикса и шага независимо выбирают
$D_q^S\sim\operatorname{Bernoulli}(1-p_S)$, $0\le p_S<1$, и используют
\begin{equation}\label{eq:modern-search-dropout}
\widetilde M_q^S=D_q^SM_q^S
\end{equation}
во всех поисковых ветвях нового кодировщика: проекции и её маске,
весах намерений, прямой примеси и резервном представлении только
по поиску. Отключённый вектор заменяется нулём.
Одно значение $D_q^S$ применяется к обоим контрастивным представлениям
данного префикса; базовые источники и целевые метки не меняются.
При $p_S>0$ поток $j=5$ вызывается один раз через \code{random()}
для каждого запроса в порядке пакета, включая запросы без поиска;
$D_q^S=\ind\{U\ge p_S\}$. При $p_S=0$ и при обслуживании вызовов нет.

Допуск $m_q^N$ и знаменатель обучающих весов вычисляются до случайного
отключения. Если после него не осталось входа, $u_q^N=0$ и пример
сохраняется с равномерным распределением вероятностей по товарам.
Минимизируется ожидание выбранной функции потерь по этой случайности.
В стартовой конфигурации $p_S=0$; ненулевое значение выбирается
на $C_{\mathrm{select}}$ только в отдельном сравнении. При обслуживании сигнал не отключается. Такая проверка устойчивости
не предписывает подавлять поиск при смене интереса и не гарантирует
разнообразия или роста охвата.

\subsection{Стоимость и поведение при отказе}
\label{sec:modern-cost}

Пусть $N_G=|\mathcal I_G|$, а
$d_0\le|\mathcal R_G|d_G$ --- максимальная исходящая степень
объединённого базового графа.
Память сохранённых путей составляет $O(N_Gd_P)$.
Ограниченный поиск раскрывает не более $N_GB_P$ состояний
и рассматривает не более $N_GB_Pd_0$ исходящих рёбер;
расход на очередь зависит от её реализации.
До усечения точный поиск в худшем случае существенно дороже.
Доля вершин с усечением публикуется вместе с качеством кандидатов.

Для одного полного пересчёта товарных представлений
проекции требуют порядка
$N_G[d(d_t+d_v)+d^2]$ операций,
а графовая агрегация после предварительного преобразования
векторов --- порядка $N_Gd_Pd$.
После фиксации модели результат хранится в матрице $N_G\times d$
и индексе; при обслуживании графовая нейросеть не пересчитывается.
Расчёт сессии требует порядка
$L_Hd^2+J_IL_Hd+d d_t$ операций,
если текущий поиск нужно проецировать заново.
Точный просмотр каталога требует $O(N_Gd)$;
приближённый поиск проверяется по тому же протоколу согласия
с точным поиском, что и базовые источники по содержанию.

При полном знаменателе стоимость нейросетевого обучения для
$N_Q^N$ запросов за эпоху включает $O(N_Q^NN_Gd)$;
блочное вычисление уменьшает память, но не число произведений.
Добавляются пересчёт графовых представлений и проекций на каждом
обновлении параметров и контрастивные пары. Сохранённые векторы
нельзя подставлять вместо представлений текущего кодировщика.
Пределы времени, памяти и числа эпох фиксируются до обучения;
превышение означает вычислительный отказ.

Отказы нового источника обрабатываются только по общим контрактам
\ref{sec:ops-partial-data} и~\ref{sec:ops-deadlines}. Незарегистрированная
деградация включает $\mathcal C_q^{\mathrm{fallback}}$ с её фиксированным
порядком, без ранжировщика основной модели.

\begin{invariant}[Фиксированное знание и доступный префикс]
При зафиксированных $\theta_N$ и артефактах $A_G$ изменение событий
после $t_q$ не меняет $h_{qk}$, $\omega_{qk}$, $g_q^{SN}$,
$u_q^I$, $u_q^N$ и $s_{qi}^N$.
Перестановка строк одной временной группы не меняет результат.
Отключение всех новых источников восстанавливает базовую политику
при восстановлении её исходных квот и ранжировщика.
\end{invariant}

\clearpage\part{Приложения: вычислительные варианты и дополнительные исследования}\label{part:appendices}
\appendix
\section{Приближённое обучение с выборкой товаров}
\label{sec:sampled-training}
\label{sec:sampled-overview}

Это необязательный вычислительный режим нейросетевого расширения;
базовая система его не использует. Эталон --- полный
знаменатель~\eqref{eq:modern-softmax}. Все положительные товары
учитываются точно, остальные --- случайной выборкой с весами
\cite{candidatesampling}. Смещение требует проверки
\cite{sampledsoftmax2025}; метод MIDX не воспроизводится.
Совпадение градиентов и качества с точным режимом не предполагается.
Метод NCE, решающий другую задачу различения наблюдений и шума, не используется.

\subsection{Положительные товары и распределение выборки}
\label{sec:sampled-proposal}

Рассмотрим один допустимый запрос $q$ из $A_N$ и одну задачу
$a\in\{\mathrm{buy},\mathrm{next}\}$ с непустой целью.
Обозначим её множество товаров через $T_q^a$:
$T_q^{\mathrm{buy}}=Y_q\cap\mathcal I_q$,
$T_q^{\mathrm{next}}=D_q^+\cap\mathcal I_q$.
Маски допустимости и веса запросов сохраняются
из разделов~\ref{sec:modern-objective}--\ref{sec:modern-next}.
На $T_q^a$ задано прежнее равномерное распределение $t_{qi}^a$.
Дополнение $\mathcal N_q^a=\mathcal I_q\setminus T_q^a$ содержит
отрицательные товары обучающей задачи; отсутствие цели не доказывает
просмотра и отказа от товара.

Пусть $N_G=|\mathcal I_G|>0$, а $\pi_j^{\mathrm{pop}}$ ---
распределение популярности из~\eqref{eq:model-popularity},
построенное только по $A_G$. При фиксированном $\delta\in(0,1]$
определим
\begin{equation}\label{eq:sampled-base-proposal}
Q_G(j)=(1-\delta)\pi_j^{\mathrm{pop}}+\frac{\delta}{N_G},
\qquad j\in\mathcal I_G.
\end{equation}
Равномерная примесь обеспечивает положительную вероятность
каждому товару каталога. Если $\mathcal N_q^a\ne\varnothing$,
условное распределение для данного запроса равно
\begin{equation}\label{eq:sampled-proposal}
Q_q^a(j)=
\frac{Q_G(j)}
{\sum_{v\in\mathcal N_q^a}Q_G(v)},
\qquad j\in\mathcal N_q^a.
\end{equation}
Разрешённую массу $b_q^a=\sum_{j\in\mathcal N_q^a}Q_G(j)$
в начале каждой включённой задачи с непустым дополнением вычисляют
прямой компенсированной суммой Кэхэна в FP64, в порядке UTF-8 ID;
вычитание массы исключённых товаров из единицы не используется.
В математической модели выборка из $Q_G$ с отбрасыванием запрещённых
товаров даёт $Q_q^a$ и требует в среднем $m/b_q^a$ попыток для $m$
товаров; малая масса снижает выгоду и регистрируется.
Распределение $Q_G$ независимо от обучаемых весов и фиксировано по $A_G$.
Метки исключают цели только при обучении на $A_N$.

В идеальной выборке независимо выбираются $m\ge1$ товаров
$J_1,\ldots,J_m$ с возвращением из $Q_q^a$.
Число попаданий товара $j$ обозначим
$C_{qj}^a=\sum_{b=1}^{m}\ind\{J_b=j\}$.
Повторы сохраняют кратность, иначе изменится ожидание суммы.
Выборка обновляется при новом предъявлении запроса;
используется поток $j=4$ профиля~\ref{sec:modern-training-recipe}.
В порядке пакета обрабатываются задачи покупки, затем следующего
просмотра каждого запроса. Кумулятивные суммы $Q_G$ в порядке UTF-8 ID
однократно строят в FP64, делят на последнюю сумму и сохраняют как $F$
с последним элементом 1.

До обучения в выборочном режиме проверяют конечность и положительность
вычисленных $Q_G(j)$ и $0=F_0<F_1<\cdots<F_{N_G}=1$.
При нумерации товаров в этом порядке сетке
$U\in\{0,1,\ldots,2^{53}-1\}/2^{53}$ вызова \code{random()}
в FP64 соответствуют массы
\[
\rho_j=2^{-53}\left(\left\lceil2^{53}F_j\right\rceil
-\left\lceil2^{53}F_{j-1}\right\rceil\right).
\]
Для каждого товара требуются $\rho_j>0$ и
$|\rho_j-Q_G(j)|\le10^{-6}Q_G(j)$; допуск является проектным,
а сравнение выполняется с вычисленными в FP64 $Q_G(j)$.
Нарушение этих условий или неконечная либо неположительная $b_q^a$
при непустом дополнении вызывает отказ
\code{negative_sampling_precision_unsupported} до расходования потока
выборки и обновления параметров. Товар не удаляют из опоры,
вероятности молча не обнуляют. Проверка не доказывает точную выборку
из вещественного $Q_G$: равенства ожиданий и предельные утверждения
ниже относятся к идеальной выборке, реализация является её численным
приближением с указанным допуском для базовых вероятностей.

На попытку вызывают один \code{random()} и
\code{searchsorted(F, U, side="right")};
запрещённый товар отбрасывают до получения $m$ допустимых результатов.
Вероятности рассчитываются в FP64. Выключенная задача, пустая цель,
пустое дополнение и точный режим не расходуют этот поток.
Выборку не сокращают при исчерпании бюджета обучения: это отказ запуска.

\subsection{Оценка знаменателя и производные}
\label{sec:sampled-loss}

Для краткости положим $s_i=s_{qi}^{N}$.
Точный знаменатель и его случайная оценка имеют вид
\begin{align}\label{eq:sampled-normalizer}
Z_q&=\sum_{i\in T_q^a}e^{s_i}
       +\sum_{j\in\mathcal N_q^a}e^{s_j},\\
\widehat Z_q^a&=\sum_{i\in T_q^a}e^{s_i}
       +\frac1m\sum_{b=1}^{m}
        \frac{e^{s_{J_b}}}{Q_q^a(J_b)}.
\nonumber
\end{align}
Каждая положительная цель входит один раз, даже если их больше $m$.
Если $\mathcal N_q^a=\varnothing$, выборка не выполняется:
$\widehat Z_q^a=\sum_{i\in T_q^a}e^{s_i}=Z_q$.
При пустом $T_q^a$ действует прежняя нулевая маска задачи;
положительная цель искусственно не создаётся.

При фиксированной модели получаем
\begin{equation}\label{eq:sampled-unbiased-normalizer}
\E\!\left[\frac{e^{s_{J_b}}}{Q_q^a(J_b)}\right]
=\sum_{j\in\mathcal N_q^a}e^{s_j},
\qquad
\E[\widehat Z_q^a]=Z_q.
\end{equation}
Приближённая потеря одного запроса равна
\begin{equation}\label{eq:sampled-query-loss}
\widehat\ell_q^a
=\log\widehat Z_q^a-\sum_{i\in T_q^a}t_{qi}^a s_i.
\end{equation}
Все положительные цели участвуют в усреднении.

Положим $\kappa_i=1$ для $i\in T_q^a$ и
$\kappa_j=C_{qj}^a/(mQ_q^a(j))$ для $j\in\mathcal N_q^a$.
Для фиксированной выборки получаем
\begin{equation}\label{eq:sampled-gradient}
\widehat p_i=\frac{\kappa_i e^{s_i}}{\widehat Z_q^a},
\qquad
\frac{\partial\widehat\ell_q^a}{\partial s_i}
=\widehat p_i-t_{qi}^a.
\end{equation}
Вне $T_q^a$ положим $t_{qi}^a=0$. У отрицательного товара вне выборки $\kappa_i=0$;
$\sum_i\widehat p_i=1$, но это распределение обучающего вычисления,
не полного каталога. Производные по параметрам следуют правилу цепочки
при фиксированных $Q_q^a$ и индексах выборки.

Для устойчивого вычисления $\log\widehat Z_q^a$ используются баллы
$s_i+\log\kappa_i$ при $\kappa_i>0$ с вычитанием максимума.
В развёрнутом виде поправка равна
\begin{equation}\label{eq:sampled-adjusted-score}
\widetilde s_i=
\begin{cases}
s_i, & i\in T_q^a,\\
s_i+\log C_{qi}^a-\log m-\log Q_q^a(i),
& i\in\mathcal N_q^a,\ C_{qi}^a>0.
\end{cases}
\end{equation}
Для этой схемы формула $s_i-\log Q(i)$ недостаточна: она не учитывает кратность,
размер выборки и точное включение положительных целей. Различие поправок
для обязательной цели и случайных товаров существенно~\cite{logq2025}.
Поправки применяются только к обучающему знаменателю;
при отборе кандидатов используется исходный $s_{qi}^N$.

Несмещённость~\eqref{eq:sampled-unbiased-normalizer}
не означает несмещённости потери. Из вогнутости логарифма следует
\begin{align}\label{eq:sampled-bias}
\E[\log\widehat Z_q^a]&\le\log Z_q,\\
\E[\nabla_\theta\widehat\ell_q^a]
&\ne\nabla_\theta\ell_q^a
\quad\text{в общем случае}.\nonumber
\end{align}
Здесь $\theta=\theta_N$, а
$\ell_q^a=\log Z_q-\sum_i t_{qi}^as_i$.
Точные и приближённые потери нельзя напрямую сравнивать
при выборе архитектуры.
При фиксированных параметрах, конечном каталоге и положительном
$Q_q^a$ закон больших чисел даёт
$\widehat Z_q^a\to Z_q$ при $m\to\infty$.
Потеря и её градиент в точках дифференцируемости также сходятся.
Сходимость при фиксированной модели не гарантирует совпадения
результатов двух процессов обучения.

\subsection{Включение в обучение и стоимость}
\label{sec:sampled-protocol}

Потеря $\widehat{\mathcal L}_a^N$ получается заменой
потери одного запроса в~\eqref{eq:modern-buy-loss}
на $\widehat\ell_q^a$ при прежних $w_q^N$ и $a_q^a$.
В~\eqref{eq:modern-total-loss} заменяются только
$\mathcal L_{\mathrm{buy}}^N$ и, при её использовании,
$\mathcal L_{\mathrm{next}}^N$.
Выборки для этих двух задач строятся по их собственным целям.
Контрастивная потеря, графовая регуляризация и штраф параметров неизменны.
Приближение не распространяется на отдельную полнокаталожную потерю SID.

Знаменатель требует не более $|T_q^a|+m$ скалярных произведений на запрос.
Прямой расчёт разрешённой массы дополнительно требует
$O(|\mathcal N_q^a|)$ скалярных операций на включённую задачу.
Дополнительно с текущими параметрами пересчитываются история,
положительные и выбранные товары и их графовые соседи; устаревшие
векторы не подставляются. Графовый регуляризатор по всем строкам
и остальные общие части сохраняют полную стоимость.

Для первого приближённого запуска принимаются $\delta=0.1$ и $m=512$;
$m\in\{1024,2048\}$ проверяются при доступном бюджете выбора на $C_{\mathrm{select}}$.
Любое $m$ означает выборку с возвращением; полный перебор --- точный режим.
Их сравнение на небольшом общем каталоге и одинаковых сессиях проверяет
вычисления и цену приближения, не гарантируя перенос на полный каталог.

Режимы следуют протоколу~\ref{sec:modern-time} при одинаковых каталоге,
квотах и метках проверки; Hit@20 и $\operatorname{Coverage}_{\mathrm{cand}}$
вычисляются без выборки отрицательных товаров.
Отчёт содержит время эпохи и всего обучения, число обновлений,
пиковую память, долю повторов в выборке и охват покупок.
Выигрыш во времени оценивается вместе с изменением охвата.

Трудные отрицательные товары допускаются только как отдельная гипотеза.
Для применения~\eqref{eq:sampled-normalizer} требуется известная фактическая
вероятность выборки каждого товара и положительная опора на всём
$\mathcal N_q^a$, например через равномерную примесь. При вычислении
градиента распределение фиксируется; его зависимость от параметров
и иной способ дифференцирования требуют отдельного вывода.
Произвольное добавление найденных индексом товаров с прежним $Q_q^a$
нарушает обоснование оценки знаменателя.

\subsection{Проверки вычислений}
\label{sec:sampled-checks}

Нормативные примеры, производные и граничные случаи сведены
в X02 (раздел~\ref{app:tests-extra}).

\section{Последовательная модель для контрольного сравнения}
\label{sec:sequence-control}
\label{sec:sequence-overview}

После нейросетевого этапа взвешенное объединение истории сравнивают
с одним блоком причинного внимания по мотивам SASRec~\cite{sasrec2018}.
Эта компактная адаптация заменяет кодировщик сессии, сохраняя данные,
покупочную цель, каталог и тип ранжировщика. Товарный граф остаётся
в обеих моделях; его вклад проверяют отдельным отключением графовой ветви.

\subsection{Одинаковые входы и группы одновременных событий}
\label{sec:sequence-input}

Используется история $\mathcal D_q^N$ доступных \code{detail}
из раздела~\ref{sec:modern-intents} с прежними $L_H$, обработкой
неизвестных идентификаторов и допущением о статических векторах.
Последовательности действий корзины, поиска и просмотра страниц
требуют отдельного сравнения: они меняют и архитектуру, и доступные входы.
Признаки корзины сохраняются у ранжировщика.

Товарный кодировщик задаётся формулами
\eqref{eq:modern-edge-attention}--\eqref{eq:modern-item-fusion},
но обучается заново на $A_N$ совместно с последовательным кодировщиком.
Обозначим его результат $e_i^{\mathrm{seq}}$.
Граф, пути и каталог из $A_G$ сохраняются; обученные параметры
основной модели не переносятся. Поисковая проекция $z_q^{\mathrm{seq}}$ получается по
\eqref{eq:modern-query-vector} с собственной обучаемой матрицей
$T_{\mathrm{seq}}$, общей для текстовой ветви и поиска.

Разобьём $\mathcal D_q^N$ на непустые временные группы
$\mathcal G_{q1},\ldots,\mathcal G_{qT_q}$ с временами
$t_{q1}<\cdots<t_{qT_q}\le t_q$, где $T_q\le L_H$.
Индекс $b$ обозначает группу, а не строку исходного файла.
Определим её входной вектор:
\begin{align}\label{eq:sequence-token}
n_{qb,i}^{\mathrm{detail}}
&=\sum_{\ell:\,e_{s(q)\ell}\in\mathcal H_q,\ t_{s(q)\ell}\le t_{qb}}
\ind\{i_{s(q)\ell}=i,\ \kappa_{s(q)\ell}=\mathrm{detail}\},\\
\bar e_{qb}&=\frac1{|\mathcal G_{qb}|}
\sum_{\ell\in\mathcal G_{qb}}e_{i_{s(q)\ell}}^{\mathrm{seq}},\nonumber\\
f_{qb}&=\frac1{|\mathcal G_{qb}|}
\sum_{\ell\in\mathcal G_{qb}}
\log(1+n_{qb,i_{s(q)\ell}}^{\mathrm{detail}}),\nonumber\\
x_{qb}&=F_{\mathrm{seq}}\bar e_{qb}+p_b
+a_{\mathrm{seq}}^t
\log\!\left(1+\frac{t_q-t_{qb}}{t_{\mathrm{unit}}}\right)
+a_{\mathrm{seq}}^f f_{qb}
+a_{\mathrm{seq}}^m\log|\mathcal G_{qb}|.\nonumber
\end{align}
Здесь $F_{\mathrm{seq}}\in\R^{d\times d}$;
$p_b,a_{\mathrm{seq}}^t,a_{\mathrm{seq}}^f,a_{\mathrm{seq}}^m\in\R^d$
обучаются на $A_N$, а $p_1,\ldots,p_{L_H}$ --- таблица позиций групп.
Частота учитывает события до конца группы, возраст --- время $t_q$.
Содержимое последующих групп недоступно. Порядок строк внутри группы
не влияет на вход; нулевые обученные векторы не удаляют события.

\subsection{Один блок причинного внимания}
\label{sec:sequence-attention}

Для $v\in\R^d$ определим нормировку координат
\begin{equation}\label{eq:sequence-layernorm}
\operatorname{LN}_{\gamma,\beta}(v)
=\gamma\odot
\frac{v-\bar v\boldsymbol 1}
{\sqrt{d^{-1}\sum_{a=1}^d(v_a-\bar v)^2+\eps_{\mathrm{LN}}}}
+\beta,\qquad
\bar v=d^{-1}\sum_{a=1}^dv_a.
\end{equation}
Знак $\odot$ означает покоординатное умножение;
$\gamma,\beta\in\R^d$ обучаются; начальное $\eps_{\mathrm{LN}}=10^{-5}$
входит в полную конфигурацию последовательного варианта.
Три применения $\operatorname{LN}_1,\operatorname{LN}_2,
\operatorname{LN}_3$ ниже имеют разные обучаемые параметры.

Для каждой существующей группы вычислим запрос, ключ и значение:
\begin{equation}\label{eq:sequence-qkv}
\tilde x_{qb}=\operatorname{LN}_1(x_{qb}),\qquad
Q_{qb}=W_Q\tilde x_{qb},\quad
K_{qb}=W_K\tilde x_{qb},\quad
V_{qb}=W_V\tilde x_{qb},
\end{equation}
где $W_Q,W_K,W_V\in\R^{d\times d}$.
Одна голова внимания использует причинную маску: группе доступны
она сама и предшествующие группы:
\begin{align}\label{eq:sequence-attention}
a_{qbc}
&=\frac{\exp(Q_{qb}^{\top}K_{qc}/\sqrt d)}
{\sum_{v=1}^{b}\exp(Q_{qb}^{\top}K_{qv}/\sqrt d)},
&&1\le c\le b\le T_q,\\
a_{qbc}&=0,&&c>b.\nonumber
\end{align}
Фиктивные группы при пакетном дополнении исключаются из вычислений.
Внимание к самой себе обеспечивает непустой знаменатель.
Далее применяются остаточные связи и покоординатная нелинейность:
\begin{align}\label{eq:sequence-ffn}
y_{qb}&=x_{qb}+W_O\sum_{c=1}^{b}a_{qbc}V_{qc},\\
h_{qb}^{\mathrm{seq}}
&=y_{qb}
+W_2\operatorname{ReLU}\!\left(
W_1\operatorname{LN}_2(y_{qb})+b_1\right)+b_2.\nonumber
\end{align}
Здесь $W_O\in\R^{d\times d}$,
$W_1\in\R^{d_f\times d}$, $W_2\in\R^{d\times d_f}$,
$b_1\in\R^{d_f}$, $b_2\in\R^d$ и
$\operatorname{ReLU}(v)_a=\max(0,v_a)$.
Стартовое значение $d_f=2d$ является вычислительным выбором.

Представление истории считывается в конце наблюдаемого префикса:
\begin{equation}\label{eq:sequence-readout}
u_q^{H,\mathrm{seq}}
=\mathsf n_\eps\!\left(
\operatorname{LN}_3(h_{qT_q}^{\mathrm{seq}})\right),
\qquad T_q>0.
\end{equation}
Если $T_q=0$ и $M_q^S=1$, положим
$u_q^{H,\mathrm{seq}}=z_q^{\mathrm{seq}}$;
при отсутствии и истории, и поиска положим его равным нулю.
Поиск добавляется тем же ограниченным смешиванием
\eqref{eq:modern-search-gate}, где $u_q^I$ заменено на
$u_q^{H,\mathrm{seq}}$, а $z_q^N$ --- на $z_q^{\mathrm{seq}}$.
Коэффициент смешивания обучается заново при прежних $\bar g_N$
и начальной настройке. Результат $u_q^{\mathrm{seq}}$ при наличии
только поиска равен $z_q^{\mathrm{seq}}$. Нулевые нормы не меняют
фиксированную маску обучающего примера.

\subsection{Обучение, кандидаты и условия сравнения}
\label{sec:sequence-training}

Для того же $i\in\mathcal I_q$ зададим
\begin{equation}\label{eq:sequence-score}
s_{qi}^{\mathrm{seq}}
=\frac{\langle u_q^{\mathrm{seq}},e_i^{\mathrm{seq}}\rangle}{\tau_N}.
\end{equation}
Покупочные цели, веса сессий и вспомогательные следующие просмотры
совпадают с разделами~\ref{sec:modern-objective}--\ref{sec:modern-next}.
Обучается сумма~\eqref{eq:modern-total-loss}
с заменой кодировщика и его параметров по
рецепту~\ref{sec:modern-training-recipe}.
Графовая регуляризация действует на собственный товарный кодировщик.
В первом сравнении обе архитектуры используют $\lambda_C=0$,
одинаковые способ нормирования, число отрицательных товаров и бюджет
обучения. При последующем включении контрастивной части совпадают
преобразования примеров и правила выбора пар. Число параметров
архитектур может различаться.

После обучения товарные векторы фиксируются и индексируются.
Один запрос $u_q^{\mathrm{seq}}$ к индексу возвращает до $m_N=90$
товаров по~\eqref{eq:sequence-score}
с нормой $e_i^{\mathrm{seq}}$, превышающей $\eps_N$.
Это вся нейросетевая квота; базовым источникам остаётся 210 товаров
при общем пределе $M=300$.
При норме запроса не больше $\eps_N$ нейросетевой источник пуст;
удаление повторов и заполнение остатка совпадают с основной политикой.
Допустимость кандидатов и запрет добавления целей следуют
\ref{sec:model-candidates}.

Ранжировщик получает балл~\eqref{eq:sequence-score}, поисковое
сходство с маской, коэффициент смешивания, происхождение и ранг кандидата.
Признаки нескольких намерений этой модели не приписываются.
Переобучение ранжировщика и выбор следуют~\ref{app:ablation};
архитектуры сравнивают при одинаковом критерии и заранее заданном бюджете.

\subsection{Стоимость и проверяемые свойства}
\label{sec:sequence-cost}

Группировка требует порядка $L_Hd$ операций над готовыми векторами.
Для одного префикса блок требует
$O(T_qd^2+T_q^2d+T_qdd_f)$ операций
и $O(T_q^2+T_qd+T_qd_f)$ промежуточной памяти.
При стартовом $d_f=2d$ последнее выражение равно $O(T_q^2+T_qd)$.
Отдельно учитываются проекция поиска, чтение товарных векторов,
поиск в индексе и окончательное ранжирование.
Товарный граф и его пути не пересчитываются при каждом запросе.
Стоимость обучения включает тот же пересчёт товарных представлений
и тот же выбранный способ вычисления знаменателя.
Стандартное сохранение ключей между запросами здесь не предполагается:
возраст события зависит от $t_q$, а ограничение истории меняет позиции.

Проверяются неизменность ответа при перестановке строк одной
временной группы и при изменении событий после $t_q$,
отсутствие внимания к дополнению пакета и сохранение общего бюджета $M$.
Перестановка товаров между различными временами может менять результат.
Публикуются охват покупок, время обучения, память и измеренная задержка.

\section{Старение связей и обновляемый граф}\label{sec:freshness}

Обновляемый граф --- отдельная политика, учитывающая изменение связей
между допустимыми товарами. Основное сравнение сохраняет снимок $A_G$
и временной протокол~\ref{sec:eval-temporal}.

\subsection{Два различных масштаба времени}\label{sec:freshness-weights}

Параметр $\tau_G$ в~\eqref{eq:model-session-edge} описывает расстояние
между действиями одной сессии. Старение записанных сессий зададим
окном $T_{\mathrm{win}}>0$ и масштабом $\tau_{\mathrm{age}}>0$ в секундах.

Пусть $c_k$ --- граница данных для $k$-го обновления,
а $r_k>c_k$ --- фактическое время публикации готового пакета.
Обозначим через $\mathcal S^{\mathrm{obs}}(c_k)$ выбранные исторические
сессии, для которых выполнены $t_{0,s}+H+g<c_k$ и
$t_{\mathrm{last},s}^{\mathrm{obs}}+g<c_k$.
Все опубликованные записи выбранной сессии прочитаны и доступны к $c_k$;
сессию с более поздним хвостом исключают целиком. Эти условия не подтверждают
полноту телеметрии. Правило одинаково для сессий с зарегистрированной покупкой
и без неё: известная цель не разрешает досрочное включение.

Оставим пары из~\eqref{eq:model-event-pairs}, у которых оба события
находятся в календарном окне, а оба товара принадлежат
$\mathcal I_G$. Для $i,j\in\mathcal I_G$ определим
\begin{equation}\label{eq:freshness-pairs}
\mathcal P^{r,\mathrm{fr}}_{sij}(c_k)
=\{(\ell,m)\in\mathcal P^r_{sij}(c_k):
c_k-T_{\mathrm{win}}\le t_{s\ell}<t_{sm}<c_k\}.
\end{equation}
Для допустимой пары введём
\begin{equation}\label{eq:freshness-pair-weight}
\xi_{s\ell m}(c_k)=
\exp\!\left(-\frac{t_{sm}-t_{s\ell}}{\tau_G}
            -\frac{c_k-t_{sm}}{\tau_{\mathrm{age}}}\right).
\end{equation}
Веса и поддержка равны
\begin{align}
w^{r,\mathrm{fr}}_{sij}(c_k)
&=
\begin{cases}
\displaystyle\max_{(\ell,m)\in\mathcal P^{r,\mathrm{fr}}_{sij}(c_k)}
\xi_{s\ell m}(c_k),
&\mathcal P^{r,\mathrm{fr}}_{sij}(c_k)\ne\varnothing,\\
0,&\text{иначе},
\end{cases}\label{eq:freshness-session}\\
W^{r,\mathrm{fr}}_{ij}(c_k)
&=\sum_{s\in\mathcal S^{\mathrm{obs}}(c_k)}
w^{r,\mathrm{fr}}_{sij}(c_k),\nonumber\\
n^{r,\mathrm{fr}}_{ij}(c_k)
&=\sum_{s\in\mathcal S^{\mathrm{obs}}(c_k)}
\ind\{w^{r,\mathrm{fr}}_{sij}(c_k)>0\}.
\label{eq:freshness-weight}
\end{align}
Сохраняется граница $0\le w^{r,\mathrm{fr}}_{sij}\le1$ и, следовательно,
$0\le W^{r,\mathrm{fr}}_{ij}\le n^{r,\mathrm{fr}}_{ij}$.
Поддержка считает различные сессии текущего окна, а не все сессии
за время существования графа. Разреживание, порог $n_G$,
нормировка и регуляризация строк применяются к этим новым величинам.
Равенство~\eqref{eq:model-row-regularization} остаётся справедливым.
Большая $\tau_{\mathrm{age}}$ ослабляет старение, но не отменяет
жёсткую границу окна.

\subsection{Публикация согласованной версии}\label{sec:freshness-publication}

Каталог $\mathcal I_G$ сохраняется. Новые товары требуют отдельного
правила допуска по времени первого наблюдения, инициализации параметров
и проверки совместимости. Статический файл метаданных не подтверждает
доступность товара в прошлом.

Вместе публикуются граф, ограниченные пути, товарные векторы,
индексы, признаки и совместимый ранжировщик.
Новый граф может менять векторы при прежних нейросетевых весах.
До проверки нового пакета используется предыдущий; запрос выбирает
версию по времени публикации $r_k<t_q$, а не по границе данных $c_k$.

В первом режиме параметры обученных моделей фиксированы.
Ранжировщик заранее обучен и проверен на ранних периодах
с воспроизведением принятого правила обновления графа.
Эта проверка проводится отдельно от статической модели.

Во втором режиме необходимые компоненты переобучаются на новых
исторических блоках с сохранением их ролей. Пакет публикуется после
всех этапов. Параметры, обучающие цели и результаты проверок доступны
до $r_k$; календарный допуск целей сохраняет запас $H+g$.

\subsection{Протокол сравнения и проверки}\label{sec:freshness-evaluation}

Обновляемую политику оценивают на отдельном последующем интервале
$D_{\mathrm{roll}}$. До его открытия фиксируют календарь обновлений,
$T_{\mathrm{win}},\tau_{\mathrm{age}}$, режим обучения, правило задержки
публикации, допустимые данные для каждого выпуска и резервный пакет.
Выбор этих величин выполняют на более ранних интервалах.
Период $D$ основного статического сравнения не используется
для их подбора или переобучения.
Обучение на ранних сессиях $D_{\mathrm{roll}}$ разрешено только после
выполнения календарного допуска; результаты не доступны предшествующим запросам.
Оценивается заранее заданная обучающаяся политика.

Статическую и обновляемую политики воспроизводят на одинаковых
запросах $D_{\mathrm{roll}}$ с одинаковыми $K,M,H,Q_{\max}$.
Для отделения эффекта затухания от пользы новой истории дополнительно
сравнивают веса на одном окне данных.
Публикуют Hit@20, $\operatorname{Coverage}_{\mathrm{cand}}$, возраст данных $t_q-c_k$,
задержку публикации $r_k-c_k$, стоимость пересчёта и число пропущенных
обновлений. В A/B-тесте обучение использует заранее определённый общий внешний поток
либо раздельные данные групп. Их смешивание требует анализа
взаимодействия групп.

Время, веса и неизменность опубликованного графа проверяет X01
(раздел~\ref{app:tests-extra}).

\section{Генератор кандидатов по семантическим кодам}
\label{sec:sid-overview}

Семантический идентификатор (SID) кодирует товар последовательностью
номеров небольших словарей. Дополнительный источник после нейросетевого
этапа предсказывает эти номера по очереди.
Заимствованные принципы DIGER и GRID указаны в~\ref{sec:research}.

Эксперимент использует векторы Coveo, позиционное мягкое квантование
и рекуррентный декодер: следующий номер зависит от предыдущих.
От DIGER он отличается отсутствием языкового кодировщика, остаточного
квантования и адаптивного изменения шума. Назначения кодов обучаются
по отдельной полнокаталожной рекомендательной потере; графовый
кодировщик, ранжировщик и бизнес-критерий сохраняются.

\subsection{Обучающие данные и фиксированные представления}
\label{sec:sid-data}

Каталог $\mathcal I_G$ и частоты фиксируются по $A_G$.
Обученный на $A_N$ нейросетевой модуль замораживается;
его $e_i\in\R^d$ и $u_q^N\in\R^d$ становятся входами SID.
SID также обучается на $A_N$, который не служит проверкой обобщения.
Размерности, число эпох и коэффициенты выбираются на $C_{\mathrm{select}}$,
без подгонки обучаемых параметров по этому периоду.

Пусть $x_i^{\mathrm{SID}}$ объединяет нормированные векторы описания
и изображения с масками наличия; отсутствующая модальность даёт
нулевой блок. Категории и исходные изображения не требуются.
Фиксированный вход равен
\begin{equation}\label{eq:sid-fixed-input}
 z_i^{\mathrm{SID}}=[e_i;x_i^{\mathrm{SID}}].
\end{equation}
Объединение $[\,;\,]$ использует статические данные
с ограничением раздела~\ref{sec:data}.

Для запроса $q$ обучающего блока положим
$Y_q^G=Y_q\cap\mathcal I_q\cap\mathcal I_G$.
В множество $\mathcal Q_N^{\mathrm{SID}}$ входят первые запросы сессий
$\mathcal S_{A_N}^{+}$ с $Y_q^G\ne\varnothing$. Положим
$w_q^{\mathrm{SID}}=1$ и
$W_N^{\mathrm{SID}}=\sum_{q\in\mathcal Q_N^{\mathrm{SID}}}w_q^{\mathrm{SID}}$.
Суммарный вес каждой сессии равен единице. При пустом каталоге
или $W_N^{\mathrm{SID}}=0$ обучение пропускается со статусом
\code{insufficient}. Запросы без зарегистрированных покупок не создают
обучающих целей и не входят в Hit@20. Покупки вне $\mathcal I_G$
не расширяют каталог и остаются возможными промахами при оценке.

\subsection{Мягкие назначения и жёсткий код}
\label{sec:sid-codes}

Пусть $L_{\mathrm{SID}}$ --- длина кода,
$V_{\mathrm{SID}}\ge2$ --- размер каждого словаря,
$d_{\mathrm{SID}}$ --- размерность его векторов.
Для позиции $\ell$ обучаются матрица $A_\ell^{\mathrm{SID}}$,
смещение $b_\ell^{\mathrm{SID}}$ и векторы
$c_{\ell b}^{\mathrm{SID}}\in\R^{d_{\mathrm{SID}}}$,
$b=1,\ldots,V_{\mathrm{SID}}$. Логиты назначения равны
\begin{equation}\label{eq:sid-logits}
 a_{i\ell b}^{\mathrm{SID}}
 =-\norm{A_\ell^{\mathrm{SID}}z_i^{\mathrm{SID}}
       +b_\ell^{\mathrm{SID}}-c_{\ell b}^{\mathrm{SID}}}^{2}.
\end{equation}
Равные максимумы разрешаются меньшим номером элемента словаря.
На шаге обучения $t$ выбираются независимые
$U_{i\ell b}\sim\mathrm{Uniform}(0,1)$ и
$G_{i\ell b}=-\log(-\log U_{i\ell b})$.
Для численной устойчивости $U$ ограничивается интервалом
$[\eps_U,1-\eps_U]$, где $0<\eps_U<1/2$.
Мягкие назначения задаются формулой
\begin{equation}\label{eq:sid-soft}
 \rho_{i\ell b}^{\mathrm{SID}}(t)=
 \frac{\exp\{(a_{i\ell b}^{\mathrm{SID}}
                 +\sigma_t^{\mathrm{SID}}G_{i\ell b})/
                    \tau_t^{\mathrm{SID}}\}}
 {\sum_{v=1}^{V_{\mathrm{SID}}}
  \exp\{(a_{i\ell v}^{\mathrm{SID}}
                 +\sigma_t^{\mathrm{SID}}G_{i\ell v})/
                    \tau_t^{\mathrm{SID}}\}}.
\end{equation}
Здесь $\tau_t^{\mathrm{SID}}>0$ задаёт мягкость распределения,
а $\sigma_t^{\mathrm{SID}}\ge0$ --- масштаб шума.
Все мягкие потери шага используют один шум для товара.
Шум воспроизводится потоком $j=3$ профиля
\ref{sec:modern-training-recipe}: на совместном шаге обходят весь
каталог по UTF-8 ID, затем $\ell,b$, используя \code{random()} и
$\eps_U=10^{-7}$. Логарифмы шума вычисляют в FP64, результат приводят
к FP32. Расписание~\eqref{eq:app-sid-schedule} использует
$t=1,\ldots,T_{\mathrm{SID}}$, где $T_{\mathrm{SID}}$ --- число шагов
трёх совместных эпох, вычисленное по сохранённому плану пакетов.
На предварительном этапе $\sigma=0,\tau=1$; шум не генерируется.

Позиционные векторы и их объединение равны
\begin{equation}\label{eq:sid-soft-vector}
 \widetilde c_{i\ell}^{\mathrm{SID}}
 =\sum_b\rho_{i\ell b}^{\mathrm{SID}}c_{\ell b}^{\mathrm{SID}},
 \qquad
 \widetilde v_i^{\mathrm{SID}}
 =[\widetilde c_{i1}^{\mathrm{SID}};\ldots;
   \widetilde c_{iL_{\mathrm{SID}}}^{\mathrm{SID}}].
\end{equation}
Конкатенация сохраняет различие позиций. В рабочем индексе
используются только детерминированные коды
\begin{equation}\label{eq:sid-hard}
 k_{i\ell}^{\mathrm{SID}}=\argmax_b a_{i\ell b}^{\mathrm{SID}},
 \qquad
 \mathbf k_i^{\mathrm{SID}}
 =(k_{i1}^{\mathrm{SID}},\ldots,k_{iL_{\mathrm{SID}}}^{\mathrm{SID}}).
\end{equation}
Шум в обслуживании запросов и исторической проверке выключен.

\subsection{Потери и прохождение градиента к назначениям}
\label{sec:sid-loss}

Обучаемая матрица $R^{\mathrm{SID}}$ переводит сессионный вектор
в пространство объединённых кодов. При $\tau_R^{\mathrm{SID}}>0$
зададим полнокаталожное распределение
\begin{align}\label{eq:sid-full-prob}
 a_{qi}^{R,\mathrm{SID}}
 &=\frac{(R^{\mathrm{SID}}u_q^N)^\top
                  \widetilde v_i^{\mathrm{SID}}}
                {\tau_R^{\mathrm{SID}}},\\
 p_{qi}^{R,\mathrm{SID}}
 &=\frac{\exp(a_{qi}^{R,\mathrm{SID}})}
 {\sum_{j\in\mathcal I_q\cap\mathcal I_G}
             \exp(a_{qj}^{R,\mathrm{SID}})}.
\end{align}
Потеря будущей покупки имеет вид
\begin{equation}\label{eq:sid-rec-loss}
 \mathcal L_R^{\mathrm{SID}}
 =-\frac{1}{W_N^{\mathrm{SID}}}
 \sum_{q\in\mathcal Q_N^{\mathrm{SID}}}
 \frac{w_q^{\mathrm{SID}}}{|Y_q^G|}
       \sum_{i\in Y_q^G}\log p_{qi}^{R,\mathrm{SID}}.
\end{equation}
Остальные товары образуют нормирующий каталог;
отсутствие записи о покупке не означает отрицательной реакции на показ.

Для одного запроса и распределения целей
$\pi_{qi}^{\mathrm{SID}}=\ind\{i\in Y_q^G\}/|Y_q^G|$
производная по логиту равна
$p_{qi}^{R,\mathrm{SID}}-\pi_{qi}^{\mathrm{SID}}$.
Поскольку
\begin{equation}\label{eq:sid-soft-gradient}
 \frac{\partial\rho_{i\ell b}^{\mathrm{SID}}}
      {\partial a_{i\ell v}^{\mathrm{SID}}}
 =\frac{\rho_{i\ell b}^{\mathrm{SID}}}
            {\tau_t^{\mathrm{SID}}}
       \bigl(\ind\{b=v\}-\rho_{i\ell v}^{\mathrm{SID}}\bigr),
\end{equation}
существует цепочка производных
$\mathcal L_R^{\mathrm{SID}}\to\widetilde v_i^{\mathrm{SID}}
\to\rho_{i\ell b}^{\mathrm{SID}}\to a_{i\ell b}^{\mathrm{SID}}
\to A_\ell^{\mathrm{SID}},c_{\ell b}^{\mathrm{SID}}$.
Эта производная обучает мягкие назначения, но не дифференцирует
дискретный $\argmax$. Улучшение мягкой потери не гарантирует
улучшения жёсткого поиска.

Для сохранения информации модальностей используются линейные
декодеры $D_m^{\mathrm{SID}}$, $m\in\{\mathrm{text},\mathrm{image}\}$.
Пусть $N_m=\sum_{i\in\mathcal I_G}M_i^m$. Определим
\begin{equation}\label{eq:sid-reconstruction}
 \mathcal L_X^{\mathrm{SID}}
 =\sum_{m:N_m>0}\frac{1}{N_m}
 \sum_{i\in\mathcal I_G}M_i^m
 \norm{D_m^{\mathrm{SID}}\widetilde v_i^{\mathrm{SID}}
                        -\bar v_i^m}^{2}.
\end{equation}
При полностью отсутствующей модальности слагаемое равно нулю.
Восстанавливаются предоставленные векторы, поскольку исходные тексты
и изображения недоступны.

Средняя загрузка кода и балансирующая потеря определяются как
\begin{equation}\label{eq:sid-balance}
 \bar\rho_{\ell b}^{\mathrm{SID}}
 =\frac{1}{|\mathcal I_G|}\sum_{i\in\mathcal I_G}
          \rho_{i\ell b}^{\mathrm{SID}},
 \qquad
 \mathcal L_B^{\mathrm{SID}}
 =\frac{1}{L_{\mathrm{SID}}}\sum_{\ell,b}
  \bar\rho_{\ell b}^{\mathrm{SID}}
            \log(V_{\mathrm{SID}}\bar\rho_{\ell b}^{\mathrm{SID}}).
\end{equation}
Для мягких назначений при расчёте с градиентами сохраняют
$r_{\ell b}=\operatorname{LSE}_{i\in\mathcal I_G}
(\log\rho_{i\ell b}^{\mathrm{SID}})-\log|\mathcal I_G|$,
где $\log\rho$ получают непосредственно через \code{logsoftmax}.
Слагаемое балансировки считают как $\exp(r_{\ell b})
(\log V_{\mathrm{SID}}+r_{\ell b})$ по~\ref{sec:modern-log-probabilities}.
В контрольных вариантах с фиксированными кодами баланс считают по
ненулевым целочисленным частотам; нулевые ячейки пропускают,
градиент балансирующей потери равен нулю.
Регуляризатор не гарантирует разнообразия жёстких кодов: равномерные
мягкие назначения дают нулевую потерю, но после разрешения равенств
весь каталог может получить один код. Его загрузка и коллизии
проверяются по пункту~\ref{sec:sid-validation}; чрезмерная регуляризация
может разрушить полезные группы товаров.

\subsection{Авторегрессивный декодер и порядок обучения}
\label{sec:sid-decoder}

Состояние декодера $h_{q\ell}^{\mathrm{SID}}\in\R^{d_h^{\mathrm{SID}}}$
инициализируется вектором сессии. Для частично сформированного кода
$k_1,\ldots,k_{\ell-1}$ задаются
\begin{align}\label{eq:sid-decoder}
 h_{q1}^{\mathrm{SID}}&=\tanh(U^{\mathrm{SID}}u_q^N+b_0^{\mathrm{SID}}),\\
 P_\ell^{\mathrm{SID}}(b\given q,k_{<\ell})
 &=\frac{\exp\{(O_\ell^{\mathrm{SID}}h_{q\ell}^{\mathrm{SID}}
                         +o_\ell^{\mathrm{SID}})_b\}}
 {\sum_v\exp\{(O_\ell^{\mathrm{SID}}h_{q\ell}^{\mathrm{SID}}
                         +o_\ell^{\mathrm{SID}})_v\}},\\
 h_{q,\ell+1}^{\mathrm{SID}}
 &=\tanh(F^{\mathrm{SID}}h_{q\ell}^{\mathrm{SID}}
       +T^{\mathrm{SID}}c_{\ell k_\ell}^{\mathrm{SID}}
       +p_\ell^{\mathrm{SID}}).
\end{align}
Функция $\tanh$ применяется к каждой координате.
Все матрицы и смещения в этих формулах обучаются.
Вероятность полной последовательности --- произведение
условных вероятностей её позиций.

В начале эпохи для всех товаров вычисляются жёсткие коды
\eqref{eq:sid-hard}, которые фиксируются как цели этой эпохи.
При обучении следующей позиции в декодер подаётся правильный
предыдущий код. Потеря равна
\begin{equation}\label{eq:sid-gen-loss}
 \mathcal L_G^{\mathrm{SID}}
 =-\frac{1}{W_N^{\mathrm{SID}}L_{\mathrm{SID}}}
 \sum_{q\in\mathcal Q_N^{\mathrm{SID}}}
 \frac{w_q^{\mathrm{SID}}}{|Y_q^G|}
 \sum_{i\in Y_q^G}\sum_{\ell=1}^{L_{\mathrm{SID}}}
 \log P_\ell^{\mathrm{SID}}
       (k_{i\ell}^{\mathrm{SID}}\given q,k_{i,<\ell}^{\mathrm{SID}}).
\end{equation}
Потеря~\eqref{eq:sid-gen-loss} не дифференцирует индексы целей;
связь с $A_\ell^{\mathrm{SID}}$ и $b_\ell^{\mathrm{SID}}$ обеспечивает
\eqref{eq:sid-rec-loss}. Градиент через входные векторы декодера
изменяет таблицы $c_{\ell b}^{\mathrm{SID}}$ и возможные будущие назначения,
но не является производной по номеру кода.

Совместно минимизируется
\begin{align}\label{eq:sid-objective}
 \mathcal L^{\mathrm{SID}}
 &=\mathcal L_G^{\mathrm{SID}}
  +\lambda_R^{\mathrm{SID}}\mathcal L_R^{\mathrm{SID}}
  +\lambda_X^{\mathrm{SID}}\mathcal L_X^{\mathrm{SID}}
  \nonumber\\
 &\quad+\lambda_B^{\mathrm{SID}}\mathcal L_B^{\mathrm{SID}}
  +\lambda_2^{\mathrm{SID}}\norm{\theta_{\mathrm{SID}}}^{2},
\end{align}
где $\lambda_R^{\mathrm{SID}}>0$, остальные коэффициенты неотрицательны,
а $\theta_{\mathrm{SID}}$ включает только параметры этого модуля.
Профиль~\ref{sec:modern-training-recipe} применяется во всех трёх фазах.
В предварительной фазе обучаются только $A_\ell^{\mathrm{SID}}$,
$b_\ell^{\mathrm{SID}}$, $c_{\ell b}^{\mathrm{SID}}$ и $D_m^{\mathrm{SID}}$;
остальные параметры SID заморожены. Выполняется одна эпоха по
$\mathcal L_X^{\mathrm{SID}}+\lambda_B^{\mathrm{SID}}\mathcal L_B^{\mathrm{SID}}$:
каталог, исходно отсортированный по UTF-8 ID, переставляется вызовом
\code{permutation} потока порядка; пакеты содержат до 256 товаров.
Сумма товарных слагаемых восстановления с глобальными $N_m$
умножается на число пакетов эпохи; баланс считается по всему каталогу.
Затем все параметры $\theta_{\mathrm{SID}}$ обучаются три совместные
эпохи по~\eqref{eq:sid-objective},
с пакетами до 64 запросов разных сессий. Покупочные и генеративные
слагаемые сохраняют глобальные веса и множитель $T_{\mathrm{ep}}$,
как в~\eqref{eq:modern-batch-loss}; восстановление, баланс и L2
вычисляются полностью на каждом шаге.
В конце фиксируются кодировщик и таблицы кодов, пересчитываются
все жёсткие коды, затем только декодер дополнительно обучается
по \eqref{eq:sid-gen-loss} на том же $A_N$ с окончательными целями.
В последней фазе обучаются только $U^{\mathrm{SID}},b_0^{\mathrm{SID}}$,
$O_\ell^{\mathrm{SID}},o_\ell^{\mathrm{SID}},F^{\mathrm{SID}}$,
$T^{\mathrm{SID}},p_\ell^{\mathrm{SID}}$ из~\eqref{eq:sid-decoder}.
Она длится две эпохи с тем же правилом пакетов
и глобальной нормировкой; остальные потери выключены.
В начале каждой фазы моменты и счётчик Adam обнуляются,
состояния генераторов продолжаются. Списки обучаемых и замороженных
тензоров каждой фазы сохраняются в паспорте; исходная нейросеть
остаётся замороженной во всех фазах SID.
После обучения декодера весь модуль замораживается до обработки $B$.

\subsection{Допустимые последовательности, коллизии и размер выдачи}
\label{sec:sid-serving}

Разные товары могут получить одинаковый код.
Для последовательности $\mathbf k$ определим корзину соответствий
\begin{equation}\label{eq:sid-bucket}
 \mathcal B^{\mathrm{SID}}(\mathbf k)
 =\{i\in\mathcal I_G:\mathbf k_i^{\mathrm{SID}}=\mathbf k\}.
\end{equation}
В таблице соответствий товары сортируются по популярности из $A_G$,
затем по идентификатору. Обслуживаемая часть
$\mathcal B_R^{\mathrm{SID}}(\mathbf k)$ содержит первые $R_{\mathrm{SID}}$
товаров; доля исключённых сообщается отдельно. Товар однозначно
определяет пара <<код, позиция в таблице>>, поскольку код может повторяться.

Префиксное дерево содержит коды каталога: ребро задаёт следующий
номер, путь к листу --- существующий код. Декодер использует только
рёбра этого дерева.
На глубине $\ell$ сохраняются не более $B_{\mathrm{SID}}$ путей
с наибольшими суммами исходных логарифмов вероятностей
из \eqref{eq:sid-decoder}. После запрета ребра вероятности
не перенормируются: иначе изменились бы сравниваемые оценки путей.
Равенство разрешается лексикографически. Ограниченная ширина поиска
не гарантирует глобально лучших $B_{\mathrm{SID}}$ последовательностей.

Пусть $\mathcal K_q^{\mathrm{SID}}$ --- найденные листья.
Из множества
\begin{equation}\label{eq:sid-candidates}
 \mathcal U_q^{\mathrm{SID}}
 =\mathcal I_q\cap\bigcup_{\mathbf k\in\mathcal K_q^{\mathrm{SID}}}
                     \mathcal B_R^{\mathrm{SID}}(\mathbf k)
\end{equation}
отбираются первые $m_{\mathrm{SID}}$ различных товаров по оценке
$\sum_\ell\log P_\ell^{\mathrm{SID}}
(k_{i\ell}^{\mathrm{SID}}\given q,k_{i,<\ell}^{\mathrm{SID}})$,
затем по той же исторической популярности и идентификатору.
Следовательно, $|\mathcal C_q^{(\mathrm{SID})}|
\le\min(m_{\mathrm{SID}},B_{\mathrm{SID}}R_{\mathrm{SID}})$.
Квота $m_{\mathrm{SID}}$ входит в прежний предел $M$;
обычные источники и общее правило дополнения сохраняются.

Дополнительными признаками могут быть ранговые позиции источника,
логарифм вероятности кода и размер таблицы соответствий. Вероятность
кода не оценивает покупку или причинный эффект показа. Переобучение
ранжировщика следует~\ref{app:ablation}, ограничения кандидатов ---
\ref{sec:model-candidates}, обработка отказов --- \ref{sec:modern-cost}.

Генератор включается в обслуживание только после измеренной проверки
задержки, памяти и стоимости на целевом узле. При нарушении требований
источник исключается; исследовательские результаты сохраняются отдельно.

\subsection{Диагностика, сравнения и вычислительные ограничения}
\label{sec:sid-validation}

После выключения шума измеряются доли
$f_{\ell b}^{\mathrm{SID}}=|\mathcal I_G|^{-1}
\sum_i\ind\{k_{i\ell}^{\mathrm{SID}}=b\}$,
эффективное число кодов
$\exp\{-\sum_b f_{\ell b}^{\mathrm{SID}}\log f_{\ell b}^{\mathrm{SID}}\}$,
доля различных полных кодов и доля товаров за пределами
$R_{\mathrm{SID}}$. Нулевые слагаемые энтропии равны нулю.
Сопоставляются мягкая потеря и качество жёсткого источника.
До запуска фиксируются пороги загрузки, коллизий и времени работы;
регуляризатор не гарантирует их соблюдения.

Диагностика выполняется по всему $\mathcal I_G$ до ограничения
$R_{\mathrm{SID}}$. Пусть $N_G=|\mathcal I_G|>0$,
$b_{\mathbf k}=|\mathcal B^{\mathrm{SID}}(\mathbf k)|$,
а сумма ниже берётся по всем непустым таблицам соответствий:
\begin{align}
\operatorname{Col}_{\mathrm{item}}^{\mathrm{SID}}
&=\frac1{N_G}\sum_{\mathbf k}
b_{\mathbf k}\ind\{b_{\mathbf k}>1\},\label{eq:sid-catalog-collisions}\\
\operatorname{Drop}_{R}^{\mathrm{SID}}
&=\frac1{N_G}\sum_{\mathbf k}
\max(0,b_{\mathbf k}-R_{\mathrm{SID}}).\nonumber
\end{align}
Показатели различают неуникальные коды и исключение товаров;
нулевой $\operatorname{Drop}_{R}^{\mathrm{SID}}$ не исключает коллизий.
При $N_G>V_{\mathrm{SID}}^{L_{\mathrm{SID}}}$ не менее
$N_G-V_{\mathrm{SID}}^{L_{\mathrm{SID}}}$ назначений повторяются;
для контроля без повторений ёмкость увеличивают.
Обратное неравенство допускает уникальные коды, но не гарантирует их.

Обязательны три варианта при одинаковых $M,L_{\mathrm{SID}},V_{\mathrm{SID}}$
и бюджете подбора: случайные фиксированные коды, коды после
реконструкции без рекомендательного дообучения и совместно обучаемые
коды. В первых двух вариантах таблица назначений фиксирована:
кодовые векторы могут обучаться как входные представления декодера,
но их изменение не пересчитывает назначенные номера.
В формулах векторов этих вариантов вместо мягкого назначения
используется $\rho_{i\ell b}^{\mathrm{SID}}=\ind\{b=k_{i\ell}^{\mathrm{SID}}\}$.
Для случайного варианта требуется
$V_{\mathrm{SID}}^{L_{\mathrm{SID}}}\ge|\mathcal I_G|$.
Поток $m=2,j=6$ однократно вызывает
\code{choice(C, size=N_G, replace=False, shuffle=True)}, где
$C=V_{\mathrm{SID}}^{L_{\mathrm{SID}}}\le2^{63}-1$ и $N_G=|\mathcal I_G|$.
Так равномерно без возвращения выбираются различные числа из $0,\ldots,C-1$.
Числа сопоставляются товарам
в порядке их идентификаторов и записываются в системе
с основанием $V_{\mathrm{SID}}$ с ведущими нулями до длины
$L_{\mathrm{SID}}$; к каждой цифре прибавляется единица.
Так проверяется ценность семантической организации. При разной
частоте коллизий сравнивается и доступная доля каталога.

Производные, фазы обучения, допустимость выдачи и коллизии проверяют
X04--X05 (раздел~\ref{app:tests-extra}).

Память таблиц кодов составляет
$O(L_{\mathrm{SID}}V_{\mathrm{SID}}d_{\mathrm{SID}})$,
индекса --- $O(|\mathcal I_G|L_{\mathrm{SID}})$.
Для мини-пакета из $b_N$ запросов полнокаталожное нормирование требует
$O(b_N|\mathcal I_G|L_{\mathrm{SID}}d_{\mathrm{SID}})$ операций
после вычисления мягких векторов товаров. Каталог обрабатывается
блоками с точным накоплением нормирующей суммы: это экономит память,
но не сокращает число скалярных произведений.
Стоимость декодирования ограничивается
$O(L_{\mathrm{SID}}B_{\mathrm{SID}}
[(d_h^{\mathrm{SID}})^2+V_{\mathrm{SID}}d_h^{\mathrm{SID}}
+d_h^{\mathrm{SID}}d_{\mathrm{SID}}])$,
число читаемых товарных соответствий ---
$B_{\mathrm{SID}}R_{\mathrm{SID}}$.
Словарь и ширина поиска изменяются при прежнем ограничении задержки.
Допуск в обслуживаемую конфигурацию проверяется по пункту~\ref{sec:sid-serving};
исторический охват и причинный эффект оцениваются раздельно.

\section{Дополнительные способы отбора и формирования выдачи}\label{sec:model-extensions}

Эти варианты не входят в первый запуск. Каждый меняет отдельную часть
системы и сравнивается с фиксированной базой при том же пределе кандидатов.

\subsection{Ослабление рёбер с малым числом сессий}\label{sec:model-edge-optional}

Отдельное исследовательское сравнение может дополнительно уменьшать
вклад рёбер с малым числом поддерживающих сессий:
\begin{align}\label{eq:model-edge-ablation}
\chi^r_{ij}&=\frac{n^r_{ij}}{n^r_{ij}+\lambda_{\mathrm{edge}}},
\qquad \lambda_{\mathrm{edge}}\ge0,\nonumber\\
g^{r,\mathrm{edge}}_{qj}
&=\sum_{\ell\in\mathcal D_q}\alpha_{q\ell}
\rho^r_{i_{s(q)\ell}}\chi^r_{i_{s(q)\ell},j}
P^r_{i_{s(q)\ell},j}.
\end{align}
Коэффициент $\chi^r_{ij}$ вычисляют только для сохранённых рёбер,
где $n^r_{ij}\ge1$; при $\lambda_{\mathrm{edge}}=0$ исходный балл
восстанавливается. Этот вариант сохраняет те же $P^r$, топологию,
пути и квоту кандидатов, но меняет графовые баллы и требует нового
обучения ранжировщика на $B$. Он может подавлять полезные редкие
связи, поэтому не включается в основную конфигурацию автоматически.

\subsection{Дополнительное разнообразие выдачи}\label{sec:model-diversity-optional}

Дополнительная политика может уменьшать избыточное сходство товаров.
Следующий вероятностный вариант не входит в текущую постановку:
он потребовал бы отдельно определённых и проверенных оценок
$\widehat p_{qi}$ вероятности покупки. Баллы $f(x_{qi})$ не заменяют их.
Пусть $S_{ij}\in[0,1]$ --- симметричное сходство корректных
содержательных векторов: отрицательные косинусы заменены нулём,
доступные модальности объединены нормированной смесью.
При отсутствии общей модальности $S_{ij}=0$ с маской неизвестности.
Нулевое значение в этом случае не доказывает различие товаров.
Для множества $L\subseteq\mathcal C_q$ определим суррогат
\begin{equation}\label{eq:model-slate-surrogate}
U_q(L)=\sum_{i\in L}\widehat p_{qi}
-\frac{\lambda_D}{\max(1,K-1)}
\sum_{\{i,j\}\subseteq L}S_{ij},\qquad\lambda_D\ge0.
\end{equation}
Первое слагаемое предпочитает высокие наблюдаемые частоты покупки,
второе штрафует повторение сходных товаров.
$U_q$ не измеряется в денежных единицах и не является конверсией.
В основной версии этот функционал не вычисляется: список сортируют
по баллу ранжировщика, $\lambda_D=0$. Для вероятностного расширения
нужны отдельный план обучения и выбора и новый итоговый период;
подбор по результатам $D$ запрещён.

При $\lambda_D>0$ применяется жадное добавление:
\begin{equation}\label{eq:model-greedy-slate}
i_k=\argmax_{i\in\mathcal C_q\setminus L_{k-1}}
\left[\widehat p_{qi}-
\frac{\lambda_D}{\max(1,K-1)}\sum_{j\in L_{k-1}}S_{ij}\right],
\quad L_k=L_{k-1}\cup\{i_k\},\quad L_0=\varnothing.
\end{equation}
Добавление выполняется до $k=K_q$, в том числе при отрицательном
приращении: длина выдачи зафиксирована протоколом.
Этот алгоритм является эвристикой; глобальный максимум
\eqref{eq:model-slate-surrogate} не гарантируется.
Стоимость последовательного построения составляет порядка
$K|\mathcal C_q|$ вычислений сходства при обновлении накопленных сумм.

\subsection{Почему сумма товарных вероятностей не равна конверсии}

Даже если были бы известны истинные вероятности будущей покупки
каждого товара при одном и том же условии $\mathcal H_q$, события
покупки разных товаров зависимы. Для вероятности хотя бы одной
покупки товара из $L$ справедливы только общие границы
\begin{equation}\label{eq:model-union-bounds}
\max_{i\in L}p_i\le
\Prob\!\left(\bigcup_{i\in L}\{y_{qi}=1\}\given\mathcal H_q\right)
\le\min\!\left(1,\sum_{i\in L}p_i\right).
\end{equation}
Формула $1-\prod_{i\in L}(1-p_i)$ требует условной независимости,
которая для заменяющих и дополняющих товаров обычно не обоснована.
Оценки $\widehat p_{qi}$ дополнительно условны по баллу и процедуре
кандидатов, поэтому подстановка в эти границы служит лишь
иллюстрацией и не создаёт гарантированных вероятностных границ.
Ни одна из формул не оценивает причинный прирост покупки
при смене политики рекомендаций.

\subsection{Несколько шагов по графу}

Одношаговый граф допускает проверяемое расширение случайным
блужданием с возвратом к истории сессии.
Пусть $T$ --- смесь матриц отношений, нормированная по строкам;
нулевые строки заменяются историческим распределением $\pi^{\mathrm{pop}}$.
Пусть $h_q$ --- распределение начального веса по товарам
истории, полученное суммированием $\alpha_{q\ell}$.
На общем каталоге снимка зададим
\begin{equation}\label{eq:model-ppr}
z_q^{(n+1)}=\eta_Rh_q+(1-\eta_R)T^\top z_q^{(n)},
\qquad z_q^{(0)}=h_q,\quad 0<\eta_R\le1.
\end{equation}
На каждом шаге часть массы возвращается к известной истории,
остальная проходит по рёбрам. При пустой известной истории
этот источник отключается. Неизвестные графу товары не получают
фиктивных начальных вершин; оставшиеся веса перенормируются.

Поскольку $T$ стохастическая по строкам, преобразование сжимает
расстояние в норме суммы модулей с коэффициентом $1-\eta_R$.
Следовательно, существует единственное предельное распределение,
и для него при $n\ge1$
\begin{equation}\label{eq:model-ppr-error}
\norm{z_q^{(n)}-z_q^*}_1\le2(1-\eta_R)^n.
\end{equation}
Это граница вычислительной ошибки, а не качества рекомендаций.
Переходы выполняются на разреженной матрице;
вклад заменённых нулевых строк вычисляется одним умножением
вектора $\pi^{\mathrm{pop}}$ на их суммарную массу.
Полный шаг с явным вектором стоит $O(N_G+E_{\mathrm{sparse}})$,
где $N_G$ --- размер каталога, $E_{\mathrm{sparse}}$ --- число сохранённых
рёбер до дополнения нулевых строк. Линейный обход нужен и при отсутствии рёбер,
поэтому многократные обходы всего каталога могут нарушить бюджет
задержки. Источник принимается только после измерения добавочного
покрытия покупок и стоимости относительно одного шага.

\subsection{Поведенческие векторы товаров}

Item2Vec можно исследовать как дополнительный источник,
использующий совместные взаимодействия~\cite{item2vec}.
Пусть $\mathcal P_A$ --- пары товаров из временных и позиционных
окон периода $A_N$, ограниченные каталогом $\mathcal I_G$. Для каждого товара обучаются
два вектора $u_i^{B},v_i^{B}\in\R^{d_B}$: центральный и контекстный.
Один из вариантов функции обучения имеет вид
\begin{align}\label{eq:model-item2vec}
\mathcal L_B=&-\sum_{(i,j)\in\mathcal P_A}
\left[\log\sigmoid((u_i^{B})^\top v_j^{B})
+m_B\E_{k\sim Q_B}\log\sigmoid(-(u_i^{B})^\top v_k^{B})\right]
\nonumber\\
&+\frac{\lambda_B}{2}\sum_i(\norm{u_i^{B}}^2+\norm{v_i^{B}}^2),
\qquad Q_B(k)=\frac{n_k^{3/4}}{\sum_jn_j^{3/4}}.
\end{align}
Здесь $n_k$ --- число взаимодействий с товаром в $A_G$,
$m_B$ --- число случайных контекстов на положительную пару,
$\lambda_B>0$ --- регуляризация. Случайный контекст помогает
отличать характерные соседства от фоновой частоты;
он не утверждает, что пользователь отверг товар $k$.
Математическое ожидание заменяется средним по независимым
выборкам из зафиксированного $Q_B$.

Пусть $\mathcal D_q^B\subseteq\mathcal D_q$ содержит просмотры товаров,
для которых обучены поведенческие векторы. При непустом множестве
представление запроса равно
\begin{equation}\label{eq:model-item2vec-query}
u_q^B=\frac{\sum_{\ell\in\mathcal D_q^B}
\alpha_{q\ell}u_{i_{s(q)\ell}}^B}
{\sum_{\ell\in\mathcal D_q^B}\alpha_{q\ell}},
\qquad s_{qi}^B=(u_q^B)^\top v_i^B.
\end{equation}
При пустом $\mathcal D_q^B$ источник недоступен.
Использование разных центральных и контекстных векторов
согласовано с функцией обучения. Замена скалярного произведения
косинусом или усреднение $u_i^{B},v_i^{B}$ образуют отдельный вариант,
который требуется оценить. Поведенческие и содержательные векторы
не смешиваются покоординатно без специально обученного отображения.

Расширения отбора сравниваются при одинаковом общем бюджете
$M$ и на том же расписании запросов. Увеличение сложности
принимается только при воспроизводимом улучшении выбранного
показателя покрытия покупок и соблюдении бюджета обслуживания.
Их задача --- измерить вклад каждого способа отбора,
а не увеличить число компонентов обязательной системы.

\section{Эксперимент с реальными показами}\label{sec:eval-online}

На открытых данных протокол откладывают; до реальных показов он
обязателен вместе с проверкой исторического качества, технических
ограничений на целевом оборудовании и учёта исходов. Контроль ---
действующая политика или отсутствие блока; его фиксируют
до запуска \cite{experiments}.

\subsection{Единица назначения и причинный эффект}\label{sec:eval-itt}

До запуска фиксируют единицу назначения, допустимую популяцию,
горизонт воздействия и соответствующую единицу анализа.
Если сервис располагает устойчивым разрешённым account/user/device key,
вариант $Z_u\in\{0,1\}$ назначают по этому ключу при первом допуске,
$\Prob(Z_u=1)=1/2$, и закрепляют на весь эксперимент и созревание исходов.
$1$ --- новая система, $0$ --- контроль. Все сессии того же ключа
наследуют назначение; версия алгоритма назначения, ключ и результат
записываются до вызова ранжировщика. Правила смены ключа, нескольких
устройств и потери cookie фиксируют заранее; такие случаи не допускают
скрытой повторной рандомизации. Допуск определяется до воздействия.

Основная сессионная метрика сохраняет единицу из~\eqref{eq:eval-itt-estimand}.
В подтверждающую когорту входит только первая допущенная сессия
$s(u)$ каждого нового для эксперимента ключа; это правило применяют
до назначения. Основной исход $Y_{s(u)}(z)$ --- покупка после назначения
в этой же сессии до $t_{0,s(u)}+H$. Каждый ключ даёт ровно одну
сессию в основной знаменатель:
\begin{align}\label{eq:eval-unit-itt}
\Delta_{1000,\mathrm{entry}}^{\mathrm{ITT}}
&=1000\bigl(\E[Y_{s(u)}(1)]-\E[Y_{s(u)}(0)]\bigr),\nonumber\\
\widehat\Delta_{1000,\mathrm{entry}}^{\mathrm{ITT}}
&=1000\left(\frac{\sum_{u:Z_u=1}Y_{s(u)}}{N_1}
             -\frac{\sum_{u:Z_u=0}Y_{s(u)}}{N_0}\right),\quad N_0,N_1>0.
\end{align}
Здесь $N_z$ одновременно число назначенных ключей и сессий входной
когорты; единица --- покупающие входные сессии на тысячу, не все визиты
и не люди. Device key не тождественен человеку. Позднейшие сессии
наследуют вариант, но не пополняют основной знаменатель. Эффект на
любую покупку ключа за $H$ и отношение покупающих сессий ко всем
последующим сессиям --- только отдельные вторичные estimand;
число последующих сессий само может меняться под воздействием.
Они не подменяют основной критерий $S_{\mathrm{business}}$.
Для входной когорты формулы ниже применяют с $s=s(u)$ и $n_z=N_z$
лишь при независимых одинаково распределённых исходах между ключами
и отсутствии существенного межключевого воздействия. Если эти условия
не обоснованы, заранее задают анализ по единицам назначения с учётом
календарной зависимости; независимый биномиальный SE не используют.
Порог, мощность и долю охваченного трафика задают для этой когорты.

Только при отсутствии устойчивого ключа допустим ограниченный
краткосрочный сессионный протокол ниже. При первом допустимом запросе
назначают $Z_s\in\{0,1\}$, $\Prob(Z_s=1)=1/2$, и сохраняют вариант
до конца сессии и её окна исхода. Он не оценивает эффект на пользователя.

Пусть $Y_s(z)$ --- потенциальный исход: наличие хотя бы одной покупки
после назначения в той же сессии и до $t_{0,s}+H$, если сессии назначен
вариант $z$. Для конкретной сессии наблюдается только
\begin{equation}\label{eq:eval-observed-outcome}
 Y_s=Z_sY_s(1)+(1-Z_s)Y_s(0).
\end{equation}
\begin{assumption}[Условия сессионного эксперимента]
\label{ass:eval-randomization}
Назначение независимо от потенциальных исходов.
Политика одного варианта однозначно определена.
Нет существенного влияния назначения в одной сессии на исход другой:
например, через запоминание рекомендаций пользователем, общий дефицит
товаров или меняемую экспериментом популярность.
Учёт покупок одинаково полон в обоих вариантах.
Для независимых биномиальных интервалов дополнительно требуется
независимость одинаково распределённых исходов между сессиями внутри
каждой группы; сама рандомизация этого условия не гарантирует.
\end{assumption}

В Coveo нет устойчивого межсессионного идентификатора человека, и
это не определяет доступность ключей в будущем сервисе. При сессионном
назначении до запуска задают допустимые privacy-safe агрегированные
сигналы повторных визитов и план чувствительности. Доля вероятных
повторов --- диагностика, а не доказательство отсутствия переноса эффекта.
Если обоснованы верхние границы $r_z$ долей исходов, на которые могло
повлиять межсессионное воздействие в группе $z$, то при бинарном исходе
интервал разности расширяют с обеих сторон на $1000(r_0+r_1)$,
пересекая с $[-1000,1000]$. Оценённые по данным границы требуют отдельного
учёта их неопределённости и общего уровня ошибки. При отсутствии
обоснованной границы и правдоподобия предположения о невмешательстве
подтверждающий сессионный эффект не заявляют; результат описательный.

При условиях предположения~\ref{ass:eval-randomization} для входной
когорты либо ограниченного сессионного протокола эффект и оценка имеют вид
\begin{align}
 \Delta_{1000}^{\mathrm{ITT}}
 &=1000\bigl(\E[Y_s(1)]-\E[Y_s(0)]\bigr),
 \label{eq:eval-itt-estimand}\\
 \widehat\Delta_{1000}^{\mathrm{ITT}}
 &=1000\left(\frac{k_1}{n_1}-\frac{k_0}{n_0}\right),
 \quad n_z=\sum_s\ind\{Z_s=z\},\quad
 k_z=\sum_s\ind\{Z_s=z\}Y_s,\qquad n_0,n_1>0.
 \label{eq:eval-itt-estimator}
\end{align}
ITT --- анализ по первоначальному назначению. В $n_z$ сохраняют сессии
с ошибками, пустыми ответами и без фактического показа; $Y_s$ определяется
реальной покупкой, не успехом обслуживания. Учитывают все покупки
без привязки заказа к показу. Отбор только кликнувших или увидевших
блок нарушает исходное случайное сравнение.

\subsection{Интервал для бизнес-эффекта}\label{sec:eval-itt-ci}

Следующие формулы относятся к указанной основной когорте с независимыми
исходами и заранее фиксированным завершением. При $n_0,n_1>0$ и
достаточно больших количествах покупающих и непокупающих сессий
в каждой группе используют приближение
\begin{align}
 \widehat p_z&=k_z/n_z,\qquad
 \widehat{\operatorname{se}}_{\mathrm{ITT}}
 =1000\sqrt{\frac{\widehat p_1(1-\widehat p_1)}{n_1}
              +\frac{\widehat p_0(1-\widehat p_0)}{n_0}},
 \label{eq:eval-itt-se}\\
 I_{\mathrm{normal}}
 &=\left[\widehat\Delta_{1000}^{\mathrm{ITT}}
       -z_{1-\alpha/2}\widehat{\operatorname{se}}_{\mathrm{ITT}},
         \widehat\Delta_{1000}^{\mathrm{ITT}}
       +z_{1-\alpha/2}\widehat{\operatorname{se}}_{\mathrm{ITT}}\right].
 \label{eq:eval-itt-normal-ci}
\end{align}
Здесь $z_u$ --- квантиль стандартного нормального распределения.
Уровень интервала не является вероятностью истинности гипотезы.
При пустой группе оценка и интервал не определены: статус
\texttt{insufficient}, без подстановки нулевого эффекта.

При редких исходах можно заранее выбрать консервативный биномиальный
интервал. Для $k$ покупающих сессий из $n>0$, $0\le k\le n$,
границы обращают биномиальные хвосты через проверенную реализацию
квантилей бета-распределения:
\begin{subequations}\label{eq:eval-exact-beta}
\begin{align}
\ell&=\begin{cases}
0,&k=0,\\
Q_{\mathrm{Beta}(k,n-k+1)}(\alpha/4),&k>0,
\end{cases}\label{eq:eval-exact-lower}\\
u&=\begin{cases}
1,&k=n,\\
Q_{\mathrm{Beta}(k+1,n-k)}(1-\alpha/4),&k<n.
\end{cases}\label{eq:eval-exact-upper}
\end{align}
\end{subequations}
Здесь $0<\alpha<1$, $Q$ --- квантиль указанного распределения;
проверяют конечность результата и $0\le\ell\le u\le1$.
Для двух групп консервативный интервал эффекта равен
\begin{equation}\label{eq:eval-exact-diff}
 I_{\mathrm{exact}}=1000[\ell_1-u_0,\;u_1-\ell_0].
\end{equation}
Непокрытие каждой групповой вероятности не превосходит $\alpha/2$,
а разности --- $\alpha$ по неравенству для объединения событий.
Условия независимости заданы в~\ref{sec:eval-itt}
и допущении~\ref{ass:eval-randomization}; несколько сессий одного ключа
требуют кластерного анализа вне этого биномиального контракта.

\subsection{Правило принятия бизнес-результата}\label{sec:eval-acceptance}

Пусть $[\ell_B,u_B]$ --- заранее выбранный доверительный интервал
для $\Delta_{1000}^{\mathrm{ITT}}$ после завершения эксперимента.
Различают два утверждения:
\begin{align}
 \ell_B>0
 &\quad\Longrightarrow\quad
 \text{подтверждён положительный эффект},
 \label{eq:eval-positive-acceptance}\\
 \ell_B>\delta_{\mathrm{accept}}
 &\quad\Longrightarrow\quad
 \text{подтверждено превышение бизнес-порога}.
 \label{eq:eval-business-acceptance}
\end{align}
Порог $\delta_{\mathrm{accept}}\ge0$ в покупающих сессиях на тысячу
задают заранее. Утверждения действуют при предположениях и уровне
выбранной процедуры; превышения порога одной точечной оценкой недостаточно.
Остальные условия приёмки заданы в разделе~\ref{sec:acceptance}.
При назначении по ключу эти условия применяют к
$\Delta_{1000,\mathrm{entry}}^{\mathrm{ITT}}$ и порогу в покупающих
входных сессиях на тысячу; пользовательский вторичный исход
не подменяет этот сессионный критерий.

\subsection{Размер и длительность эксперимента}\label{sec:eval-power}

До запуска задают обнаруживаемое увеличение
$\delta_{\mathrm{plan}}>\delta_{\mathrm{accept}}$ покупающих сессий на тысячу, уровень ошибки $\alpha$
и желаемую мощность $1-\beta$ --- вероятность обнаружить эффект
предполагаемого размера. Пусть $p_0$ --- ожидаемая частота в контроле,
$\delta=\delta_{\mathrm{plan}}/1000$, $p_1=p_0+\delta$ и
$\bar p=(p_0+p_1)/2$.
Для равных групп при проверке нулевого эффекта ориентировочное число
независимых сессий в каждой группе равно
\begin{equation}\label{eq:eval-sample-size}
 n\approx\frac{
 \left[z_{1-\alpha/2}\sqrt{2\bar p(1-\bar p)}
       +z_{1-\beta}\sqrt{p_0(1-p_0)+p_1(1-p_1)}\right]^2}
 {\delta^2}.
\end{equation}
Применимость приближения проверяют для ожидаемой редкости исходов
и способа анализа. $p_0$ берут из сопоставимого потока или пилота;
частота Coveo не переносится автоматически в другой магазин.
При назначении по ключу планируют число входных сессий, равное числу
новых ключей; $p_0$ также относится к входной когорте. При доказанной
независимости этих исходов формула применима, но подстановка числа
всех последующих сессий недопустима.

Формула~\eqref{eq:eval-sample-size} планирует проверку отсутствия эффекта,
не нижнюю границу выше $\delta_{\mathrm{accept}}$.
Для бизнес-порога проверяют
$H_0:\Delta_{1000}^{\mathrm{ITT}}\le\delta_{\mathrm{accept}}$
против $\delta_{\mathrm{plan}}>\delta_{\mathrm{accept}}$;
нужный размер зависит от разрыва этих величин.

Для правила нижней границы нормального интервала получается простой
ориентир. Обозначим $d_*=(\delta_{\mathrm{plan}}-
\delta_{\mathrm{accept}})/1000$ и
$v_*=p_0(1-p_0)+p_1(1-p_1)$, где $p_1=p_0+\delta_{\mathrm{plan}}/1000$
должно лежать в $(0,1)$. При равных группах стандартная ошибка
разности долей приблизительно равна $\sqrt{v_*/n}$. Следовательно,
\begin{equation}\label{eq:eval-threshold-power}
n\approx\frac{v_*(z_{1-\alpha/2}+z_{1-\beta})^2}{d_*^2}.
\end{equation}
Условие $\sqrt{n}\,d_*/\sqrt{v_*}-z_{1-\alpha/2}=z_{1-\beta}$
планирует нижнюю границу выше порога при близости стандартной ошибки
к плановой. Для редких исходов мощность выбранного правила проверяют
суммированием вероятностей пар $(k_0,k_1)$ или моделированием.

При $v$ новых допущенных единицах основной когорты в секунду
ориентир набора равен $2n/v$,
затем требуется наблюдение $H+g$. Длительность также покрывает заданные
календарные циклы; меняющийся трафик и зависимость сессий влияют на мощность.

Срок или правило завершения набора фиксируют заранее; анализ проводят
после созревания всех исходов по правилу~\ref{sec:eval-srm-stop}.

\subsection{Проверка назначения и правило остановки}
\label{sec:eval-srm-stop}

До первого назначения сохраняют план $L$ календарных проверок,
уровень контроля SRM $\alpha_{\mathrm{SRM}}=0.001$ и ожидаемые доли $1/2$.
На каждой проверке считают уникальные назначенные единицы, включая
ошибки и отсутствие показа. Если план содержит $G$ проверяемых страт,
принадлежность $S_u=\phi(X_u^{\mathrm{pre}})$ вычисляют до назначения
из доступных к этому моменту признаков. Значение $S_u$ записывают
в журнал назначения и не меняют. План задаёт поля $X_u^{\mathrm{pre}}$,
версию и хеш правила $\phi$, обработку пропусков и ожидаемую долю
каждого варианта в каждой страте. Клики, показы, резервное обслуживание
и исходы после назначения не определяют SRM-страты; такие разрезы
допустимы лишь для отдельно обоснованного анализа.

Стандартный режим требует независимых назначений с
$\Prob(Z_u=1\mid X_u^{\mathrm{pre}})=1/2$. На каждой плановой проверке
дополнительно требуется: условно на состав и размер $N$ общей когорты
или каждой проверяемой страты число назначений варианта 1 имеет закон
$\mathrm{Bin}(N,1/2)$. Состав не адаптируют по истории назначений
или воздействию эксперимента. При этих условиях применяют точный
двусторонний биномиальный тест; $N$ включает всех назначенных участников
проверяемого состава независимо от обслуживания и исхода.
Без страт порог каждой проверки равен $\alpha_{\mathrm{SRM}}/L$;
при $G$ стратах порог для каждого теста, включая общий, равен
$\alpha_{\mathrm{SRM}}/[L(G+1)]$. Зависимое назначение или адаптивное
формирование состава требуют заранее заданного теста фактического
механизма с контролем общей ошибки, а не биномиальной подстановки.
План содержит правило состава общей когорты, метод теста,
обоснование условного закона и контроля общей ошибки.

SRM означает несоответствие долей назначения. При его срабатывании
приостанавливают новый набор, сохраняют все назначения и расследуют
потери логов, повторные ключи, изменения допуска и ошибки маршрутизации.
Исходы не удаляют ради восстановления пропорции. Дополнительно проверяют
баланс заранее выбранных признаков до воздействия, целостность логов
назначения и соответствие назначения фактическому обслуживанию.
Резерв не меняет вариант ITT. Баланс после показа не используется
для фильтрации участников; отсутствие SRM не доказывает корректность
учёта покупок или отсутствие вмешательства.

Защитный монитор получает только заранее объявленные показатели;
до запуска фиксируют их окна, знаменатели, минимальные объёмы, пороги,
ответственное действие и расписание проверок. Пересечение аварийного
порога останавливает воздействие; доступность такого канала сама по себе
не делает итоговый интервал устойчивым к зависимой от данных остановке.
Основной протокол имеет один итоговый анализ без промежуточных решений
по эффективности. Любая зависимая от данных досрочная
остановка, включая заранее предусмотренную защитную остановку или остановку
из-за SRM, а также любое незапланированное прерывание переводят причинный результат
в описательный: обычные интервалы и $p$-значения не объявляют
подтверждающими и не применяют правила~\eqref{eq:eval-positive-acceptance}
и~\eqref{eq:eval-business-acceptance}. После расследования для нового
подтверждающего вывода нужен новый заранее зарегистрированный эксперимент.
Сохранение подтверждающего анализа после запланированных промежуточных
остановок допустимо только по отдельно заданной до запуска процедуре
с общим контролем ошибки (например, alpha-spending или always-valid);
она не выбирается задним числом.

\subsection{Неполные исходы и защитные показатели}
\label{sec:eval-missing}

Созревший календарный горизонт не гарантирует полноты телеметрии.
Если в группе $z$ имеется $m_z$ сессий с неизвестным исходом,
а $k_z^{\mathrm{obs}}$ покупок достоверно зарегистрированы среди остальных,
то при $n_z>0$ без дополнительных предположений
\begin{equation}\label{eq:eval-missing-rate}
 \frac{k_z^{\mathrm{obs}}}{n_z}\le p_z^{\mathrm{sample}}
 \le\frac{k_z^{\mathrm{obs}}+m_z}{n_z}.
\end{equation}
Соответствующий диапазон совместимых выборочных эффектов равен
\begin{equation}\label{eq:eval-missing-effect}
 1000\left[
 \frac{k_1^{\mathrm{obs}}}{n_1}
 -\frac{k_0^{\mathrm{obs}}+m_0}{n_0},\quad
 \frac{k_1^{\mathrm{obs}}+m_1}{n_1}
 -\frac{k_0^{\mathrm{obs}}}{n_0}\right].
\end{equation}
Это диапазон совместимых результатов, не доверительный интервал:
статистическую неопределённость оценивают отдельно. Неизвестные исходы
не заменяют нулями и не удаляют из знаменателя ITT без обоснования.

До эксперимента задают пределы ухудшения частоты ошибок, задержки,
пустых выдач, регистрации покупок и доли недопустимых товаров.
Остановка по этим показателям допустима; статус причинного анализа
определяет раздел~\ref{sec:eval-srm-stop}. Неизвестные исходы,
правила зрелости и полноты проверяют по той же единице назначения:
при назначении по ключу $n_z,m_z,k_z^{\mathrm{obs}}$ относятся к
входным сессиям этих ключей, а не ко всем их сессиям или покупкам.

Длину сессии, клики и добавления в корзину показывают как вторичные
результаты. Несколько подтверждаемых эффектов требуют заранее заданной
поправки на множественность. Разрезы по признакам после назначения
не допускают причинной интерпретации без отдельного обоснования.

\section{Почему исторические данные не определяют прирост покупок}
\label{sec:eval-identification}

Неидентифицируемость установлена в утверждении~\ref{prop:business-nonidentification}.
Например, Hit@20 пустой выдачи равен нулю, поэтому его превышение новой
системой не доказывает прироста покупок относительно отсутствия рекомендаций.

\subsection{Условия оценки иной политики по журналу показов}
\label{sec:eval-ope}

Для оценки новой политики по журналу прежней в задаче с одним решением
нужны записи $(X_s,A_s^{\mathrm{act}},Y_s)$: контекст до назначения,
фактически показанный список и исход. Вероятность выбора списка
$b(a\mid x)$ прежней политикой должна быть известна.

Нужны отсутствие скрытых причин выбора списка, влияющих на исход,
корректный учёт показов и условие поддержки:
\begin{equation}\label{eq:eval-support}
 \pi(a\mid x)>0\ \Longrightarrow\ b(a\mid x)>0.
\end{equation}
При этих условиях оценка с обратными вероятностями равна
\begin{equation}\label{eq:eval-ips}
 \widehat V_{\mathrm{IPS}}(\pi)
 =\frac1N\sum_s
  \frac{\pi(A_s^{\mathrm{act}}\mid X_s)}
       {b(A_s^{\mathrm{act}}\mid X_s)}Y_s.
\end{equation}
Редкие списки с малой вероятностью $b$ получают большой вес,
поэтому дисперсия может быть неприемлемой \cite{ips}.
Обрезка весов уменьшает дисперсию ценой смещения.

Если дополнительно обучена оценка
$\widehat m(x,a)\approx\E[Y\mid X=x,A^{\mathrm{act}}=a]$,
возможно сочетание модели исхода и весов:
\begin{align}
 \widehat V_{\mathrm{DR}}(\pi)
 =\frac1N\sum_s\Bigg[&\sum_a\pi(a\mid X_s)\widehat m(X_s,a)
 \notag\\[-0.2em]
 &+\frac{\pi(A_s^{\mathrm{act}}\mid X_s)}
       {b(A_s^{\mathrm{act}}\mid X_s)}
       \bigl(Y_s-\widehat m(X_s,A_s^{\mathrm{act}})\bigr)\Bigg].
 \label{eq:eval-dr}
\end{align}
Модель исхода обучают без оцениваемой записи. Двойная устойчивость
означает согласованность при верном задании хотя бы одного из двух
компонентов и условиях причинной идентификации и регулярности.
Отсутствие поддержки этим не устраняется \cite{dr}.

Формулы~\eqref{eq:eval-ips}--\eqref{eq:eval-dr} относятся к одному решению
и фиксированному последующему режиму. Последовательные рекомендации
меняют будущий контекст; подстановка отдельных запросов с общей
сессионной покупкой требует другого обоснования.

В Coveo нет нужных вероятностей назначения и полного журнала показов
рекомендаций. Результаты поиска с кликами их не заменяют, поэтому IPS
и DR исключены из основного расчёта. Предсказанные вероятности покупки
также не заменяют вероятности назначения $b$.

\subsection{Что может подтвердить воспроизведение событий}
\label{sec:eval-replay}

Воспроизведение проверяет доступность признаков, состояние сессии,
поиск кандидатов, задержку и ограничения. Искусственный пользователь
проверяет сценарии при объявленной модели поведения; оба способа
не подтверждают прироста реальных покупок.

Различение охвата и эффекта проверяет тест~\ref{test:eval-effect}.


\section{Конфигурации запуска и реестр параметров}\label{app:parameters}

\subsection{Фиксированная конфигурация первого запуска}\label{app:first-config}

Первый запуск --- GCS+Rk: три отношения графа, описания, изображения,
поисковый запрос и LambdaMART. Используют прямые связи и неизменный
граф в протоколе \code{observed_purchases_v1}, без нейросетевых
источников и подбора параметров. Ниже заданы
начальные настройки, а не измеренные оптимумы на Coveo.

\paragraph{Данные и время.}
До обучения фиксируют границы семи частей, отбор сессий,
версии файлов и порядок идентификаторов. Граф, каталог и статистики
строят по $A_G$, ранжировщик обучают на $B$. Зарезервированный $A_N$
не используют в базе. Первый отчёт получают на $C_{\mathrm{smoke}}$;
$C_{\mathrm{select}},E,D$ закрыты.
\begin{align}
H&=86\,400\ \text{с},&g&=1800\ \text{с},&K&=20,\notag\\
Q_{\max}&=1,&M&=300,&d_G&=100,\notag\\
\Delta_G&=1800\ \text{с},&\tau_G&=600\ \text{с},&n_G&=3,\notag\\
\tau_H&=300\ \text{с},&\Delta_S&=600\ \text{с},&\lambda_G&=10.
\label{eq:app-start}
\end{align}
Граф использует $B_G=10$ предыдущих товарных временных групп,
его представление сессии --- $L_H=20$ последних просмотров
по разделу~\ref{sec:model-graph}. Хранят $L_{\mathrm{hist}}=20$
временных групп; сглаживание популярности: $\eta_P=1$.
Признак покупочной частоты не используется.

\paragraph{Отбор и признаки.}
Поиск по векторам описаний, изображений и запроса выполняют точно
по всему допустимому каталогу. Квоты фиксированы:
\begin{center}
\begin{tabular}{L{0.63\textwidth}r}
\toprule Источник & Квота \\\midrule
Каждое из трёх отношений графа & 50 \\
Описание товара & 40 \\
Изображение & 30 \\
Поисковый запрос & 50 \\
Недавние просмотры & 20 \\
Историческая популярность & 10 \\
\midrule Всего & 300 \\
\bottomrule
\end{tabular}
\end{center}
Устранение повторений, маски недоступных источников и заполнение
по популярности выполняют по~\eqref{eq:model-candidate-fill}.
Ранжировщик получает отдельные графовые и содержательные баллы;
смеси с коэффициентами $\omega_r$ и $\gamma_m$ не используются.
Дополнительный множитель $\chi_{ij}^r$ выключен.
Допуски вычислений двойной точности: $\eps_v=10^{-12}$,
$\eps_f=10^{-8}$. Порядок равных баллов и категориальные обозначения
фиксируют вместе со схемой признаков.

\paragraph{Ранжировщик и выдача.}
Версия LightGBM, функция производных и все настройки заданы
в~\ref{sec:model-ranker-implementation}. Бюджет первого запуска
не выбирают по $C_{\mathrm{select}}$. Штраф сходства выключен: $\lambda_D=0$.

\paragraph{Проверка и завершение опыта.}
Первый запуск завершается проверками пункта~\ref{app:tests-base}
и отчётом на $C_{\mathrm{smoke}}$ по пункту~\ref{app:diagnostic-output}.
Выполнение нагрузочных требований точным поиском проверяют измерением;
приближённый поиск вводят отдельным сравнением до открытия $D$.

Калибровка отключена. Итоговая оценка следует
разделам~\ref{sec:eval-observed} и~\ref{sec:eval-primary-comparison}.
Начальные $\alpha=0.05$ и
$B_{\mathrm{boot}}=2000$ не заменяют план точности и условия
применимости временных интервалов.

\subsection{Параметры нейросетевого расширения}\label{app:neural-config}

Граф, каталог и статистики сохраняют из $A_G$; представления обучают
на $A_N$. Квоты перераспределяют при общем пределе $M=300$:
\begin{longtable}{L{0.38\textwidth}L{0.15\textwidth}L{0.15\textwidth}L{0.17\textwidth}}
\toprule Источник & База & Нейросеть & С SID \\\midrule\endhead
Каждое из трёх отношений графа & 50 & 30 & 30 \\
Описание товара & 40 & 30 & 30 \\
Изображение & 30 & 20 & 20 \\
Поисковый запрос & 50 & 40 & 40 \\
Недавние просмотры & 20 & 20 & 20 \\
Популярность & 10 & 10 & 10 \\
Смешанный нейросетевой вектор & 0 & 30 & 20 \\
Каждое из двух намерений & 0 & 30 & 10 \\
Семантический код & 0 & 0 & 50 \\
\midrule Всего & 300 & 300 & 300 \\
\bottomrule
\end{longtable}
Неактивное намерение возвращает пустой список. После устранения
повторений свободные места заполняются по популярности согласно
формуле~\eqref{eq:model-candidate-fill}; её начальная квота
не ограничивает дополнительное заполнение.
При последующей проверке $J_I=4$ квоты одного намерения равны 15
для нейросетевой архитектуры и 5 для варианта с SID; общие нейросетевые
квоты 90 и 40 сохраняются. При $J_I=1$ единственное направление
получает соответственно 60 и 20 мест сверх смешанного источника.


Начальная конфигурация задаётся как
\begin{align}
d&=64,&J_I&=2,&H_P&=2,&d_P&=50,\notag\\
B_P&=100,&R_P&=20,&c_G&=0.1,&\eta_N&=1,\notag\\
\tau_P&=1,&\tau_N&=0.1,&\tau_C&=0.1,&\tau_{SN}&=300\ \text{с}.
\label{eq:app-modern-start}
\end{align}
Для трёх отношений исходно $\zeta_r=1/3$, предел прямой поисковой
примеси $\bar g_N=0.5$. Начальный фактический вес при доступном
поиске задаётся как $g_0=0.1$: исходно
$v_g=w_g=0$, $a_g^t=0$, а свободный член равен
\begin{equation}\label{eq:app-search-init}
b_g=\log\frac{g_0}{\bar g_N-g_0}=\log\frac14.
\end{equation}
При отсутствии поискового сигнала маска по-прежнему даёт
$g_q^{SN}=0$. Эти значения не ограничивают отдельно влияние
$\beta_S$ на распределение весов намерений.
Веса потерь:
\begin{equation}\label{eq:app-modern-loss-start}
\lambda_{\mathrm{next}}=0.2,\qquad \lambda_C=0.05,\qquad
\lambda_G^N=10^{-3},\qquad \lambda_\theta=10^{-5}.
\end{equation}
В исходном расширении $p_S=0$: случайное отключение поиска
не применяется. Диагностический порог совпадения намерений
фиксируют как $\delta_I=0.05$.
Для вычислений с 32-битной плавающей точкой исходно принимают
$\eps_N=10^{-6}$. Вероятность отключения каждой доступной ветви
при построении контрастивных представлений равна $0.2$ с восстановлением
ветви по правилу раздела~\ref{sec:modern-contrastive}.
Оптимизатор, инициализация, пакеты и три эпохи обучения заданы
в~\ref{sec:modern-training-recipe}. Контрольные точки после одной и пяти эпох
оставляют для отдельно запланированного продолжения; их проверка
также входит в бюджет выбора на $C_{\mathrm{select}}$.

По фактическим $|\mathcal I_G|$ и числу обучающих запросов измеряют
стоимость точной эпохи. При превышении бюджета применяют отдельно
обозначенный режим приложения~\ref{sec:sampled-training}: $m=512$,
$\delta=0.1$. Значения $m=1024,2048$ требуют отдельного плана;
выборочный знаменатель меняет обучающую процедуру.


\subsection{Начальные настройки отложенного исследования SID}\label{app:sid-config}

Отложенное исследование по приложению~\ref{sec:sid-overview}:
\begin{align}
L_{\mathrm{SID}}&=4,&V_{\mathrm{SID}}&=32,&d_{\mathrm{SID}}&=16,\notag\\
d_h^{\mathrm{SID}}&=64,&B_{\mathrm{SID}}&=32,&R_{\mathrm{SID}}&=4,
\label{eq:app-sid-start}
\end{align}
и $m_{\mathrm{SID}}=50$. Начальные веса
$\lambda_R^{\mathrm{SID}}=1$,
$\lambda_X^{\mathrm{SID}}=0.1$,
$\lambda_B^{\mathrm{SID}}=0.01$,
$\lambda_2^{\mathrm{SID}}=10^{-5}$; $\tau_R^{\mathrm{SID}}=1$.
Для заранее заданных $T_{\mathrm{SID}}$ совместных шагов допустимо
расписание
\begin{equation}\label{eq:app-sid-schedule}
\sigma_t^{\mathrm{SID}}=\pos{1-t/T_{\mathrm{SID}}},\qquad
\tau_t^{\mathrm{SID}}=0.2+0.8\pos{1-t/T_{\mathrm{SID}}}.
\end{equation}
Фазы обучения заданы в~\ref{sec:modern-training-recipe}
и разделе~\ref{sec:sid-decoder}. Проверяют загрузку жёстких кодов
всего каталога: $32^4=1\,048\,576$ возможных кодов не исключают
коллизий при меньшем числе товаров. Увеличение $L_{\mathrm{SID}}$
или $V_{\mathrm{SID}}$ обосновывают измеренными потерями отбора.


\subsection{Паспорт запуска и этапные проверки}
\label{app:run-manifest}

Один паспорт хранит закреплённые входы,
код и настройки опыта в JSON; результаты этапов сохраняются отдельно
и ссылаются на этот паспорт. Отдельный снимок паспорта и файл SHA-256
для каждого этапа не обязательны. До открытия $C_{\mathrm{select}}$,
$D$ и до назначения вариантов A/B соответствующий план сохраняют
неизменяемым; ссылка на версию либо контрольную сумму связывает его
с результатами. До расчётов проверяют, что код, настройки и входы
совпадают с планом; результат не может изменить уже закреплённый план.
Нормативный шаблон
\code{documents/math/coveo_session_recommender_run_manifest.template.json}
не готов к запуску: неизвестные даты, хеши и измерения имеют значение
\code{null}. Их не заменяют вымышленными данными.
Шаблон \code{manifest_version=2.10} задаёт режим полем
\code{evaluation_protocol}; паспорт фиксирует
правило первого запроса, зарегистрированную цель, совокупность оценки,
пороги точности в долях и отключённые компоненты
по разделу~\ref{sec:eval-observed}. Неизвестное покрытие не подменяется полным.
Типы, состояния, временные связи, целостность входов и условия применимого этапа
проверяются обычными проверками данных и тестами. Отдельная система
регистрации тестов не требуется; версия проверяющего кода воспроизводима.
Проверка B13 включает допустимые и недопустимые паспорта.

\begin{longtable}{L{0.24\textwidth}L{0.66\textwidth}}
\toprule Группа полей & Обязательное содержимое \\\midrule\endhead
Данные & Версии и SHA-256 трёх таблиц, схема, размерности векторов,
покрытие и разрывы наблюдения; UTC и целочисленные миллисекунды,
правила канонизации и сессионной подвыборки. \\
Периоды & Границы всех семи частей как $[a,b)$; непересечение,
целые сессии и допуск~\eqref{eq:eval-admission}.
Календарные поля имеют префикс \code{period_}, матрицы --- \code{matrix_}. \\
Код и модели & Хеш исходного кода, файл закреплённых зависимостей
с версиями и хешами пакетов, схема признаков, функция производных,
полные настройки активных модулей, оборудование, типы чисел,
генераторы случайных чисел и их состояния. \\
Выбор и оценка & \code{evaluation_protocol=observed_purchases_v1},
варианты и правило выбора, основное сравнение,
чувствительность к метаданным, план точности и временной зависимости.
Обязательность определяется таблицей этапов ниже. \\
Обслуживание & Нагрузка и лимиты, правила доставки и восстановления,
доверенные источники, сроки хранения, мониторинг, ответственная роль;
обязательны при проверке сервиса и допуске к A/B. \\
Проверки & Стандартные тесты сопоставлены с ID требований; их код
задаёт вход или рецепт, независимый эталон и допуски. Отдельные результаты
этапов содержат фактический исход, версию кода, конфигурацию и seed.
Отдельный хешированный паспорт каждого теста не требуется. \\
\bottomrule
\end{longtable}

\paragraph{Условия по этапам.}\label{app:run-stages}
Поле \code{stage} задаёт проверяемый этап. Таблица определяет
обязательность полей; данные будущих и неприменимых этапов не требуются.
Каждая запись результата содержит ID требования/теста, состояние,
причину, фактический отчёт либо ссылку на него и воспроизводимые
версию кода, конфигурацию и seed. \code{null} допустим до срока исполнения,
но не заменяет обязательный результат при завершении этапа.
Для отключённого модуля записывают \code{not_used} и причину;
при его включении применяются те же сроки.
$E$ зарезервирован и не используется; калибровка отключена.

\begin{longtable}{L{0.18\textwidth}L{0.72\textwidth}}
\toprule Этап & Условия завершения \\\midrule\endhead
\code{preflight} &
До обучения: данные, периоды, код, параметры активных моделей,
численные допуски заполнены; применимые проверки данных, канонизации,
допуска и математических свойств подготовки выполнены до расчётов.
Обученных моделей, плановых параметров точности оценки $D$
и будущих результатов нагрузки не требуется. \\
\code{post_training} &
Сохранены обученные модели и их хеши; выполнены применимые
математические проверки обучения и разделения данных.
Реальный протокол A/B и B18 на этом этапе не обязательны. \\
\code{post_load} &
Отдельная приёмка сервиса, не предпосылка локальной исторической оценки.
Сохранены три нагрузочных прогона, память, задержки, отказы,
проверки воспроизведения и восстановления. Результат испытания
может быть отрицательным; это запрещает приёмку, но не запись результата. \\
\code{pre_selection} &
До открытия $C_{\mathrm{select}}$ зафиксированы варианты, бюджет
и правило выбора; имеются результаты обучения, локальные измерения
стоимости и применимые проверки участвующих в выборе вариантов.
План включает конфигурацию варианта без метаданных.
Нагрузочная приёмка сервиса не требуется.
Доработка по $C_{\mathrm{select}}$ запрещена. \\
\code{pre_D} &
Зафиксированы основная политика либо пара и весь план оценки:
полуширина в долях, плановая дисперсия, допущения и множитель зависимости,
происхождение пилота на $C_{\mathrm{smoke}}$, календарная сетка и блоки;
выполнены условия выбора и применимые проверки оценки, включая
полный вариант без метаданных.
Ошибки ранжировщика и нарушение исторического протокола блокируют расчёт;
готовность сервиса для \code{ideal_deterministic} не требуется.
Результатов $D$ ещё нет. \\
\code{post_D} &
Сохранены итоговые оценки, интервалы и состояния
раздела~\ref{sec:acceptance}, включая отрицательные и недостаточные
результаты. План задним числом не изменяется. \\
\code{pre_AB} &
До назначения вариантов заполнены протокол эксперимента,
ключ назначения, расчёт объёма и срока, SRM, аварийные условия,
ответственные, доверенные ключи и производственные лимиты.
План A/B сохранён неизменяемым; B18, обслуживание и применимые защитные
проверки пройдены, выполнена готовность~\eqref{eq:acceptance-pilot}. \\
\code{post_AB} &
После созревания исходов сохранены результаты по назначенным
вариантам, проверки целостности и сведения об остановках. \\
\bottomrule
\end{longtable}

Сроки применимых тестов заданы в~\ref{app:tests-normative}.
После выбора окончательную модель фиксируют неизменяемо
и проверяют её связь с результатом выбора.
Подготовка, обучение и оценка перед расчётами проверяют относящиеся
к ним данные, математические свойства и неизменность применимого плана.
Сохранённые проверки повторяют при изменении затронутого кода или входов.
Для локального маршрута \code{post_load} не является предшественником
\code{pre_selection} или \code{pre_D}.
Отсутствие обязательного отчёта
даёт \code{not_run}, нарушение выполненного теста --- \code{failed}.
Пригодность данных, точность оценки и качество модели учитываются
раздельно по разделу~\ref{sec:acceptance}.

\subsection{Предварительный календарный план}

Предварительные доли времени для $A_G,A_N,B,C,E,D$:
$20\%$, $20\%$, $25\%$, $15\%$, $10\%$, $10\%$.
После календарного допуска публикуют фактические размеры
$N_J^{\mathrm{obs}}$ и $N_J^+$; $E$ остаётся резервом.
Внутри $C$ исходно отводят равные календарные половины
$C_{\mathrm{smoke}}$ и $C_{\mathrm{select}}$; к каждой отдельно применяется
допуск по $H+g$. Доли служат только исходной сеткой планирования.
Если они не обеспечивают допущенные сессии либо требуемое число временных
блоков, до обучения меняют план по раннему потоку или сохраняют
исследовательский статус. Нехватку продолжительности нельзя исправить
делением зависимого периода на искусственно короткие независимые блоки.
Период короче $H+g$ непригоден. Границы и правило дополнительного
сбора задают до обучения по предшествующим данным; покупающие сессии
выборочно не добавляют и результаты $D$ для этого не используют.


\subsection{Источники параметров}\label{app:parameter-roles}

Первый запуск использует рецепт~\ref{app:first-config} без подбора.
В следующих сравнениях гиперпараметры выбирают только на
$C_{\mathrm{select}}$; параметры модели обучают на назначенном периоде.
Таблица указывает место определения, не повторяет значения.

\begin{longtable}{L{0.34\textwidth}L{0.57\textwidth}}
\toprule Группа & Где определяется и фиксируется \\\midrule\endhead
Постановка: $H,g,K,Q_{\max}$ &
Правила запросов~\ref{sec:requests}; фиксируются до отбора сессий. \\
Граф и история: $\Delta_G,B_G,\tau_G,\tau_H,n_G,d_G,\lambda_G$,
$L_H,L_{\mathrm{hist}}$ &
Раздел~\ref{sec:model-graph}; статистики только из $A_G$.
Пределы просмотров и временных блоков различаются. \\
Источники: $M,m_a,\Delta_S,\eta_P$ &
Рецепт~\ref{app:first-config}; сравнения сохраняют общий бюджет.
Квоты не заменяют заполнение после устранения повторов. \\
Ранжировщик &
Раздел~\ref{sec:model-ranker-implementation}; обучение на $B$.
Штраф $\lambda_D=0$ в базе. \\
Точность и интервалы &
Разделы~\ref{sec:eval-precision}, \ref{sec:eval-time-block}.
План и предпосылки фиксируют до открытия $D$ по разрешённому пилоту
на $C_{\mathrm{smoke}}$; для \code{preflight} и обучения они не требуются. \\
Обслуживание и численные допуски &
Разделы~\ref{sec:ops-load-profile}, \ref{sec:ops-numerics}.
Пороговые значения фиксируют до выбора модели. \\
Нейросеть: размерности, пути, температуры, веса потерь &
Приложение~\ref{app:neural-config}; обучение на $A_N$ при
неизменных графе, каталоге и статистиках $A_G$.
Правила оптимизации --- в~\ref{sec:modern-training-recipe}. \\
Поисковая примесь: $\bar g_N,g_0,p_S$ &
$0\le\bar g_N\le1$; при положительной границе
$0<g_0<\bar g_N$. Нулевая прямая примесь не отключает влияние поиска
на веса намерений. Начальные значения --- в~\ref{app:neural-config}. \\
Выборочный знаменатель: $m,\delta$ &
Приложение~\ref{sec:sampled-training}; отдельная обучающая процедура,
предложение строится по $A_G$. \\
Последовательность и старение &
Приложения~\ref{sec:sequence-control}, \ref{sec:freshness};
окно и старение истории не заменяют внутрисессионное затухание. \\
SID: размерности, словари, расписания, потери &
Приложение~\ref{app:sid-config}; обучение на $A_N$ при
фиксированной нейросети. В обучаемом варианте
$\lambda_R^{\mathrm{SID}}>0$, остальные штрафы неотрицательны. \\
Простые смеси и штраф редких рёбер &
Раздел~\ref{sec:model-extensions}; отдельные сравнения.
В базе $\omega_r,\gamma_m,\lambda_{\mathrm{edge}}$ не используются. \\
Эксперимент с показами &
Раздел~\ref{sec:eval-online}; $\alpha$, мощность, бизнес-порог
$\delta_{\mathrm{accept}}$ и альтернатива
$\delta_{\mathrm{plan}}>\delta_{\mathrm{accept}}$
фиксируются до назначения вариантов. \\
\bottomrule
\end{longtable}

\subsection{Размеры периодов и достаточность данных}

Выборка из 200--500 тысяч сессий может быть достаточна для инженерного
прототипа, но число сессий с зарегистрированной целью заранее неизвестно.
При открытии периода с целями регистрируют
\begin{equation}\label{eq:app-positive-counts}
N_J^{\mathrm{obs}},\qquad N_J^+=|\mathcal S_J^+|,\qquad
N_J^{\mathrm{pair}}=\sum_{s\in\mathcal S_J^+}\ind\{|\mathcal J_{q_s^{(1)}}|>0\}.
\end{equation}
Первые два числа считают сессии, третье --- запросы с обучающими
парами; строки кандидатов не являются независимыми покупками.
Для $A_N$ также указывают число запросов с $Y_q\cap\mathcal I_G\ne\varnothing$,
число пригодных контрастивных пар и долю целей вне каталога.
Достаточность данных проверяют для каждого этапа.

Условия точности охвата заданы
в~\ref{sec:eval-precision} и~\ref{sec:eval-information}.
Дополнительный сбор допускается только по предварительному плану,
не до достижения удобного числа покупок.

\clearpage
\section{Порядок сравнений и отделение вклада компонентов}\label{app:ablation}

После запуска~\ref{app:first-config} выбирают следующее сравнение
по разделу~\ref{sec:eval-next-step}. До запуска фиксируют вопрос,
данные и бюджет. Таблицы ниже перечисляют допустимые исследования;
выполнять все строки не требуется. P и R --- обязательные описательные
контроли первой оценки; G--GCS служат диагностикой на $C_{\mathrm{smoke}}$.
Основной контроль исследования улучшений --- GCS+Rk.

\begin{longtable}{L{0.12\textwidth}L{0.40\textwidth}L{0.38\textwidth}}
\toprule Вариант & Состав & Проверяемый вопрос \\\midrule\endhead
P & Историческая популярность & Превосходит ли система доступный
неперсонализированный ответ? \\
R & Повторные просмотры и популярность & Объясняется ли результат
простым напоминанием о просмотренном? \\
G & Отношения графа и резервные источники & Полезны ли переходы между
товарами сверх повторных просмотров? \\
GC & Граф и признаки товара & Улучшает ли содержательная близость охват при
том же пределе кандидатов? \\
GCS & Граф, признаки товара и запрос поиска & Даёт ли доступный поисковый
контекст дополнительную информацию? \\
GCS+Rk & Те же источники и обученный ранжировщик & Уменьшает ли
ранжирование разрыв между охватом кандидатов и итоговой выдачи? \\
\bottomrule
\end{longtable}

Вторая группа сравнений отделяет вклад частей обучаемого модуля.
Базовые источники сохраняются, если в строке не оговорено иное.
\begin{longtable}{L{0.25\textwidth}L{0.65\textwidth}}
\toprule Вариант & Изменение относительно полной обучаемой модели \\\midrule\endhead
Без графовой ветви & Исключается $e_i^G$; остальные ветви и обучение
на $A_N$ сохраняются. \\
Только прямые связи & $H_P=1$; проверяется вклад непрямых путей
при прежнем числе сохраняемых соседей. \\
Одно намерение & $J_I=1$; суммарная квота источников по намерениям
переходит единственному направлению. \\
Без содержательных признаков товара & Исключаются ветви $e_i^T,e_i^V$ нового
кодировщика; идентификаторы, граф и доступный поисковый контекст сохраняются. \\
Без поиска в кодировщике & Маска $M_q^S$ обнуляется во всём новом
модуле, включая резервное представление только по поиску;
базовый поисковый источник не меняется. \\
Без прямой примеси поиска & Только $\bar g_N=0$; влияние поиска
на веса намерений и остальные правила сохраняются. \\
Без контрастивной части & $\lambda_C=0$; остальные цели и доступные
модальности сохраняются. \\
Без следующего просмотра & $\lambda_{\mathrm{next}}=0$; остаётся
обучение по покупкам и выбранные регуляризаторы. \\
Полная модель & Все включённые компоненты раздела~\ref{sec:modern-model};
квоты основной архитектуры из стартового плана. \\
Последовательный контроль & Замена кодировщика сессии одним
блоком причинного внимания при сохранении графового товарного
кодировщика; те же входы, покупочная цель и бюджет кандидатов. \\
Выборочный знаменатель & Приближённая потеря
приложения~\ref{sec:sampled-training}; точная версия служит ориентиром
на доступном масштабе. \\
Три варианта SID & Случайные фиксированные коды, фиксированные коды
после восстановления, совместно обучаемые коды; квоты архитектуры с SID. \\
\bottomrule
\end{longtable}

В сравнении одного компонента сохраняют $A_N$, прочие цели
и расписание обучения. Последовательный кодировщик сравнивают
с основной моделью при $\lambda_C=0$ у обеих. При изменении $J_I$
сохраняют общую квоту и $M$. Допущение о статических векторах
проверяют отдельным исключением всех предоставленных векторов,
включая базовые содержательные источники и поиск: это сравнение
доступной информации, а не одной ветви.

При сравнении отбора фиксируют алгоритм и бюджет подбора ранжировщика,
но обучают его заново на кандидатах каждого варианта в $B$.
Сравнивают кандидатную границу и выдачу.
Для сравнения только ранжировщиков кандидаты
фиксируют одинаковыми; коэффициенты простой смеси выбирают на $C_{\mathrm{select}}$.
Намерения проверяют по разделу~\ref{sec:modern-intent-diagnostics};
высокая энтропия весов не доказывает различие направлений,
штраф разнообразия исходно выключен.

Граф, каталог и статистики строят по одному $A_G$; дополнительное
обучение на $A_N$ отмечают в отчёте. Вклад архитектуры сравнивают
при одинаковых обучающих данных. Одинаковое $M$ не означает равную
стоимость: до запуска ограничивают число конфигураций, время и память;
сообщают оборудование, завершённые запуски и причины отказов.

Следующий этап, отдельный от первого запуска, допускает не более
12 полных конфигураций суммарно. Исходный план из девяти:
GCS+Rk и нейросетевые $J_I=1,2$, каждая с бюджетом
\code{num_boost_round}$\in\{100,200,400\}$; фактическое $N_T$
не превышает выбранного бюджета.
Нейросетевые варианты обучают три эпохи одинаковой процедурой;
случайное состояние --- 17. Дополнения регистрируют до оценки.
Состояния 29 и 43 предназначены для отдельного исследования
устойчивости; лучшую реализацию по $D$ не выбирают.

Формат отчёта задан ниже. Улучшение NDCG не заменяет основной критерий
и эксплуатационные ограничения. До открытия $D$ фиксируют одну
основную пару: выбранная на $C_{\mathrm{select}}$ модель и GCS+Rk. Прочие сравнения
исследовательские; подтверждение нескольких превосходств требует
заранее заданного семейства проверок и поправки на множественность.
Выбор улучшения по результатам $D$ требует новой итоговой выборки.
Бизнес-эффект оценивают отдельным экспериментом с реальными показами.

\subsection{Обязательная чувствительность к метаданным}
\label{sec:metadata-sensitivity}

Основное сравнение с предоставленными векторами условно относительно
их неизвестной даты построения. Обязательный вариант
\code{graph_recent_popularity_no_provided_metadata} исключает все
предоставленные текстовые, визуальные и поисковые векторы, категорию,
цену, их маски и производные признаки из отбора и ранжировщика.
Поисковые клики/списки и контекст запроса также не используются.
Остаются ID, доступные события, граф, повторные просмотры и статистики
популярности, построенные исключительно на $A_G$; нейросеть и SID
выключены. Каталог $\mathcal I_G$, сессии, расписание запросов,
правила допуска, $K$ и $M$ сохраняются, чтобы не менять знаменатель.

Квоты трёх графовых отношений равны 50 каждая, повторных просмотров ---
20, начальная квота популярности --- 10; оставшиеся до $M=300$ места
после устранения повторов заполняются популярностью. Новый ранжировщик
обучается на своих кандидатах $B$ тем же бюджетом и callback,
с заранее зафиксированной схемой без метаданных; $\lambda_D=0$.
Конфигурация и её хеш фиксируются до $C_{\mathrm{select}}$.
Полная реализация и применимые проверки этого варианта обязательны
перед его расчётами и к итоговой оценке, но не блокируют первый $C_{\mathrm{smoke}}$.
На $D$ обязательно публикуют цепочку $V,U,\operatorname{Hit}$,
разность с базой, временные интервалы чувствительности и стоимость.

Это обязательная проверка зависимости результата от информации,
а не вторая автоматически подтверждающая гипотеза. Основная пара
остаётся единственной; если заявляют несколько превосходств, до $D$
фиксируют семейство и поправку на множественность. Устойчивость этой
проверки не доказывает отсутствие утечки в неизвестном снимке;
неустойчивость ограничивает применимость исходного результата.
Замена предоставленных метаданных при неизменных событиях не должна
менять кандидатов, признаки или ответ варианта без метаданных.

\subsection{Обязательный отчёт первого запуска и сравнений}\label{app:diagnostic-output}

Таблица задаёт поля отчёта, а не результаты на Coveo. На первом запуске
её заполняют для $C_{\mathrm{smoke}}$; на $D$ --- после фиксации протокола. Недоступные
величины получают состояние с причиной, а не подставленное число.
Для локального \code{ideal_deterministic} стоимость означает измеренные
локальные время и память; показатели сервисной нагрузки остаются
неизмеренными до \code{post_load} и не блокируют исторический отчёт.
Протокол \code{observed_purchases_v1} указывают явно.
Все показатели качества относятся к $\mathcal S_J^+$; затраты
учитывают и сессии без зарегистрированной цели. Полнота телеметрии
остаётся неизвестной; конверсию на всех сессиях не оценивают.

\begin{longtable}{L{0.25\textwidth}L{0.65\textwidth}}
\toprule Блок & Обязательные величины \\\midrule\endhead
Данные и версия & Период, версия политики, $N_J^{\mathrm{obs}}$, $N_J^+$,
число сессий без зарегистрированной цели и причины исключений. \\
Цепочка охвата & $\sum_s V_s$, $\sum_s U_s$, $\sum_s\operatorname{Hit}_s$;
три потери~\eqref{eq:eval-loss-decomposition},
условный охват~\eqref{eq:eval-candidate-conditional} и $\sum_s V_s/N_J^+$.
При нулевом знаменателе доля не определена. \\
Итоговая выдача & $\operatorname{Hit@20}_J$,
$\operatorname{Coverage}_{\mathrm{cand},J}$; парная разность с базой
и её доверительный интервал при сравнении моделей. \\
Источники & Доступность, размер списков, пересечение, уникальный охват;
потери ограничения квот и $\Delta_a^{\mathrm{fixed}}$
из~\eqref{eq:eval-source-removal}. Исключение с заполнением резерва
и новым обучением проводят лишь как отдельный выбранный опыт
(раздел~\ref{sec:eval-source-contribution}). \\
Условия запроса & Число сессий и охват по заранее заданным группам:
поиск, корзина, наличие модальностей, история товара и сессии.
Отдельно --- покупки из наблюдаемой корзины и вне её. \\
Стоимость & Время обучения, часы ускорителя, максимумы оперативной
и видеопамяти; оборудование, нагрузка, P95, P99, частоты ошибок,
превышений срока и резервных ответов; согласие приближённого поиска с точным. \\
Следующее действие & Причина выбора следующего сравнения
по разделу~\ref{sec:eval-next-step}, фиксируемые данные и бюджет. \\
\bottomrule
\end{longtable}

Сохраняют короткие списки источников до объединения, происхождения
товаров, выдачу и причины резервных ответов. Потери квот оценивают
по рангам целей либо при заранее ограниченной глубине диагностических
списков; неизвестные ранги не восстанавливают предположением.
Технический отказ не образует четвёртое слагаемое разложения:
он может затронуть несколько этапов, а резерв --- сохранить попадание.

\section{Численный пример}\label{app:example}

Рассмотрим три допустимых товара $a,b,c$. После нормирования строки
одного графового отношения получены $P_{ab}=0.8$, $P_{ac}=0.2$;
суммарный вес строки равен $d_a=20$, а $\lambda_G=10$.
В истории только просмотр $a$, поэтому
\begin{equation}\label{eq:app-graph-example}
\rho_a=\frac23,\qquad \alpha_{q1}=1,\qquad
g_{qb}=\frac{8}{15},\qquad g_{qc}=\frac{2}{15}.
\end{equation}
Эти числа не являются вероятностями покупки.
Пусть нормированные текстовые векторы и запрос равны
\begin{equation}\label{eq:app-content-example}
\bar d_a=(1,0),\quad \bar d_b=(0.8,0.6),\quad
\bar d_c=(0,1),\quad\bar z_q=(0,1).
\end{equation}
Текстовое сходство с просмотренным товаром равно $0.8$ для $b$
и $0$ для $c$, а поисковое --- $0.6$ для $b$ и $1$ для $c$.
Граф и история предпочитают $b$, поиск --- $c$.
Это осмысленное расхождение признаков; его разрешает обучаемый
ранжировщик, а не неоговорённое сложение разных шкал.

Если зарегистрированная будущая покупка --- $c$, а кандидаты $\{b,c\}$, то
$y_{qb}=0$, $y_{qc}=1$. При равных исходных баллах,
$\beta_R=1$ и весе пары $\omega=1$ производные равны
\begin{equation}\label{eq:app-gradient-example}
g^R_{qc}=-\frac12,\quad g^R_{qb}=\frac12,
\qquad h^R_{qc}=h^R_{qb}=\frac14.
\end{equation}
Шаг минимизации повышает балл $c$ относительно $b$.
Числа иллюстрируют производные одной пары и не описывают
готовую модель на Coveo.

Теперь пусть из 1000 допущенных записанных сессий 30 имеют
зарегистрированную будущую покупку, а первая выдача охватила такой
товар в 12 сессиях. Кандидаты первого запроса содержали его в 20 сессиях.
Тогда
\begin{equation}\label{eq:app-proxy-example}
N_J^{\mathrm{obs}}=1000,\quad N_J^+=30,\quad
\operatorname{Hit@20}_J=\frac{12}{30}=0.4,\quad
\operatorname{Coverage}_{\mathrm{cand},J}=\frac{20}{30}.
\end{equation}
Если допустимая цель существовала в 24 сессиях, три потери в долях равны
$(30-24)/30$, $(24-20)/30$ и $(20-12)/30$. Их сумма $0.6$
дополняет Hit@20 до единицы. Остальные 970 сессий не становятся
подтверждёнными отрицательными исходами.
Условный охват отбора равен $20/24$; ни разрыв, ни эта доля
не измеряют дополнительные покупки от рекомендаций.

В отдельном рандомизированном эксперименте 3000 покупающих сессий из 100 тысяч в контроле
и 3200 из 100 тысяч в новом варианте дают
\begin{equation}\label{eq:app-business-example}
\widehat\Delta_{1000}^{\mathrm{ITT}}=2,\qquad
\widehat{\operatorname{se}}_{\mathrm{ITT}}\approx0.775.
\end{equation}
Нормальный 95\%-й интервал приблизительно равен $[0.48,3.52]$.
Он поддерживает положительный эффект при выполнении предположений
эксперимента, но не доказывает превышение порога одна дополнительная
покупающая сессия на тысячу: нижняя граница меньше единицы.
Все числа приложения вымышлены для проверки формул.

\section{Проверяемые свойства реализации}\label{app:tests}

Проверяют активные модули. Изменение данных, признаков, модели или
правил выдачи требует повторной проверки затронутых свойств.

\subsection{Сопоставление требований и порядок проверок}\label{app:tests-normative}

Каждое требование ниже имеет постоянный ID, сопоставленный со стандартным
тестом в его имени или метке. Код теста задаёт вход или рецепт его
построения, независимый эталон, допуски и начальное состояние генератора.
Результат сохраняет фактический отчёт, версию кода, конфигурацию и seed;
отдельная предварительная регистрация или хешированный паспорт каждого
теста не обязательны. Срок определяется применимостью по~\ref{app:run-stages}.
Отсутствие эталона или отчёта не считается успешной проверкой.

Счётчики, множества, маски, версии и порядок сравниваются точно.
Для чисел требуется
$|x-x_{\mathrm{ref}}|\le\mathrm{atol}+\mathrm{rtol}|x_{\mathrm{ref}}|$;
допуски заданы в~\ref{sec:ops-numerics}, для производных деревьев ---
в~\ref{sec:model-ranker-implementation}. Неожиданный NaN/Inf означает
отказ; исключения, например задержка $+\infty$, перечислены явно.
Генеративные проверки содержат не менее 1000 случаев на свойство,
начальное состояние генератора --- 17. Сохраняют его алгоритм, версию
и минимальный контрпример. Численные примеры ниже обязательны.

Сроки: B01, B02, B06, B09, B13 и \code{B10.data} --- \code{preflight};
B04, B05, B07, B08 и \code{B10.ranker} --- до применения соответствующих
расчётов, с проверкой фактических артефактов на \code{post_training}.
B14--B17 и \code{B03.service} (доставка и состояние) относятся
к \code{post_load}; \code{B03.data} (схема и канонизация) --- к \code{preflight}.
Применимые свойства воспроизведения локальной оценки проверяются
до её расчёта без обязательной реализации доставки и восстановления сервиса.
B11 проверяет отключение калибровки до первого расчёта выдачи.
\code{B10.metadata} обязателен для полного варианта
без метаданных перед его расчётами и перед $D$; конфигурация варианта
фиксируется до $C_{\mathrm{select}}$.
B12 выполняют до расчёта заявленных временных интервалов
и на фактических входах перед $D$.
Для выборочного обучения \code{X02.support} выполняют на
\code{preflight}, а \code{X02.loss} --- на \code{post_training}.
Прочие математические проверки включённых исследовательских компонентов
выполняют до применения соответствующих расчётов и подтверждают
на \code{post_training}. Применимые B10--B12 повторяют на фактических
артефактах до \code{pre_D}. B18 и реальный протокол эксперимента
обязательны к \code{pre_AB}, но не к \code{post_training}.
Результаты на $D$ и реальные исходы A/B
записывают только на \code{post_D} и \code{post_AB} соответственно.
Тесты включённого ANN завершают до его допуска к обслуживанию.

\subsection{Базовая система}\label{app:tests-base}

\begin{testcase}[B01. Время и канонизация]
На сетке миллисекунд проверяют события ровно в $t_q$ и после него,
покупку в первом блоке и на $t_{0,s}+H$, начало на границе допуска,
границу записанной истории $t_{\mathrm{last},s}^{\mathrm{obs}}+g=b_{A_G}$.
Равенство исключает сессию из $A_G$; календарный запас не подтверждает
полноту телеметрии.
Эталон --- неравенства~\ref{sec:time} и~\ref{sec:requests}, без округления.
Перестановка одновременных строк сохраняет результат.
Две одинаковые строки с разными ID имеют сырую кратность 2:
\code{set_v1} даёт 1, \code{multiset_v1} --- 2; повтор того же ID
не увеличивает кратность. Будущие события не меняют прошлые входы.
Для последней группы с девятью просмотрами $A$ и одним $B$
(разные event ID) и сходствами с кандидатом $1$ и $0$
среднее~\eqref{eq:model-last-group-mean} равно $0.5$ в \code{set_v1}
и $0.9$ в \code{multiset_v1}. Без вектора $B$ оба средних равны $1$;
без обоих векторов значение и маска равны $0$.
Смена семантики кратностей сохраняет множество различных товаров
группы и максимумы, но может менять отбор по усреднённым баллам.
\end{testcase}

\begin{testcase}[B02. Поисковая неоднозначность]
Два вектора $(1,0)$ и $(0,1)$ при одинаковых $(s,t)$ сохраняют
неоднозначность после канонизации и любой перестановки строк.
Два одинаковых пригодных вектора не создают ложной неоднозначности.
Непригодный или отсутствующий вектор в последнем поисковом блоке
отключает контекст даже при наличии в нём одного пригодного вектора.
Замена списков \code{clicked_skus_hash}, \code{product_skus_hash}
не меняет входы и ответ основного варианта.
\end{testcase}

\begin{testcase}[B03. Повторы, закрытые блоки и пределы]
Принять 100 событий и повторить тот же блок через 1 мс:
недостаток токенов не меняет счётчики, полноту и прежний ответ.
То же верно для проверяемого повтора после $\delta_{\mathrm{late}}$.
При $\delta_{\mathrm{future}}=0$ новый блок с $t=a_{\mathrm{block}}$
проходит проверку времени; опережение на 1 мс или час даёт
\code{timestamp_future} без изменения истории и $t_{\mathrm{last},s}$.
Проверяют обе границы календарного диапазона, неверные единицы,
переполнение и повтор принятого блока после изменения допустимого диапазона.
Проверить изменённую нагрузку при прежнем ID, список ID той же длины
с другим содержимым, чужого владельца,
пустой блок, неверное число событий и неполный конверт.
Конверт из двух частей по 60 событий не принимается по частям;
новые ID для уже закрытого времени дают \code{closed_block_conflict}.
При 950 принятых событиях новый блок из 100 отклоняется целиком.
Принятая история сохраняется; разрыв новой доверенной части отражается
по~\ref{sec:ops-transport-block}. Недоверенный ввод не портит чужую сессию.
Два \code{remove} не дают отрицательного количества; одновременные
\code{add}/\code{remove} без порядка оставляют корзину неизвестной.
Только \code{add} даёт соответствующее состояние, только \code{remove}
--- состояние удаления; отсутствие обоих --- неизвестность.
\end{testcase}

\begin{testcase}[B04. Граф]
Сто повторений пары в одной сессии дают поддержку 1 и вес не более 1.
Непустая нормированная строка суммируется в 1, пустая остаётся нулевой.
При фиксированном $W_{ij}^r>0$ увеличение остальных весов строки
уменьшает $\rho_i^rP_{ij}^r=W_{ij}^r/(d_i^r+\lambda_G)$.
Отдельный штраф ребра не подменяет опорную матрицу путей.
\end{testcase}

\begin{testcase}[B05. Числа и производные ранжировщика]
Отсутствующий, нулевой, неконечный или неверно размерный внешний
вектор получает нулевую маску и отдельную причину; исправные ветви работают.
Неконечный внутренний признак или потеря останавливает обучение,
а при обслуживании включает резерв. Проверяют противоположные
векторы, координаты порядка $10^{200}$, пустые маски и баллы
$-10^4,-100,0,100,10^4$. Эталон использует масштабированную норму
и устойчивые softplus/LSE. Исходные пространства модальностей не смешивают.
Для включённой нейросети энтропию и KL проверяют также при насыщении
по N06; конечного softmax без конечных производных недостаточно.

Для пары $(1,0)$ с разностями баллов $-5,-0.1,0,0.1,5$
сравнивают производные с центральными разностями, шаг $10^{-4}$,
при фиксированных рангах и весах: абсолютный допуск $10^{-7}$
для градиента, $10^{-6}$ для гессиана. Проверяют прямую сумму пар,
приведение FP32, нулевую группу и нулевую сумму парных градиентов.
На малом дереве при $\gamma_T=0$ проверяют значение листа и
$\operatorname{Gain}_{\mathrm{LGBM}}=2\operatorname{Gain}_{\mathrm{math}}$;
внешние веса Dataset выключены.
Отдельный пример содержит не менее пяти категорий: одну с $G=H=0$,
остальные --- с информативными парами и положительной суммой гессианов.
Проверка должна пройти ветвь группового категориального разбиения;
при \code{cat_smooth=1} внутренние ключи и результат конечны.
\end{testcase}

\begin{testcase}[B06. Зарегистрированные цели и допуск]
Три допущенные сессии имеют
неизвестное покрытие и зарегистрированные цели $\{a,b\}$, $\{z\}$,
$\varnothing$, где $z\notin\mathcal I_G$.
Первая выдача содержит $a$ только для первой сессии:
$N_J^{\mathrm{obs}}=3$, $N_J^+=2$,
$\operatorname{Hit@20}_J=1/2$.
Покупка $z$ остаётся в целях и учитывается как промах;
будущий каталог не подставляется.
Пустая цель не становится подтверждённым отрицательным исходом;
положительная цель не подтверждает полноту телеметрии.
Повторная покупка того же SKU не меняет множество, последующие выдачи
не меняют метрику. Покупка в блоке первого просмотра исключает сессию
из подходящих. Для подходящей сессии покупка в $t_{0,s}+H$
включена в цель, позже --- исключена. При $t_{0,s}+H+g=b_J$ сессия исключается
независимо от наличия покупки. Пропуск опубликованных строк выбранной
сессии считается ошибкой чтения, а не неизвестностью телеметрии.
\end{testcase}

\begin{testcase}[B07. Кандидаты, доставка и разложение потерь]
Квоты базы суммируются в 300; после фильтрации и заполнения
$L_q\subseteq\mathcal C_q^{\mathrm{served}}\subseteq\mathcal I_q$,
$|\mathcal C_q^{\mathrm{served}}|\le M$.
Малый каталог не дополняют недопустимыми товарами; общий предел $K$
не требует равной фактической длины ответов.
Покупку вне кандидатов не добавляют при обучении или оценке.

Проверяют штатную ветвь, резерв, отказ до выбора и отказ доставки
после выбора. В последнем случае список пуст, выбранные кандидаты
сохраняются; до выбора пусты оба множества. Задним числом ветви не объединяют.
Четыре покупающие сессии с потерей на каждом из трёх этапов и одним
попаданием дают $(N_J^+,\sum V_s,\sum U_s,\sum\operatorname{Hit}_s)
=(4,3,2,1)$, три потери по $1/4$.
Повторное попадание не добавляет единицу. Всегда
$\operatorname{Hit}_s\le U_s\le V_s\le1$;
при $\sum V_s=0$ условный охват не определён.
\end{testcase}

\begin{testcase}[B08. Источники и резерв]
Цель с происхождениями $\{a,b\}$ остаётся после удаления $a$:
$\Delta_a^{\mathrm{fixed}}=0$. Удаление единственного источника
теряет цель; заполнение резервом может её вернуть.
Эти два расчёта регистрируют раздельно. Отказ не удаляет сессию
из знаменателя; учитывается фактически доставленный ответ.
\end{testcase}

\begin{testcase}[B09. Недоступная покупочная частота]
Признак исторической вероятности покупки отсутствует в схеме признаков;
неизвестные исходы не заменяются нулями ради его расчёта.
\end{testcase}

\begin{testcase}[B10. Разделение данных и выбор]
Части \code{B10.data}, \code{B10.ranker} и \code{B10.metadata}
проверяют данные, ранжировщик и вариант без метаданных в применимые сроки.
Изменение $D$ не меняет модель, перечень вариантов, длину блока
или правила отбора; изменение зарезервированного $E$ не меняет расчёты.
Проверяют отсутствие
общих сессий у семи частей, запрет доработки по $C_{\mathrm{select}}$
и порядок разрешения равенств. Произвольная замена предоставленных
метаданных не меняет вариант~\ref{sec:metadata-sensitivity}.
\end{testcase}

\begin{testcase}[B11. Отключённая калибровка]
Паспорт и расчёты не требуют калибратора, данные $E$ не читаются,
а баллы выдачи не объявляются вероятностями покупки.
Штраф $\lambda_D=0$; вероятностная процедура разнообразия отключена.
\end{testcase}

\begin{testcase}[B12. Временные интервалы]
При $J=4,m=2$ начала $(1,3)$ дают индексы $(1,2,3,4)$;
обе политики переносятся вместе, статистика --- отношение сумм.
Пустые календарные ячейки сохраняются, календарный запас $H+g$ не добавляется.
Нулевой знаменатель и нехватка блоков дают \code{insufficient}.
То же состояние требуется при всех $d_s=0$, всех Hit=0 или всех Hit=1,
даже если формальный интервал имеет нулевую ширину.
Проверяют границы 99/100 сессий и 19/20 блоков
по~\eqref{eq:eval-information-gate}; достаточные счётчики не разрешают
вырожденный интервал, например при одинаково сбалансированных ячейках.
Формула одного охвата совпадает с заменой $d_s$ на Hit.
Знаменатель повторной выборки равен
числу сессий с зарегистрированной целью, множителя 1000 нет;
ячейки без таких сессий остаются на календарной сетке.

Для $B_0$ независимых равновероятных знаков $\pm1$,
каждый из которых повторён $m$ раз,
дисперсия среднего равна $1/B_0$, а независимый расчёт ошибочно даёт
$1/(mB_0)$. Проверка покрытия заранее фиксирует число повторений
и биномиальный допуск ошибки моделирования. Не требуется, чтобы
каждый случайный интервал был шире интервала для независимых наблюдений.
\end{testcase}

\begin{testcase}[B13. Этапы паспорта]
При заполненных предварительных условиях отсутствие будущих
результатов обучения/нагрузки не блокирует \code{preflight}.
Локальный \code{ideal_deterministic} допускает выбор и оценку на $D$
без \code{post_load}; отсутствие сервиса не даёт готовности к пилоту.
Первый $C_{\mathrm{smoke}}$ не требует B18 или готового варианта
без метаданных. Пропуск этих проверок в их
обязательные сроки отклоняется. Результаты сохраняются отдельно;
изменение уже закреплённого плана или несовпадение его версии
с используемыми кодом, входами и настройками обнаруживается до расчётов.
На сроке соответствующего теста отсутствие результата даёт \code{not_run}.
Отрицательный отчёт не превращается в \code{passed} из-за его наличия.
Проверяют условную обязательность полей включённых модулей.
Протокол оценки задан явно; неизвестный режим отклоняется.
Отключённая калибровка не блокирует \code{pre_D}, но успешный
исторический отчёт не подтверждает полноту учёта исходов для A/B.
Неверный тип, неизвестное состояние, несогласованные времена,
нарушение целостности входов и сроков применимых проверок отклоняются.
\end{testcase}

\begin{testcase}[B14. Версии и воспроизведение]
Пакетная и последовательная обработка одного видимого префикса
дают одинаковые маски, категории, кандидаты и версии; числа ---
в объявленном допуске. Позднее событие не меняет прошлый ответ.
При $t_q=10{:}00$, $a_q=10{:}05$ пакет публикации 10:03 запрещён.
Трасса отключённого источника изображений воспроизводит ту же
маску и $\mathcal C_q^{\mathrm{served}}$; без трассы здоровый повтор
не объявляется измеренным $U_s$. Удалённые входы не заменяются их хешем.
\end{testcase}

\begin{testcase}[B15. Нагрузка, память и восстановление]
Отказы, просроченные и недоставленные запросы остаются в нагрузке
и знаменателях. Приём входа и очередь входят в срок; P99 этапов
не складываются. Своевременный резервный ответ остаётся резервным.
Точный поиск с пиком 30 ГиБ не проходит общий предел 28 ГиБ;
учитываются оба пакета при обновлении.
Проверяют переполнение очереди, Sybil/SKU-всплеск, подпись, потерю
контрольной точки, конкурирующие записи и отказ каталога.
Между просмотрами вставляют более $L_{\mathrm{hist}}$ поисковых блоков:
очередь $L_H$ просмотров сохраняется. Пределы восстановления,
совпадение признаков и RTO/RPO проверяются измерением.
\end{testcase}

\begin{testcase}[B16. Остановка и частоты тревог]
При одном запросе в секунду 20 последовательных технических отказов
валидных запросов запускают защиту, хотя за пять минут нет 1000 запросов.
Три неуспешные внутренние проверки также запускают защиту.
Некорректные запросы и неподлинный новый пакет не откатывают исправную версию.
Для $A=(2,0.5,2,0.5)$ и $h=1$ доля превышений равна $1/2$,
но срабатываний правила двух соседних окон нет.
Пустое окно не даёт нулевой частоты ошибок; пропавший статус считается отказом.
\end{testcase}

\begin{testcase}[B17. Приближённый поиск]
При включённом ANN сравнивают с точным поиском на том же допустимом
каталоге. Значения скалярного произведения $0.05$ и $0.045$ дают
потерю $0.005$ независимо от температуры ранжирующего балла.
Пустой приближённый список при непустом точном не проходит приёмку.
Дубли, недопустимые товары и неполная повторяемость обнаруживаются отдельно.
\end{testcase}

\begin{testcase}[B18. Причинный учёт]\label{test:eval-effect}
Две сессии одного устойчивого ключа получают один вариант и одну
входную сессию в основном знаменателе. Ошибка сервиса не удаляет её
или достоверную покупку. Проверяют SRM и описательный статус после
защитной остановки; пустая группа даёт \code{insufficient}.
План со стратой по клику после назначения отклоняется до запуска;
план по доступным до назначения признакам проверяет правило пропусков,
неизменность принадлежности и сохранение всех назначенных в знаменателе.
При независимых $Z_1,Z_2\sim\mathrm{Bernoulli}(1/2)$ стандартный план
$S_1=1$, $S_2=Z_1$ также отклоняется: при размере страты 1 число назначений
варианта 1 тождественно нулю, а не имеет закон $\mathrm{Bin}(1,1/2)$.
При $k=0$ и $k=n$ проверяют граничные beta-интервалы.
При выполненных остальных условиях $\ell_B=\delta_{\mathrm{accept}}<u_B$
даёт \code{insufficient}, $\ell_B\le u_B=\delta_{\mathrm{accept}}$ ---
\code{failed}, а $\ell_B>\delta_{\mathrm{accept}}$ --- \code{passed}.
При заранее заданных покупках, не зависящих от рекомендаций,
идеальная выдача может иметь высокий Hit@20, но сохраняет
$Y_s(1)=Y_s(0)$ и $\Delta_{1000}^{\mathrm{ITT}}=0$.
Объявление такого результата ростом конверсии означает ошибку
интерпретации метрик.
\end{testcase}

\subsection{Нейросетевое расширение}\label{app:tests-neural}

\begin{testcase}[N01. Графовая опора и пути]
Изменение $A_N$ при фиксированном $A_G$ сохраняет каталог,
$\mathcal N_i^P,D_{ij}^G$ и частоты. При $\eta_N=0$
непустые строки $A_{ij}^N=\nu_{ij}^P$, $\mathcal L_G^N=0$.
На малом графе полный перебор до $H_P$ переходов совпадает с поиском
по состояниям <<вершина, число переходов>> при достаточном $B_P$.
Усечённый поиск возвращает допустимый путь и признак усечения,
но не утверждает точный минимум.
\end{testcase}

\begin{testcase}[N02. Маски и нулевые проекции]
При отключении ветвей оставшиеся веса нормированы; без входов источник пуст.
Для пригодного входа, $Y_q^G\ne\varnothing$ и $u_q^N=0$ пример сохраняется:
$p_N(i\given q)=1/|\mathcal I_q|$, потеря $\log|\mathcal I_q|$.
При $M_q^S=1,m_q^{SN}=0$ поисковая примесь нулевая и история не ослабляется.
Обнуление обученной проекции не удаляет пример.
\end{testcase}

\begin{testcase}[N03. Намерения и поиск]
При $J_q=J_I$ совместная перестановка намерений, параметров и квот
сохраняет представления, оценки и кандидатов.
Для совпавших $h_{q1}=h_{q2}=(1,0)$ и равных весов энтропия равна
$\log2$, хотя направления одинаковы. При одном направлении
попарная диагностика имеет нулевую маску.
Для $h_{q1}=(1,0)$, $z_q^N=(0,1)$ и $g_q^{SN}=1/4$
получают $u_q^N=(3/4,1/4)$ и балл товара $(0,1)$, равный $1/(4\tau_N)$.
При $\bar g_N=0$ поиск ещё может менять веса разных намерений;
отключение всех его путей устраняет это влияние. Маска $p_S$
не генерируется при обслуживании.
\end{testcase}

\begin{testcase}[N04. Покупочная и следующая цель]
Прямой и блочный расчёты полного знаменателя совпадают.
Две покупки имеют массы $1/2$, сумма производных по баллам нулевая;
нецелевой балл влияет на знаменатель.
Для двух просмотров в 0 и 60 без покупки последний просмотр
входит в следующую цель; проверяют также $t_{0,s}+H$ и просмотр
одновременно с первой покупкой по~\eqref{eq:modern-next-time}.
\end{testcase}

\begin{testcase}[N05. Контрастивные сравнения]
В пакете нет двух запросов одной сессии. Одинаковые сигнатуры
и общая покупка исключаются по заданному правилу.
Без допустимых сравнений потеря нулевая, входы и кандидаты не меняются.
Сохраняются порядок пакетов и маски; вариант \code{label_free_v1}
удаляет только фильтр общей покупки, без чтения $C,E,D$.
\end{testcase}

\begin{testcase}[N06. Инициализация и обучение]
Повтор заданного профиля и состояний генераторов воспроизводит
начальные тензоры и порядок пакетов. Общая матрица инициализируется
один раз; при $U_N=0$ исходно $A^N=\nu^P$.
Два шага Adam на скаляре сверяют с независимым расчётом:
ограничение нормы предшествует моментам, L2 учитывается один раз.
При фиксированных параметрах средняя покупочная часть потери
по всем пакетам, включая короткий последний, совпадает с полной
по~\eqref{eq:modern-batch-loss}.
Для логитов $(0,-200)$ в FP32 энтропия конечна и близка к нулю,
KL к $(1/2,1/2)$ близка к $\log2$. Значения и градиенты по логитам
сверяют с устойчивым FP64-эталоном; единственная позиция даёт нулевые
энтропию и градиент, пустая маска --- предусмотренную нулевую ветвь.
Округление вероятности до нуля не вызывает NaN.
Проверяют независимость потоков: отключённый необязательный потребитель
не меняет порядок пакетов и другие случайные преобразования.
\end{testcase}

\begin{testcase}[N07. Квоты и восстановление базы]
Общий предел $M$ соблюдается при любом числе намерений.
Отключение обучаемых источников с восстановлением исходных квот,
признаков и ранжировщика воспроизводит базу; одного обнуления нового балла мало.
При незарегистрированном отказе нейроисточника результат совпадает
с общим резервом: ветвь, $\mathcal C_q^{\mathrm{served}}$ и порядок выдачи
однозначны; баллы резервной выдачи не получают вероятностную трактовку.
\end{testcase}

\subsection{Дополнительные варианты}\label{app:tests-extra}

\begin{testcase}[X01. Старение, публикация и обход графа]
Сдвиг пары событий в прошлое при неизменном интервале сохраняет
внутрисессионный множитель и уменьшает множитель старения;
за пределами окна связь исключается. Граф, пути, векторы и индекс
публикуются согласованным пакетом.
При фиксированных $c_k$ и допущенных целых сессиях события вне окна или после
$c_k$ не меняют пакет. Повторы пары не увеличивают вклад сессии сверх 1;
на сохранённых рёбрах $\rho_i^rP_{ij}^r=W_{ij}^r/(d_i^r+\lambda_G)\le P_{ij}^r$.
Пакет с $r_k\ge t_q$ недоступен; обновление не меняет старый ответ и его версию.
Для включённого PPR на графе без рёбер проверяют
$z^{(1)}=\eta_Rh_q+(1-\eta_R)\pi^{\mathrm{pop}}$
и измеряют стоимость полного шага при росте каталога.
\end{testcase}

\begin{testcase}[X02. Выборочный знаменатель]
\emph{Опора (до обучения в выборочном режиме).}
В порядке $a,b$ при частотах $(1,0)$ и $\eta_P=\delta=10^{-20}$
вычисляется $Q_G\approx(1,1.5\cdot10^{-20})$, но сохранённый $F=(1,1)$.
Для цели $a$ прямая разрешённая масса положительна, а $\rho_b=0$;
до обучения и расходования потока требуется отказ
\code{negative_sampling_precision_unsupported}, не удаление $b$ и
не ожидание бюджетного отказа. Проверяют также строго возрастающий
$F=(2^{-54},3\cdot2^{-55},1)$: его второй интервал положителен,
но не содержит узлов сетки генератора и должен быть отклонён.
Начальные параметры проверяют на фактическом каталоге по всем условиям
\ref{sec:sampled-proposal}; проверка не расходует случайный поток.

\emph{Потеря (после обучения).}
В каталоге $\{a,b,c\}$ при $T_q^a=\{a\}$, нулевых баллах,
$m=1$, $Q_q^a(b)=1/4$, $Q_q^a(c)=3/4$ имеем $Z_q=3$.
Оценка принимает значения $5$ и $7/3$ с этими вероятностями:
$\E[\widehat Z_q^a]=3$, но $\E[\log\widehat Z_q^a]<\log3$.
Средний градиент потери по положительному баллу равен $-22/35$,
точный --- $-2/3$. При фиксированной выборке конечные разности
воспроизводят~\eqref{eq:sampled-gradient}.
Положительные товары входят один раз, отрицательные сохраняют
кратность выборки. Пустое дополнение даёт точный режим.
Проверяют положительную опору $Q_q^a$, исключение всех целей из
отрицательной выборки, несколько целей, нулевые векторы
и покупку вне каталога: она не расширяет его и остаётся возможным промахом.
Среднее $\widehat Z$ сверяют с точным $Z$ с учётом ошибки случайной
выборки и допущенной ошибки дискретизации распределения;
полный перебор дополнения с единичными весами воспроизводит точные
потерю и градиенты. Конечная случайная выборка этого не требует.
Повтор потока $j=4$ сохраняет выбранные товары и число отброшенных попыток;
нулевая задача и пустое дополнение не меняют состояние генератора.
\end{testcase}

\begin{testcase}[X03. Последовательность групп]
Перестановка внутри группы не меняет представление; причинная маска
исключает влияние поздних групп на ранние. Перенос между группами
может менять ответ. Контроль сохраняет типы действий и покупочную цель;
обогащённый действиями вариант обозначают отдельно.
При $\gamma=1,\beta=0$, входе $(0.001,-0.001)$ и $\eps_{\mathrm{LN}}=10^{-5}$
первая координата LN равна $0.001/\sqrt{0.000011}$, примерно $0.3015113$.
\end{testcase}

\begin{testcase}[X04. SID: производные и этапы]
При фиксированном шуме конечные разности воспроизводят
\eqref{eq:sid-soft-gradient}. Проверяют путь от $\mathcal L_R^{\mathrm{SID}}$
к $A_\ell^{\mathrm{SID}}$ через мягкие назначения (ненулевой градиент
на невырожденном примере) и от
$\mathcal L_G^{\mathrm{SID}}$ к входным кодовым векторам декодера.
При фиксированных целевых кодах $\mathcal L_G^{\mathrm{SID}}$ не даёт
производной по $A_\ell^{\mathrm{SID}}$; $e_i,u_q^N$ фиксированы.
Для каждого этапа проверяют набор изменяемых параметров, сброс состояния
Adam и повторяемость шума. На заключительном этапе изменяется только декодер.
Насыщенные назначения сохраняют конечные балансирующую потерю и градиенты;
проверка конечными разностями не применяется к дискретному $\argmax$.
\end{testcase}

\begin{testcase}[X05. Коллизии SID]
При $L_{\mathrm{SID}}=2$ два товара с кодом $(2,1)$
и $R_{\mathrm{SID}}=1$ дают первый
по заданному порядку. Покупка второго не меняет соответствие
и остаётся промахом SID-источника; другой источник может вернуть товар.
Без шума коды и кандидаты повторяются; все выданные товары входят в $\mathcal I_q$.
Случайный контроль с тем же потоком $m=2,j=6$ повторяет уникальные коды.
Несогласованные словари, декодер и дерево не допускаются к обслуживанию.
\end{testcase}

\section{Состав итогового исследовательского результата}\label{app:deliverables}

Передают паспорт, проверки и отчёт~\ref{app:diagnostic-output}
с источниками данных и знаменателями оценок, а также план A/B.
Без соответствующих наблюдений причинный и денежный результаты
имеют состояние \code{not_observed}.

\begingroup
\small
\setlength{\parskip}{0pt}
\begin{thebibliography}{99}
\setlength{\itemsep}{0.3\baselineskip}
\addcontentsline{toc}{section}{Литература и документация}

\bibitem{coveo-data}
Coveo. \emph{SIGIR eCom Data Challenge: dataset documentation}.
\url{https://github.com/coveooss/SIGIR-ecom-data-challenge}.
Описание схемы, анонимизации, признаков и поисковых данных.

\bibitem{coveo-paper}
Tagliabue J., Greco C., Roy J.-F., Bianchi F., Cassani G., Yu B., Chia P. J.
\emph{SIGIR 2021 E-Commerce Workshop Data Challenge}.
2021. \url{https://arxiv.org/abs/2104.09423}.

\bibitem{spare2025}
Peintner A., Mohammadi A. R., Zangerle E.
\emph{Efficient Session-based Recommendation with Contrastive
Graph-based Shortest Path Search}.
ACM Transactions on Recommender Systems, 2025.
\url{https://doi.org/10.1145/3701764}.
Авторская версия:
\url{https://evazangerle.at/publication/peintner-tors-2025/}.

\bibitem{miasrec2024}
Choi M., Kim H.-Y., Cho H., Lee J.
\emph{Multi-intent-aware Session-based Recommendation}.
SIGIR, 2024. \url{https://doi.org/10.1145/3626772.3657928}.
\url{https://arxiv.org/abs/2405.00986}.

\bibitem{macl2025}
Zhang X., Xu B., Ma F., Wang Z., Yang L., Lin H.
\emph{Rethinking Contrastive Learning in Session-based Recommendation}.
Pattern Recognition. Vol.~169, 2026. Article~111924.
\url{https://doi.org/10.1016/j.patcog.2025.111924}.
\url{https://arxiv.org/abs/2506.05044}.

\bibitem{grid2025}
Ju C. M., Collins L., Neves L., Kumar B., Wang L. Y., Zhao T., Shah N.
\emph{Generative Recommendation with Semantic IDs: A Practitioner's Handbook}.
2025. \url{https://arxiv.org/abs/2507.22224}.
Репозиторий авторов: \url{https://github.com/snap-research/GRID}.

\bibitem{diger2026}
Fu J., Ge X., Karatzoglou A., Arapakis I., Verberne S., Jose J. M., Ren Z.
\emph{Differentiable Semantic ID for Generative Recommendation}.
SIGIR, 2026. \url{https://doi.org/10.1145/3805712.3809641}.
\url{https://arxiv.org/abs/2601.19711}.

\bibitem{sasrec2018}
Kang W.-C., McAuley J.
\emph{Self-Attentive Sequential Recommendation}.
ICDM, 2018. \url{https://arxiv.org/abs/1808.09781}.

\bibitem{sampledsoftmax2025}
Chen J., Zhang J., Huang X., Yang Y., Lian D., Chen E.
\emph{Adaptive Sampled Softmax with Inverted Multi-Index:
Methods, Theory and Applications}.
2025. \url{https://arxiv.org/abs/2501.08563}.

\bibitem{logq2025}
Khrylchenko K., Baikalov V., Makeev S., Matveev A., Liamaev S.
\emph{Correcting the LogQ Correction: Revisiting Sampled Softmax
for Large-Scale Retrieval}.
RecSys, 2025. \url{https://doi.org/10.1145/3705328.3748033}.
\url{https://arxiv.org/abs/2507.09331}.

\bibitem{candidatesampling}
TensorFlow. \emph{Candidate Sampling Algorithms}.
\url{https://www.tensorflow.org/extras/candidate_sampling.pdf}.
Сравнение способов обучения с выборкой классов и шумовых примеров.

\bibitem{lambdamart}
Burges C. J. C.
\emph{From RankNet to LambdaRank to LambdaMART: An Overview}.
Microsoft Research Technical Report MSR-TR-2010-82, 2010.
\url{https://www.microsoft.com/en-us/research/publication/from-ranknet-to-lambdarank-to-lambdamart-an-overview/}.

\bibitem{lightgbm}
LightGBM 4.6.0. \emph{Parameters; Python training API;
custom derivatives and split gain implementation}.
\url{https://lightgbm.readthedocs.io/en/v4.6.0/Parameters.html};
\url{https://lightgbm.readthedocs.io/en/v4.6.0/pythonapi/lightgbm.train.html};
\url{https://github.com/microsoft/LightGBM/blob/v4.6.0/python-package/lightgbm/basic.py};
\url{https://github.com/microsoft/LightGBM/blob/v4.6.0/src/treelearner/feature_histogram.hpp}.
Версия \code{lightgbm==4.6.0}, хеши бинарных пакетов и callback
сохраняются в lockfile запуска.

\bibitem{item2vec}
Barkan O., Koenigstein N.
\emph{Item2Vec: Neural Item Embedding for Collaborative Filtering}.
2016. \url{https://arxiv.org/abs/1603.04259}.

\bibitem{ips}
Schnabel T., Swaminathan A., Singh A., Chandak N., Joachims T.
\emph{Recommendations as Treatments: Debiasing Learning and Evaluation}.
2016. \url{https://arxiv.org/abs/1602.05352}.

\bibitem{dr}
Dud\'ik M., Langford J., Li L.
\emph{Doubly Robust Policy Evaluation and Learning}.
2011. \url{https://arxiv.org/abs/1103.4601}.

\bibitem{experiments}
Kohavi R., Tang D., Xu Y.
\emph{Trustworthy Online Controlled Experiments:
A Practical Guide to A/B Testing}.
Cambridge University Press, 2020.
\url{https://doi.org/10.1017/9781108653985}.

\bibitem{holm1979}
Holm S. \emph{A Simple Sequentially Rejective Multiple Test Procedure}.
Scandinavian Journal of Statistics. 1979. Vol.~6, no.~2. P.~65--70.
\url{https://www.jstor.org/stable/4615733}.

\bibitem{bootstrap}
Efron B.
\emph{Bootstrap Methods: Another Look at the Jackknife}.
The Annals of Statistics. 1979. Vol.~7, no.~1. P.~1--26.
\url{https://doi.org/10.1214/aos/1176344552}.

\bibitem{kunsch1989}
Künsch H. R.
\emph{The Jackknife and the Bootstrap for General Stationary Observations}.
The Annals of Statistics. 1989. Vol.~17, no.~3. P.~1217--1241.
\url{https://doi.org/10.1214/aos/1176347265}.

\end{thebibliography}
\endgroup
\end{document}
